
\documentclass[11pt]{article}

\usepackage[T1]{fontenc}
\usepackage{geometry}
\usepackage{hyperref}
\usepackage{longtable}
\usepackage{fancyvrb}
\usepackage{amsmath}
\usepackage{amssymb}

\geometry{margin=1in}

\title{Corvus Keyword Reference}
\author{Automatically Generated}
\date{\today}

\begin{document}

\maketitle



\tableofcontents

\newpage

\section{Introduction}
% What is corvus
Corvus is a scientific workflow framework built with python with 
a focus on calculations of spectroscopy. It is written to facilitate
advanced calculations for novice users. 

Corvus is "property driven" which means that the user
specifies a property to calculate, and Corvus figures out 
how it can obtain that property, and what inputs are needed
from the user. This is done by building a dependency tree, with the root
given by the list of requested properties. The dependency tree is built
by checking the dependencies of (other properties needed to calculate) the 
requested property, then asking if these can be satisfied by any
of the available codes. This procedure is repeated until either all dependency
requirements are met, in which case a workflow is built and run. If all requirements
are not met, the user is notivied of any necessary user input.

Most of the workflows available to standard users are meant to facilitate calculations
using the FEFF10 code. These workflows include loops over specified parameters, configurational
averaging, RIXS, and creation of input files from multiple available structural input such as
.cif, .xyz, materials project id, or vasp output. 


\section{Download and Installation}
\subsection{Getting Corvus}
To get corvus, go to 
\url{https://github.com/times-software/Corvus/releases/latest}
and download the archived .zip or .tgz file.

\subsubsection{Requirements}
Corvus is portable and runs on macOS, Linux, and Windos. 
To install corvus, you will need python. Versions > 3.12 and < 3.14 are supported. Earlier versions may work as well. We also strongly suggest that you use
an environment manager such as conda to install corvus. Our suggestion is to
use the free miniforge package, which can be found at \url{https://conda-forge.org/download/}

Python's venv environment management package can also be used.

To actually run corvus, you will need at least FEFF10. To get FEFF10, go to \url{https://times-software.github.io/feff10/} to register. Registration is free,
and is only used for communication about bugs and or updates, and for usage statistics. After registration, you will receive an email with directions for 
download and installation.

\subsection{Installation}

\subsubsection{Simple Installation using install script}
Once you have a version of Corvus from the github website, you can follow the following simple directions to install. Note that in the below "version" should be
replaced with the version that you downloaded, e.g., Corvus-3.2.0-01.zip.

\paragraph{MacOS} Unzip the downloaded file. Open a terminal and navigate to the file. Run the install.sh script.
\begin{verbatim}
cd Downloads
unzip Corvus-version.zip
cd Corvus-version/
source install.sh
\end{verbatim}

\paragraph{Windows} Unzip the download file. Open a terminal and navigate to the directory. Run the installer.bat script.
\begin{verbatim}
cd Downloads
unzip Corvus-version.zip
cd Corvus-version\Corvus-version
install.bat
\end{verbatim}

\subsubsection{Step-by-step instructions with conda}
In MacOS, just open a terminal, and it will already have conda enabled if you have it installed. In windows, you will need to open a conda terminal, (Miniforge command window if you installed miniforge conda).
\begin{enumerate}
\item Download corvus from the github page.
\item unzip the folder.
\item In the terminal, navigate to the folder
\item Create a conda environment for corvus
\item Install corvus with pip.
\end{enumerate}
\begin{verbatim}
cd Downloads/Corvus-version/
conda create -name your_corvus_env python=">=3.12,<3.14"
conda activate your_corvus_env
pip install .
\end{verbatim}


\section{Running Corvus}
Corvus is run from within a terminal (or miniforge dos command window on windows). If you used the 
install script to setup corvus, there should be an executable script named "Corvus-version.command" or "Corvus-version.bat" 
on your Desktop. To run corvus, double click this scipt, navigate to the input file you want 
to run, and type,
\begin{verbatim}
run-corvus -i your_input_file
\end{verbatim}

If you installed with the step-by-step instructions, you will need to activate your
corvus environment first. With conda, this is done by typing,
\begin{verbatim}
conda activate your_corvus_env
\end{verbatim}

\subsection{Example files}
There are a variety of example input files located in the folder
\begin{verbatim}
your_home_directory/corvus_examples
\end{verbatim}
If you installed with the install.sh or install.bat scripts, your terminal
will start in the examples folder when you use the launcher provided on 
your desktop. You can then run a set of the input files using the run_examples.py
python script,
\begin{verbatim}
python run_examples.py
\end{verbatim}

You can also compare various files with their references, which have
the same name with REFERENCE\_ prepended to it, e.g., 
REFERENCE_Corvus.xanes.out


\section{The Corvus input file}

A Corvus input file is composed of a sequence of
keyword blocks.

Each keyword is followed by a brace-delimited
block containing the data associated with that
keyword.

Example:

\begin{Verbatim}
title
{
    Example Calculation
}
\end{Verbatim}

Supported schema types:

\begin{itemize}

\item S = String

\item I = Integer

\item F = Float

\item L = Logical

\item P = Free-form text

\end{itemize}

\section{The minimal input, and most important keywords}
The starting point for any Corvus calculation is to request a property or list of properties to calculate.
This is done using the keyword target\_list.
From here, corvus will look for available codes that can calculate each target, and figure out what input is needed.
There are many possible targets for Corvus, but here I will list a few of the most important ones for calculating spectra
using FEFF10, and how to use them. In general, when using FEFF10 an average over unique sites in the structure must be taken
in order to get the full spectrum. Thus in order to perform a spectral calculation using an external structure as input (the suggested method)
you must define,
\begin{verbatim}
target_list{ cfavg} # Calculate the average of some property.
cfavg_target{ xanes } # Average the XANES, could also be xes, rixs, exafs, ... or multiple properties.
\end{verbatim}

To specify the structure, you can use a variety of file formats, e.g., .xyz, .cif, XDATCAR, or you can have 
corvus download a structure from the Materials Project.
Example:
\begin{verbatim}
cif_input{ your_structure.cif }
\end{verbatim}

You will also need to specify which element and edge you want, e.g., for the Cu K-edge,
\begin{verbatim}
absorbing_atom_type{ Cu }
feff.edge{ K }
\end{verbatim}

Thus a minimal input file for the Cu K-edge of CuO would look as follows:
\begin{verbatim}
target_list{ cfavg }
cfavg_target{ xanes }
cif_input{ CuO.cif }
absorbing_atom_type{ Cu }
feff.edge{ K }
\end{verbatim}

There are of course many useful and important options available, but this will give you a calculation that has reasonable defaults.
One special target to note is the "loop" target which allows the user to defined loops over the values associate with any keyword or set 
of keywords. For single line keywords, the structure is relatively simple:
\begin{verbatim}
target_list{ loop } # Make corvus produce a loop over some parameters.
loop_target{ cfavg } # The target property to calculate for each calculation is an average

cfavg_target{ xanes } # Average the xanes over different unique sites in the structure.

# Now define the loop:
loop_parameter{
single_line_keyword
values_line1
values_line2
values_line3
...
}
# Optionally, define labels for the loop directories. Must have the same number of entries as loop_parameter.
# loop_labels{
# label1
# label2
# label3
# ...
# }

There are is also a special value for the first line after "loop_parameter" that loops over input files so that one can loop over 
any defined sets of keyword->value combinations, e.g., 
\begin{verbatim}
loop_parameter{
file
file1.in
file2.in
file3.in
...
}
\end{verbatim}
This will cause corvus to perform calculations with all parameters loaded from the main input file + any parameters existing in fileN.in. Input in fileN.in overrides input in the 
main input file.


\section{The corvus GUI (SciGUI)}
SciGUI is a graphical user interface that can be used for Corvus. 
The installers (install.bat, install.sh) will install it automatically,
or you can get it from the github page \url{https://github.com/times-software/SciGUI/releases/latest}.
The GUI can be used to control corvus, open input files, plot output, and setup calculation
in an easy manner. If you installed with the install scripts, there will be a launcher
on your desktop called Corvus-version\_GUI.bat or Corvus-version\_GUI.command (windows, MacOS).
Once launched, you should see a main window with three panels. The left panel shows keyword buttons, which 
can be selected. The right screen shows help for the selected keyword. The top panel has some controls (run and plot) 
as well as filters for selecting which keywords to show on the left panel.

\subsection{General Use}
The GUI shows keywords on the left panel. When a user selects a keyword, 
the right panel shows help for that keyword, along with and "enable" checkbox,
and fillable fields for value input. If a user want to fill a keyword, they 
check the "enable" checkbox, then fill the fields. When the user is finished 
filling out all of the keywords needed for their calculation, they click the run
button. SciGUI will then write a corvus input file with all of the enabled keywords
filled out, then run corvus on the input file. Note that there are many possible
keywords, so there are filters to select keyword by subject or code, and also a search
bar to search for keywords. Finally, there is a "Quick Start" button to help users
get started with specific types of calculations.

\subsection{The controls and filter panel (top panel)}
This panel has the run button, the Quick Start selection, and a variety of ways to filter and search keywords:
\begin{itemize}
\item Filter by code
Which code is this a specific keyword for, or is it a general keyword.
\item Filter by subject
Keywords related to a specific property of subject.
\item Search bar
Search for specific terms within a keyword
\item "Only enabled" checkbox
Show only the enabled keywords.
\end{itemize}

\subsection{The Quick Start Menu}
If a user selects one of the options from the "Quick Start" menu, the 
target\_list keyword (and maybe some others) will automatically be enabled and filled out.
In addition, keywords will be colored according to importance:
\begin{itemize} 
\item RED: Required
\item BLUE: Most useful
\item GRAY: Other/Advanced
\end{itemize}
By default, when the "Quick Start" option is chosen, only required and useful cards are shown.
In order to fill in the minimim requirements for this type of calculation, the user then just 
needs to click on each of the red keywords shown in the keyword panel, and fill in the values
for that keyword. 

\subsection{Validation}
If a user has enabled a keyword (selecting it on the left panel, then clicking the "enable" checkbox on the right)
then they must fill all required fields or uncheck the enabled checkbox before selecting another keyword. Otherwise
a pop-up message will warn the user, and the same keyword will be selected again.


\section{Output directory structure and output files}
The output file and directory structure of Corvus is as follows:
Calculations will be run in sub-directories of your current directory named Corvus\_description
where "description" describes what is in the directory. 

Simple calculation that ask directly for a spectrum calculated by FEFF are calculated in the directory
Corvus1\_FEFF, while asking for an average over configurations (target\_list\{ cfavg \}) produces a top-level
directory named,
Corvus3\_cfavg\_spectrum, where spectrum is the named spectrum that was requested by the user, e.g., 
Corvus3\_cfavg\_xanes. Then each of the subdirectories of Corvus3\_cfavg\_xanes, will have directories named by the
index of the absorber as it was run, a label connected to which atom is the absorber, and the code (FEFF in this case). 

For the special target loop, the top level folders start with CorvusN\_loop\_loopparameter\_loopvalue. For example, if the user requests a loop over the parameter feff.corehole, with 
values none, rpa, fsr, then Corvus produces the directories Corvus1\_loop\_feff.corehole\_none, Corvus1\_loop\_feff.corehole\_rpa, Corvus1\_loop\_feff.corehole\_fsr

Finally, the main output files of Corvus are named Corvus.target\_subtarget.out. Thus for calculations that request xanes and provide the structure directly in the input file, the main output
is in Corvus.xanes.out. For configurational averages (cfavg), the main output is in
Corvus.cfavg\_xanes.out (for xanes calculations). Finally for loop calculations, the main output is inside the top-level loop directories, and named the same as it would be for a single calculation, e.g. Corvus.xanes.out or Corvus.cfavg\_xes.out etc. 


The documentation below lists every recognized
keyword, its expected schema, default value,
and documentation string.



\section{Keyword Index by Section}

\subsection*{Corvus Controls}
\begin{itemize}
\item \hyperref[kw:multiprocessing-level]{multiprocessing\_level}
\item \hyperref[kw:multiprocessing-ncpu]{multiprocessing\_ncpu}
\item \hyperref[kw:scratch]{scratch}
\item \hyperref[kw:target-list]{target\_list}
\item \hyperref[kw:title]{title}
\item \hyperref[kw:usehandlers]{usehandlers}
\item \hyperref[kw:usesaved]{usesaved}
\item \hyperref[kw:write-input-only]{write\_input\_only}
\end{itemize}

\subsection*{General Methodology}
\begin{itemize}
\item \hyperref[kw:clusterradius]{clusterradius}
\item \hyperref[kw:constant-volume]{constant\_volume}
\item \hyperref[kw:keep-symm]{keep\_symm}
\item \hyperref[kw:method]{method}
\item \hyperref[kw:molecule-vacuum-margin]{molecule\_vacuum\_margin}
\item \hyperref[kw:mt-overlap]{mt\_overlap}
\item \hyperref[kw:nkpoints]{nkpoints}
\item \hyperref[kw:nqpoints]{nqpoints}
\item \hyperref[kw:pspfiles]{pspfiles}
\item \hyperref[kw:pw-encut]{pw\_encut}
\item \hyperref[kw:scf-conv]{scf\_conv}
\item \hyperref[kw:xc]{xc}
\end{itemize}

\subsection*{Physical Properties}
\begin{itemize}
\item \hyperref[kw:absorbing-atom]{absorbing\_atom}
\item \hyperref[kw:absorbing-atom-by-label]{absorbing\_atom\_by\_label}
\item \hyperref[kw:absorbing-atom-type]{absorbing\_atom\_type}
\item \hyperref[kw:atomic-charge]{atomic\_charge}
\item \hyperref[kw:cell-angles-abc]{cell\_angles\_abc}
\item \hyperref[kw:cell-scaling-abc]{cell\_scaling\_abc}
\item \hyperref[kw:cell-scaling-iso]{cell\_scaling\_iso}
\item \hyperref[kw:cell-struc-opt-flags]{cell\_struc\_opt\_flags}
\item \hyperref[kw:cell-struc-xyz-red]{cell\_struc\_xyz\_red}
\item \hyperref[kw:cell-vectors]{cell\_vectors}
\item \hyperref[kw:charge]{charge}
\item \hyperref[kw:cif-input]{cif\_input}
\item \hyperref[kw:cluster]{cluster}
\item \hyperref[kw:fermi-shift]{fermi\_shift}
\item \hyperref[kw:ismetal]{ismetal}
\item \hyperref[kw:multiplicity]{multiplicity}
\item \hyperref[kw:number-of-atoms]{number\_of\_atoms}
\item \hyperref[kw:polarization]{polarization}
\item \hyperref[kw:species]{species}
\item \hyperref[kw:spectral-broadening]{spectral\_broadening}
\item \hyperref[kw:spin-moment]{spin\_moment}
\item \hyperref[kw:supercell-dimensions]{supercell\_dimensions}
\item \hyperref[kw:vasp-outcar]{vasp\_outcar}
\item \hyperref[kw:vasp-snapshot]{vasp\_snapshot}
\item \hyperref[kw:vasp-xdatcar]{vasp\_xdatcar}
\item \hyperref[kw:vasp-xml]{vasp\_xml}
\item \hyperref[kw:xyz-input]{xyz\_input}
\item \hyperref[kw:xyz-snapshot]{xyz\_snapshot}
\end{itemize}

\subsection*{Abinit}
\begin{itemize}
\item \hyperref[kw:abinit-ng2qpt]{abinit.ng2qpt}
\item \hyperref[kw:abinit-ngqpt]{abinit.ngqpt}
\item \hyperref[kw:abinit-verbatim]{abinit.verbatim}
\end{itemize}

\subsection*{DMDW}
\begin{itemize}
\item \hyperref[kw:dmdw-ioflag]{dmdw.ioflag}
\item \hyperref[kw:dmdw-nlanc]{dmdw.nlanc}
\item \hyperref[kw:dmdw-paths]{dmdw.paths}
\item \hyperref[kw:dmdw-tempgrid]{dmdw.tempgrid}
\end{itemize}

\subsection*{NWChem}
\begin{itemize}
\item \hyperref[kw:nwchem-basis]{nwchem.basis}
\item \hyperref[kw:nwchem-charge]{nwchem.charge}
\item \hyperref[kw:nwchem-iter]{nwchem.iter}
\item \hyperref[kw:nwchem-mult]{nwchem.mult}
\item \hyperref[kw:nwchem-qmd-dt-nucl]{nwchem.qmd.dt\_nucl}
\item \hyperref[kw:nwchem-qmd-nstep-nucl]{nwchem.qmd.nstep\_nucl}
\item \hyperref[kw:nwchem-qmd-print-xyz]{nwchem.qmd.print\_xyz}
\item \hyperref[kw:nwchem-qmd-snapstep]{nwchem.qmd.snapstep}
\item \hyperref[kw:nwchem-qmd-targ-temp]{nwchem.qmd.targ\_temp}
\item \hyperref[kw:nwchem-qmd-thermostat]{nwchem.qmd.thermostat}
\item \hyperref[kw:nwchem-xas-alpha]{nwchem.xas.alpha}
\item \hyperref[kw:nwchem-xas-nroots]{nwchem.xas.nroots}
\item \hyperref[kw:nwchem-xas-vec]{nwchem.xas.vec}
\item \hyperref[kw:nwchem-xas-xrayenergywin]{nwchem.xas.xrayenergywin}
\item \hyperref[kw:nwchem-xaselem]{nwchem.xaselem}
\item \hyperref[kw:nwchem-xc]{nwchem.xc}
\end{itemize}

\subsection*{FEFF}
\begin{itemize}
\item \hyperref[kw:debyetemp]{debyetemp}
\item \hyperref[kw:dmdw-nlanczos]{dmdw\_nlanczos}
\item \hyperref[kw:feff-absolute]{feff.absolute}
\item \hyperref[kw:feff-afolp]{feff.afolp}
\item \hyperref[kw:feff-atoms]{feff.atoms}
\item \hyperref[kw:feff-cfaverage]{feff.cfaverage}
\item \hyperref[kw:feff-cgrid]{feff.cgrid}
\item \hyperref[kw:feff-chbroad]{feff.chbroad}
\item \hyperref[kw:feff-chsh]{feff.chsh}
\item \hyperref[kw:feff-chwidth]{feff.chwidth}
\item \hyperref[kw:feff-cif]{feff.cif}
\item \hyperref[kw:feff-compton]{feff.compton}
\item \hyperref[kw:feff-config]{feff.config}
\item \hyperref[kw:feff-control]{feff.control}
\item \hyperref[kw:feff-coordinates]{feff.coordinates}
\item \hyperref[kw:feff-corehole]{feff.corehole}
\item \hyperref[kw:feff-corrections]{feff.corrections}
\item \hyperref[kw:feff-corval]{feff.corval}
\item \hyperref[kw:feff-criteria]{feff.criteria}
\item \hyperref[kw:feff-crpa]{feff.crpa}
\item \hyperref[kw:feff-danes]{feff.danes}
\item \hyperref[kw:feff-debye]{feff.debye}
\item \hyperref[kw:feff-density]{feff.density}
\item \hyperref[kw:feff-dims]{feff.dims}
\item \hyperref[kw:feff-dmdw]{feff.dmdw}
\item \hyperref[kw:feff-edge]{feff.edge}
\item \hyperref[kw:feff-egapse]{feff.egapse}
\item \hyperref[kw:feff-egrid]{feff.egrid}
\item \hyperref[kw:feff-ellipticity]{feff.ellipticity}
\item \hyperref[kw:feff-elnes]{feff.elnes}
\item \hyperref[kw:feff-eps0]{feff.eps0}
\item \hyperref[kw:feff-equivalence]{feff.equivalence}
\item \hyperref[kw:feff-equivalence-nmax]{feff.equivalence\_nmax}
\item \hyperref[kw:feff-exafs]{feff.exafs}
\item \hyperref[kw:feff-exchange]{feff.exchange}
\item \hyperref[kw:feff-exelfs]{feff.exelfs}
\item \hyperref[kw:feff-extpot]{feff.extpot}
\item \hyperref[kw:feff-fms]{feff.fms}
\item \hyperref[kw:feff-folp]{feff.folp}
\item \hyperref[kw:feff-fprime]{feff.fprime}
\item \hyperref[kw:feff-hole]{feff.hole}
\item \hyperref[kw:feff-hubbard]{feff.hubbard}
\item \hyperref[kw:feff-interstitial]{feff.interstitial}
\item \hyperref[kw:feff-ion]{feff.ion}
\item \hyperref[kw:feff-iorder]{feff.iorder}
\item \hyperref[kw:feff-jumprm]{feff.jumprm}
\item \hyperref[kw:feff-kmesh]{feff.kmesh}
\item \hyperref[kw:feff-lattice]{feff.lattice}
\item \hyperref[kw:feff-ldec]{feff.ldec}
\item \hyperref[kw:feff-ldos]{feff.ldos}
\item \hyperref[kw:feff-lfms1]{feff.lfms1}
\item \hyperref[kw:feff-lfms2]{feff.lfms2}
\item \hyperref[kw:feff-ljmax]{feff.ljmax}
\item \hyperref[kw:feff-magic]{feff.magic}
\item \hyperref[kw:feff-mbconv]{feff.mbconv}
\item \hyperref[kw:feff-mpse]{feff.mpse}
\item \hyperref[kw:feff-multipole]{feff.multipole}
\item \hyperref[kw:feff-nleg]{feff.nleg}
\item \hyperref[kw:feff-nohole]{feff.nohole}
\item \hyperref[kw:feff-nrixs]{feff.nrixs}
\item \hyperref[kw:feff-nstar]{feff.nstar}
\item \hyperref[kw:feff-numdens]{feff.numdens}
\item \hyperref[kw:feff-opcons]{feff.opcons}
\item \hyperref[kw:feff-overlap]{feff.overlap}
\item \hyperref[kw:feff-pcriteria]{feff.pcriteria}
\item \hyperref[kw:feff-pmbse]{feff.pmbse}
\item \hyperref[kw:feff-polarization]{feff.polarization}
\item \hyperref[kw:feff-potentials]{feff.potentials}
\item \hyperref[kw:feff-potentials-spin]{feff.potentials.spin}
\item \hyperref[kw:feff-preps]{feff.preps}
\item \hyperref[kw:feff-print]{feff.print}
\item \hyperref[kw:feff-rconv]{feff.rconv}
\item \hyperref[kw:feff-real]{feff.real}
\item \hyperref[kw:feff-reciprocal]{feff.reciprocal}
\item \hyperref[kw:feff-restart]{feff.restart}
\item \hyperref[kw:feff-rgrid]{feff.rgrid}
\item \hyperref[kw:feff-rhozzp]{feff.rhozzp}
\item \hyperref[kw:feff-rixs]{feff.rixs}
\item \hyperref[kw:feff-rlprint]{feff.rlprint}
\item \hyperref[kw:feff-rmultiplier]{feff.rmultiplier}
\item \hyperref[kw:feff-rpath]{feff.rpath}
\item \hyperref[kw:feff-rphases]{feff.rphases}
\item \hyperref[kw:feff-rsigma]{feff.rsigma}
\item \hyperref[kw:feff-s02]{feff.s02}
\item \hyperref[kw:feff-scf]{feff.scf}
\item \hyperref[kw:feff-scframp]{feff.scframp}
\item \hyperref[kw:feff-screen]{feff.screen}
\item \hyperref[kw:feff-scxc]{feff.scxc}
\item \hyperref[kw:feff-self]{feff.self}
\item \hyperref[kw:feff-setedge]{feff.setedge}
\item \hyperref[kw:feff-sfconv]{feff.sfconv}
\item \hyperref[kw:feff-sfse]{feff.sfse}
\item \hyperref[kw:feff-sgroup]{feff.sgroup}
\item \hyperref[kw:feff-sig2]{feff.sig2}
\item \hyperref[kw:feff-sig3]{feff.sig3}
\item \hyperref[kw:feff-siggk]{feff.siggk}
\item \hyperref[kw:feff-spin]{feff.spin}
\item \hyperref[kw:feff-ss]{feff.ss}
\item \hyperref[kw:feff-strfac]{feff.strfac}
\item \hyperref[kw:feff-symmetry]{feff.symmetry}
\item \hyperref[kw:feff-target]{feff.target}
\item \hyperref[kw:feff-tdlda]{feff.tdlda}
\item \hyperref[kw:feff-temperature]{feff.temperature}
\item \hyperref[kw:feff-title]{feff.title}
\item \hyperref[kw:feff-tolscf]{feff.tolscf}
\item \hyperref[kw:feff-unfreezef]{feff.unfreezef}
\item \hyperref[kw:feff-xanes]{feff.xanes}
\item \hyperref[kw:feff-xes]{feff.xes}
\item \hyperref[kw:feff-xmcd]{feff.xmcd}
\item \hyperref[kw:nuctemp]{nuctemp}
\item \hyperref[kw:opcons-drude]{opcons.drude}
\item \hyperref[kw:opcons-usesaved]{opcons.usesaved}
\item \hyperref[kw:run-dmdw]{run\_dmdw}
\end{itemize}

\subsection*{ORCA}
\begin{itemize}
\item \hyperref[kw:orca-basis-basis]{orca.basis.basis}
\item \hyperref[kw:orca-basis-decontract]{orca.basis.decontract}
\item \hyperref[kw:orca-ci]{orca.ci}
\item \hyperref[kw:orca-citype]{orca.citype}
\item \hyperref[kw:orca-coords-units]{orca.coords.units}
\item \hyperref[kw:orca-cpcm]{orca.cpcm}
\item \hyperref[kw:orca-expert]{orca.expert}
\item \hyperref[kw:orca-geom-optimizationquality]{orca.geom.optimizationquality}
\item \hyperref[kw:orca-method-allowrhf]{orca.method.allowrhf}
\item \hyperref[kw:orca-method-amass]{orca.method.amass}
\item \hyperref[kw:orca-method-frozencore]{orca.method.frozencore}
\item \hyperref[kw:orca-method-grid]{orca.method.grid}
\item \hyperref[kw:orca-method-gridx]{orca.method.gridx}
\item \hyperref[kw:orca-method-method]{orca.method.method}
\item \hyperref[kw:orca-method-ri]{orca.method.ri}
\item \hyperref[kw:orca-method-runtype]{orca.method.runtype}
\item \hyperref[kw:orca-method-usesymm]{orca.method.usesymm}
\item \hyperref[kw:orca-mp2]{orca.mp2}
\item \hyperref[kw:orca-mp2type]{orca.mp2type}
\item \hyperref[kw:orca-output-pdbfile]{orca.output.pdbfile}
\item \hyperref[kw:orca-output-print]{orca.output.print}
\item \hyperref[kw:orca-output-printlevel]{orca.output.printlevel}
\item \hyperref[kw:orca-output-xyzfile]{orca.output.xyzfile}
\item \hyperref[kw:orca-rel-method]{orca.rel.method}
\item \hyperref[kw:orca-rel-soctype]{orca.rel.soctype}
\item \hyperref[kw:orca-scf-autostart]{orca.scf.autostart}
\item \hyperref[kw:orca-scf-cnvdamp]{orca.scf.cnvdamp}
\item \hyperref[kw:orca-scf-cnvshift]{orca.scf.cnvshift}
\item \hyperref[kw:orca-scf-convergence]{orca.scf.convergence}
\item \hyperref[kw:orca-scf-convergencestrategy]{orca.scf.convergencestrategy}
\item \hyperref[kw:orca-scf-diis]{orca.scf.diis}
\item \hyperref[kw:orca-scf-fracocc]{orca.scf.fracocc}
\item \hyperref[kw:orca-scf-guess]{orca.scf.guess}
\item \hyperref[kw:orca-scf-hftyp]{orca.scf.hftyp}
\item \hyperref[kw:orca-scf-jmatrix]{orca.scf.jmatrix}
\item \hyperref[kw:orca-scf-kdiis]{orca.scf.kdiis}
\item \hyperref[kw:orca-scf-keepdens]{orca.scf.keepdens}
\item \hyperref[kw:orca-scf-keepints]{orca.scf.keepints}
\item \hyperref[kw:orca-scf-kmatrix]{orca.scf.kmatrix}
\item \hyperref[kw:orca-scf-maxiter]{orca.scf.maxiter}
\item \hyperref[kw:orca-scf-nr]{orca.scf.nr}
\item \hyperref[kw:orca-scf-readints]{orca.scf.readints}
\item \hyperref[kw:orca-scf-scfmode]{orca.scf.scfmode}
\item \hyperref[kw:orca-scf-smeartemp]{orca.scf.smeartemp}
\item \hyperref[kw:orca-scf-soscf]{orca.scf.soscf}
\item \hyperref[kw:orca-scf-uno]{orca.scf.uno}
\item \hyperref[kw:orca-scf-usecheapints]{orca.scf.usecheapints}
\item \hyperref[kw:orca-scf-valformat]{orca.scf.valformat}
\item \hyperref[kw:orca-tole]{orca.tole}
\item \hyperref[kw:orca-tolmaxg]{orca.tolmaxg}
\item \hyperref[kw:orca-tolrmaxd]{orca.tolrmaxd}
\item \hyperref[kw:orca-tolrmsd]{orca.tolrmsd}
\item \hyperref[kw:orca-tolrmsg]{orca.tolrmsg}
\end{itemize}

\subsection*{SIESTA}
\begin{itemize}
\item \hyperref[kw:phsf-broadening]{phsf.broadening}
\item \hyperref[kw:phsf-ekeqp]{phsf.ekeqp}
\item \hyperref[kw:phsf-numtimepoints]{phsf.numtimepoints}
\item \hyperref[kw:siesta-AtomicCoordinatesFormat]{siesta.AtomicCoordinatesFormat}
\item \hyperref[kw:siesta-Beta-Broadening]{siesta.Beta.Broadening}
\item \hyperref[kw:siesta-Block-AtomicCoordinatesAndAtomicSpecies]{siesta.Block.AtomicCoordinatesAndAtomicSpecies}
\item \hyperref[kw:siesta-Block-ChemicalSpeciesLabel]{siesta.Block.ChemicalSpeciesLabel}
\item \hyperref[kw:siesta-Block-LatticeParameters]{siesta.Block.LatticeParameters}
\item \hyperref[kw:siesta-Block-PAO-Basis]{siesta.Block.PAO.Basis}
\item \hyperref[kw:siesta-Coreresponse-Broadening]{siesta.Coreresponse.Broadening}
\item \hyperref[kw:siesta-Diag-DivideAndConquer]{siesta.Diag.DivideAndConquer}
\item \hyperref[kw:siesta-DM-InitSpin-AF]{siesta.DM.InitSpin.AF}
\item \hyperref[kw:siesta-DM-MixingWeight]{siesta.DM.MixingWeight}
\item \hyperref[kw:siesta-DM-NumberBroyden]{siesta.DM.NumberBroyden}
\item \hyperref[kw:siesta-DM-NumberBroyden]{siesta.DM.NumberBroyden}
\item \hyperref[kw:siesta-DM-NumberPulay]{siesta.DM.NumberPulay}
\item \hyperref[kw:siesta-DM-Tolerance]{siesta.DM.Tolerance}
\item \hyperref[kw:siesta-ElectronicTemperature]{siesta.ElectronicTemperature}
\item \hyperref[kw:siesta-expert]{siesta.expert}
\item \hyperref[kw:siesta-LatticeConstant]{siesta.LatticeConstant}
\item \hyperref[kw:siesta-LongOutput]{siesta.LongOutput}
\item \hyperref[kw:siesta-MaxSCFIterations]{siesta.MaxSCFIterations}
\item \hyperref[kw:siesta-MeshCutoff]{siesta.MeshCutoff}
\item \hyperref[kw:siesta-NumberOfAtoms]{siesta.NumberOfAtoms}
\item \hyperref[kw:siesta-NumberOfSpecies]{siesta.NumberOfSpecies}
\item \hyperref[kw:siesta-PAO-BasisSize]{siesta.PAO.BasisSize}
\item \hyperref[kw:siesta-PAO-BasisType]{siesta.PAO.BasisType}
\item \hyperref[kw:siesta-SolutionMethod]{siesta.SolutionMethod}
\item \hyperref[kw:siesta-SpinPolarized]{siesta.SpinPolarized}
\item \hyperref[kw:siesta-TD-CoreExcitedAtom]{siesta.TD.CoreExcitedAtom}
\item \hyperref[kw:siesta-TD-CorePerturbationCharge]{siesta.TD.CorePerturbationCharge}
\item \hyperref[kw:siesta-TD-mxPC]{siesta.TD.mxPC}
\item \hyperref[kw:siesta-TD-NumberOfTimeSteps]{siesta.TD.NumberOfTimeSteps}
\item \hyperref[kw:siesta-TD-ShapeOfEfield]{siesta.TD.ShapeOfEfield}
\item \hyperref[kw:siesta-TD-TimeStep]{siesta.TD.TimeStep}
\item \hyperref[kw:siesta-UseSaveData]{siesta.UseSaveData}
\item \hyperref[kw:siesta-XC-authors]{siesta.XC.authors}
\item \hyperref[kw:siesta-XC-functional]{siesta.XC.functional}
\end{itemize}

\subsection*{cif2cell}
\begin{itemize}
\item \hyperref[kw:cif2cell-atomicunits]{cif2cell.atomicunits}
\item \hyperref[kw:cif2cell-cartesian]{cif2cell.cartesian}
\item \hyperref[kw:cif2cell-cif-input]{cif2cell.cif\_input}
\item \hyperref[kw:cif2cell-outfile]{cif2cell.outfile}
\item \hyperref[kw:cif2cell-program]{cif2cell.program}
\item \hyperref[kw:cif2cell-supercell]{cif2cell.supercell}
\end{itemize}

\subsection*{Loop Workflows}
\begin{itemize}
\item \hyperref[kw:loop-labels]{loop\_labels}
\item \hyperref[kw:loop-parameter]{loop\_parameter}
\item \hyperref[kw:loop-target]{loop\_target}
\end{itemize}

\subsection*{Configurational Averaging}
\begin{itemize}
\item \hyperref[kw:cfavg-choose-random-absorbers]{cfavg\_choose\_random\_absorbers}
\item \hyperref[kw:cfavg-egrid]{cfavg\_egrid}
\item \hyperref[kw:cfavg-max-configurations]{cfavg\_max\_configurations}
\item \hyperref[kw:cfavg-max-configurations]{cfavg\_max\_configurations}
\item \hyperref[kw:cfavg-target]{cfavg\_target}
\end{itemize}

\subsection*{Fitting}
\begin{itemize}
\item \hyperref[kw:fit-bond]{fit.bond}
\item \hyperref[kw:fit-datafile]{fit.datafile}
\item \hyperref[kw:fit-parameters]{fit.parameters}
\item \hyperref[kw:fit-target]{fit.target}
\end{itemize}

\subsection*{OCEAN}
\begin{itemize}
\item \hyperref[kw:ocean-abpad]{ocean.abpad}
\item \hyperref[kw:ocean-acell]{ocean.acell}
\item \hyperref[kw:ocean-cnbse-broaden]{ocean.cnbse.broaden}
\item \hyperref[kw:ocean-cnbse-gmres-elist]{ocean.cnbse.gmres.elist}
\item \hyperref[kw:ocean-cnbse-gmres-erange]{ocean.cnbse.gmres.erange}
\item \hyperref[kw:ocean-cnbse-gmres-ffff]{ocean.cnbse.gmres.ffff}
\item \hyperref[kw:ocean-cnbse-gmres-gprc]{ocean.cnbse.gmres.gprc}
\item \hyperref[kw:ocean-cnbse-gmres-nloop]{ocean.cnbse.gmres.nloop}
\item \hyperref[kw:ocean-cnbse-mode]{ocean.cnbse.mode}
\item \hyperref[kw:ocean-cnbse-niter]{ocean.cnbse.niter}
\item \hyperref[kw:ocean-cnbse-rad]{ocean.cnbse.rad}
\item \hyperref[kw:ocean-cnbse-solver]{ocean.cnbse.solver}
\item \hyperref[kw:ocean-cnbse-spect-range]{ocean.cnbse.spect\_range}
\item \hyperref[kw:ocean-cnbse-strength]{ocean.cnbse.strength}
\item \hyperref[kw:ocean-cnbse-xmesh]{ocean.cnbse.xmesh}
\item \hyperref[kw:ocean-core-offset]{ocean.core\_offset}
\item \hyperref[kw:ocean-degauss]{ocean.degauss}
\item \hyperref[kw:ocean-dft]{ocean.dft}
\item \hyperref[kw:ocean-dft-energy-range]{ocean.dft\_energy\_range}
\item \hyperref[kw:ocean-diemac]{ocean.diemac}
\item \hyperref[kw:ocean-ecut]{ocean.ecut}
\item \hyperref[kw:ocean-edges]{ocean.edges}
\item \hyperref[kw:ocean-fband]{ocean.fband}
\item \hyperref[kw:ocean-k0]{ocean.k0}
\item \hyperref[kw:ocean-ldau]{ocean.ldau}
\item \hyperref[kw:ocean-metal]{ocean.metal}
\item \hyperref[kw:ocean-natom]{ocean.natom}
\item \hyperref[kw:ocean-nbands]{ocean.nbands}
\item \hyperref[kw:ocean-ngkpt]{ocean.ngkpt}
\item \hyperref[kw:ocean-nkpt]{ocean.nkpt}
\item \hyperref[kw:ocean-nspin]{ocean.nspin}
\item \hyperref[kw:ocean-nstep]{ocean.nstep}
\item \hyperref[kw:ocean-ntypat]{ocean.ntypat}
\item \hyperref[kw:ocean-occopt]{ocean.occopt}
\item \hyperref[kw:ocean-opf-fill]{ocean.opf.fill}
\item \hyperref[kw:ocean-opf-hfkgrid]{ocean.opf.hfkgrid}
\item \hyperref[kw:ocean-opf-opts]{ocean.opf.opts}
\item \hyperref[kw:ocean-opf-program]{ocean.opf.program}
\item \hyperref[kw:ocean-para-prefix]{ocean.para\_prefix}
\item \hyperref[kw:ocean-photon-energy]{ocean.photon.energy}
\item \hyperref[kw:ocean-photon-operator]{ocean.photon.operator}
\item \hyperref[kw:ocean-photon-polarization]{ocean.photon.polarization}
\item \hyperref[kw:ocean-photon-qhat]{ocean.photon.qhat}
\item \hyperref[kw:ocean-photon-q]{ocean.photon\_q}
\item \hyperref[kw:ocean-pp-list]{ocean.pp\_list}
\item \hyperref[kw:ocean-rprim]{ocean.rprim}
\item \hyperref[kw:ocean-scfac]{ocean.scfac}
\item \hyperref[kw:ocean-scratch]{ocean.scratch}
\item \hyperref[kw:ocean-screen-nbands]{ocean.screen.nbands}
\item \hyperref[kw:ocean-screen-nkpt]{ocean.screen.nkpt}
\item \hyperref[kw:ocean-screen-shells]{ocean.screen.shells}
\item \hyperref[kw:ocean-screen-energy-range]{ocean.screen\_energy\_range}
\item \hyperref[kw:ocean-ser-prefix]{ocean.ser\_prefix}
\item \hyperref[kw:ocean-smag]{ocean.smag}
\item \hyperref[kw:ocean-smag-ixc]{ocean.smag.ixc}
\item \hyperref[kw:ocean-spin-orbit]{ocean.spin-orbit}
\item \hyperref[kw:ocean-toldfe]{ocean.toldfe}
\item \hyperref[kw:ocean-tolwfr]{ocean.tolwfr}
\item \hyperref[kw:ocean-typat]{ocean.typat}
\item \hyperref[kw:ocean-xred]{ocean.xred}
\item \hyperref[kw:ocean-znucl]{ocean.znucl}
\item \hyperref[kw:ocean-zymb]{ocean.zymb}
\end{itemize}

\subsection*{Miscellaneous}
\begin{itemize}
\item \hyperref[kw:ene-int]{ene\_int}
\item \hyperref[kw:mac-diel-const]{mac\_diel\_const}
\item \hyperref[kw:magmom-by-label]{magmom\_by\_label}
\item \hyperref[kw:mp-apikey]{mp\_apikey}
\item \hyperref[kw:mp-id]{mp\_id}
\item \hyperref[kw:spectralFunction-file]{spectralFunction\_file}
\item \hyperref[kw:xanes-file]{xanes\_file}
\end{itemize}

\section{Alphabetical Keyword Index}
\begin{itemize}
\item \hyperref[kw:abinit-ng2qpt]{abinit.ng2qpt} (Abinit)
\item \hyperref[kw:abinit-ngqpt]{abinit.ngqpt} (Abinit)
\item \hyperref[kw:abinit-verbatim]{abinit.verbatim} (Abinit)
\item \hyperref[kw:absorbing-atom]{absorbing\_atom} (Physical Properties)
\item \hyperref[kw:absorbing-atom-by-label]{absorbing\_atom\_by\_label} (Physical Properties)
\item \hyperref[kw:absorbing-atom-type]{absorbing\_atom\_type} (Physical Properties)
\item \hyperref[kw:atomic-charge]{atomic\_charge} (Physical Properties)
\item \hyperref[kw:cell-angles-abc]{cell\_angles\_abc} (Physical Properties)
\item \hyperref[kw:cell-scaling-abc]{cell\_scaling\_abc} (Physical Properties)
\item \hyperref[kw:cell-scaling-iso]{cell\_scaling\_iso} (Physical Properties)
\item \hyperref[kw:cell-struc-opt-flags]{cell\_struc\_opt\_flags} (Physical Properties)
\item \hyperref[kw:cell-struc-xyz-red]{cell\_struc\_xyz\_red} (Physical Properties)
\item \hyperref[kw:cell-vectors]{cell\_vectors} (Physical Properties)
\item \hyperref[kw:cfavg-choose-random-absorbers]{cfavg\_choose\_random\_absorbers} (Configurational Averaging)
\item \hyperref[kw:cfavg-egrid]{cfavg\_egrid} (Configurational Averaging)
\item \hyperref[kw:cfavg-max-configurations]{cfavg\_max\_configurations} (Configurational Averaging)
\item \hyperref[kw:cfavg-max-configurations]{cfavg\_max\_configurations} (Configurational Averaging)
\item \hyperref[kw:cfavg-target]{cfavg\_target} (Configurational Averaging)
\item \hyperref[kw:charge]{charge} (Physical Properties)
\item \hyperref[kw:cif2cell-atomicunits]{cif2cell.atomicunits} (cif2cell)
\item \hyperref[kw:cif2cell-cartesian]{cif2cell.cartesian} (cif2cell)
\item \hyperref[kw:cif2cell-cif-input]{cif2cell.cif\_input} (cif2cell)
\item \hyperref[kw:cif2cell-outfile]{cif2cell.outfile} (cif2cell)
\item \hyperref[kw:cif2cell-program]{cif2cell.program} (cif2cell)
\item \hyperref[kw:cif2cell-supercell]{cif2cell.supercell} (cif2cell)
\item \hyperref[kw:cif-input]{cif\_input} (Physical Properties)
\item \hyperref[kw:cluster]{cluster} (Physical Properties)
\item \hyperref[kw:clusterradius]{clusterradius} (General Methodology)
\item \hyperref[kw:constant-volume]{constant\_volume} (General Methodology)
\item \hyperref[kw:debyetemp]{debyetemp} (FEFF)
\item \hyperref[kw:dmdw-ioflag]{dmdw.ioflag} (DMDW)
\item \hyperref[kw:dmdw-nlanc]{dmdw.nlanc} (DMDW)
\item \hyperref[kw:dmdw-paths]{dmdw.paths} (DMDW)
\item \hyperref[kw:dmdw-tempgrid]{dmdw.tempgrid} (DMDW)
\item \hyperref[kw:dmdw-nlanczos]{dmdw\_nlanczos} (FEFF)
\item \hyperref[kw:ene-int]{ene\_int} (Miscellaneous)
\item \hyperref[kw:feff-absolute]{feff.absolute} (FEFF)
\item \hyperref[kw:feff-afolp]{feff.afolp} (FEFF)
\item \hyperref[kw:feff-atoms]{feff.atoms} (FEFF)
\item \hyperref[kw:feff-cfaverage]{feff.cfaverage} (FEFF)
\item \hyperref[kw:feff-cgrid]{feff.cgrid} (FEFF)
\item \hyperref[kw:feff-chbroad]{feff.chbroad} (FEFF)
\item \hyperref[kw:feff-chsh]{feff.chsh} (FEFF)
\item \hyperref[kw:feff-chwidth]{feff.chwidth} (FEFF)
\item \hyperref[kw:feff-cif]{feff.cif} (FEFF)
\item \hyperref[kw:feff-compton]{feff.compton} (FEFF)
\item \hyperref[kw:feff-config]{feff.config} (FEFF)
\item \hyperref[kw:feff-control]{feff.control} (FEFF)
\item \hyperref[kw:feff-coordinates]{feff.coordinates} (FEFF)
\item \hyperref[kw:feff-corehole]{feff.corehole} (FEFF)
\item \hyperref[kw:feff-corrections]{feff.corrections} (FEFF)
\item \hyperref[kw:feff-corval]{feff.corval} (FEFF)
\item \hyperref[kw:feff-criteria]{feff.criteria} (FEFF)
\item \hyperref[kw:feff-crpa]{feff.crpa} (FEFF)
\item \hyperref[kw:feff-danes]{feff.danes} (FEFF)
\item \hyperref[kw:feff-debye]{feff.debye} (FEFF)
\item \hyperref[kw:feff-density]{feff.density} (FEFF)
\item \hyperref[kw:feff-dims]{feff.dims} (FEFF)
\item \hyperref[kw:feff-dmdw]{feff.dmdw} (FEFF)
\item \hyperref[kw:feff-edge]{feff.edge} (FEFF)
\item \hyperref[kw:feff-egapse]{feff.egapse} (FEFF)
\item \hyperref[kw:feff-egrid]{feff.egrid} (FEFF)
\item \hyperref[kw:feff-ellipticity]{feff.ellipticity} (FEFF)
\item \hyperref[kw:feff-elnes]{feff.elnes} (FEFF)
\item \hyperref[kw:feff-eps0]{feff.eps0} (FEFF)
\item \hyperref[kw:feff-equivalence]{feff.equivalence} (FEFF)
\item \hyperref[kw:feff-equivalence-nmax]{feff.equivalence\_nmax} (FEFF)
\item \hyperref[kw:feff-exafs]{feff.exafs} (FEFF)
\item \hyperref[kw:feff-exchange]{feff.exchange} (FEFF)
\item \hyperref[kw:feff-exelfs]{feff.exelfs} (FEFF)
\item \hyperref[kw:feff-extpot]{feff.extpot} (FEFF)
\item \hyperref[kw:feff-fms]{feff.fms} (FEFF)
\item \hyperref[kw:feff-folp]{feff.folp} (FEFF)
\item \hyperref[kw:feff-fprime]{feff.fprime} (FEFF)
\item \hyperref[kw:feff-hole]{feff.hole} (FEFF)
\item \hyperref[kw:feff-hubbard]{feff.hubbard} (FEFF)
\item \hyperref[kw:feff-interstitial]{feff.interstitial} (FEFF)
\item \hyperref[kw:feff-ion]{feff.ion} (FEFF)
\item \hyperref[kw:feff-iorder]{feff.iorder} (FEFF)
\item \hyperref[kw:feff-jumprm]{feff.jumprm} (FEFF)
\item \hyperref[kw:feff-kmesh]{feff.kmesh} (FEFF)
\item \hyperref[kw:feff-lattice]{feff.lattice} (FEFF)
\item \hyperref[kw:feff-ldec]{feff.ldec} (FEFF)
\item \hyperref[kw:feff-ldos]{feff.ldos} (FEFF)
\item \hyperref[kw:feff-lfms1]{feff.lfms1} (FEFF)
\item \hyperref[kw:feff-lfms2]{feff.lfms2} (FEFF)
\item \hyperref[kw:feff-ljmax]{feff.ljmax} (FEFF)
\item \hyperref[kw:feff-magic]{feff.magic} (FEFF)
\item \hyperref[kw:feff-mbconv]{feff.mbconv} (FEFF)
\item \hyperref[kw:feff-mpse]{feff.mpse} (FEFF)
\item \hyperref[kw:feff-multipole]{feff.multipole} (FEFF)
\item \hyperref[kw:feff-nleg]{feff.nleg} (FEFF)
\item \hyperref[kw:feff-nohole]{feff.nohole} (FEFF)
\item \hyperref[kw:feff-nrixs]{feff.nrixs} (FEFF)
\item \hyperref[kw:feff-nstar]{feff.nstar} (FEFF)
\item \hyperref[kw:feff-numdens]{feff.numdens} (FEFF)
\item \hyperref[kw:feff-opcons]{feff.opcons} (FEFF)
\item \hyperref[kw:feff-overlap]{feff.overlap} (FEFF)
\item \hyperref[kw:feff-pcriteria]{feff.pcriteria} (FEFF)
\item \hyperref[kw:feff-pmbse]{feff.pmbse} (FEFF)
\item \hyperref[kw:feff-polarization]{feff.polarization} (FEFF)
\item \hyperref[kw:feff-potentials]{feff.potentials} (FEFF)
\item \hyperref[kw:feff-potentials-spin]{feff.potentials.spin} (FEFF)
\item \hyperref[kw:feff-preps]{feff.preps} (FEFF)
\item \hyperref[kw:feff-print]{feff.print} (FEFF)
\item \hyperref[kw:feff-rconv]{feff.rconv} (FEFF)
\item \hyperref[kw:feff-real]{feff.real} (FEFF)
\item \hyperref[kw:feff-reciprocal]{feff.reciprocal} (FEFF)
\item \hyperref[kw:feff-restart]{feff.restart} (FEFF)
\item \hyperref[kw:feff-rgrid]{feff.rgrid} (FEFF)
\item \hyperref[kw:feff-rhozzp]{feff.rhozzp} (FEFF)
\item \hyperref[kw:feff-rixs]{feff.rixs} (FEFF)
\item \hyperref[kw:feff-rlprint]{feff.rlprint} (FEFF)
\item \hyperref[kw:feff-rmultiplier]{feff.rmultiplier} (FEFF)
\item \hyperref[kw:feff-rpath]{feff.rpath} (FEFF)
\item \hyperref[kw:feff-rphases]{feff.rphases} (FEFF)
\item \hyperref[kw:feff-rsigma]{feff.rsigma} (FEFF)
\item \hyperref[kw:feff-s02]{feff.s02} (FEFF)
\item \hyperref[kw:feff-scf]{feff.scf} (FEFF)
\item \hyperref[kw:feff-scframp]{feff.scframp} (FEFF)
\item \hyperref[kw:feff-screen]{feff.screen} (FEFF)
\item \hyperref[kw:feff-scxc]{feff.scxc} (FEFF)
\item \hyperref[kw:feff-self]{feff.self} (FEFF)
\item \hyperref[kw:feff-setedge]{feff.setedge} (FEFF)
\item \hyperref[kw:feff-sfconv]{feff.sfconv} (FEFF)
\item \hyperref[kw:feff-sfse]{feff.sfse} (FEFF)
\item \hyperref[kw:feff-sgroup]{feff.sgroup} (FEFF)
\item \hyperref[kw:feff-sig2]{feff.sig2} (FEFF)
\item \hyperref[kw:feff-sig3]{feff.sig3} (FEFF)
\item \hyperref[kw:feff-siggk]{feff.siggk} (FEFF)
\item \hyperref[kw:feff-spin]{feff.spin} (FEFF)
\item \hyperref[kw:feff-ss]{feff.ss} (FEFF)
\item \hyperref[kw:feff-strfac]{feff.strfac} (FEFF)
\item \hyperref[kw:feff-symmetry]{feff.symmetry} (FEFF)
\item \hyperref[kw:feff-target]{feff.target} (FEFF)
\item \hyperref[kw:feff-tdlda]{feff.tdlda} (FEFF)
\item \hyperref[kw:feff-temperature]{feff.temperature} (FEFF)
\item \hyperref[kw:feff-title]{feff.title} (FEFF)
\item \hyperref[kw:feff-tolscf]{feff.tolscf} (FEFF)
\item \hyperref[kw:feff-unfreezef]{feff.unfreezef} (FEFF)
\item \hyperref[kw:feff-xanes]{feff.xanes} (FEFF)
\item \hyperref[kw:feff-xes]{feff.xes} (FEFF)
\item \hyperref[kw:feff-xmcd]{feff.xmcd} (FEFF)
\item \hyperref[kw:fermi-shift]{fermi\_shift} (Physical Properties)
\item \hyperref[kw:fit-bond]{fit.bond} (Fitting)
\item \hyperref[kw:fit-datafile]{fit.datafile} (Fitting)
\item \hyperref[kw:fit-parameters]{fit.parameters} (Fitting)
\item \hyperref[kw:fit-target]{fit.target} (Fitting)
\item \hyperref[kw:ismetal]{ismetal} (Physical Properties)
\item \hyperref[kw:keep-symm]{keep\_symm} (General Methodology)
\item \hyperref[kw:loop-labels]{loop\_labels} (Loop Workflows)
\item \hyperref[kw:loop-parameter]{loop\_parameter} (Loop Workflows)
\item \hyperref[kw:loop-target]{loop\_target} (Loop Workflows)
\item \hyperref[kw:mac-diel-const]{mac\_diel\_const} (Miscellaneous)
\item \hyperref[kw:magmom-by-label]{magmom\_by\_label} (Miscellaneous)
\item \hyperref[kw:method]{method} (General Methodology)
\item \hyperref[kw:molecule-vacuum-margin]{molecule\_vacuum\_margin} (General Methodology)
\item \hyperref[kw:mp-apikey]{mp\_apikey} (Miscellaneous)
\item \hyperref[kw:mp-id]{mp\_id} (Miscellaneous)
\item \hyperref[kw:mt-overlap]{mt\_overlap} (General Methodology)
\item \hyperref[kw:multiplicity]{multiplicity} (Physical Properties)
\item \hyperref[kw:multiprocessing-level]{multiprocessing\_level} (Corvus Controls)
\item \hyperref[kw:multiprocessing-ncpu]{multiprocessing\_ncpu} (Corvus Controls)
\item \hyperref[kw:nkpoints]{nkpoints} (General Methodology)
\item \hyperref[kw:nqpoints]{nqpoints} (General Methodology)
\item \hyperref[kw:nuctemp]{nuctemp} (FEFF)
\item \hyperref[kw:number-of-atoms]{number\_of\_atoms} (Physical Properties)
\item \hyperref[kw:nwchem-basis]{nwchem.basis} (NWChem)
\item \hyperref[kw:nwchem-charge]{nwchem.charge} (NWChem)
\item \hyperref[kw:nwchem-iter]{nwchem.iter} (NWChem)
\item \hyperref[kw:nwchem-mult]{nwchem.mult} (NWChem)
\item \hyperref[kw:nwchem-qmd-dt-nucl]{nwchem.qmd.dt\_nucl} (NWChem)
\item \hyperref[kw:nwchem-qmd-nstep-nucl]{nwchem.qmd.nstep\_nucl} (NWChem)
\item \hyperref[kw:nwchem-qmd-print-xyz]{nwchem.qmd.print\_xyz} (NWChem)
\item \hyperref[kw:nwchem-qmd-snapstep]{nwchem.qmd.snapstep} (NWChem)
\item \hyperref[kw:nwchem-qmd-targ-temp]{nwchem.qmd.targ\_temp} (NWChem)
\item \hyperref[kw:nwchem-qmd-thermostat]{nwchem.qmd.thermostat} (NWChem)
\item \hyperref[kw:nwchem-xas-alpha]{nwchem.xas.alpha} (NWChem)
\item \hyperref[kw:nwchem-xas-nroots]{nwchem.xas.nroots} (NWChem)
\item \hyperref[kw:nwchem-xas-vec]{nwchem.xas.vec} (NWChem)
\item \hyperref[kw:nwchem-xas-xrayenergywin]{nwchem.xas.xrayenergywin} (NWChem)
\item \hyperref[kw:nwchem-xaselem]{nwchem.xaselem} (NWChem)
\item \hyperref[kw:nwchem-xc]{nwchem.xc} (NWChem)
\item \hyperref[kw:ocean-abpad]{ocean.abpad} (OCEAN)
\item \hyperref[kw:ocean-acell]{ocean.acell} (OCEAN)
\item \hyperref[kw:ocean-cnbse-broaden]{ocean.cnbse.broaden} (OCEAN)
\item \hyperref[kw:ocean-cnbse-gmres-elist]{ocean.cnbse.gmres.elist} (OCEAN)
\item \hyperref[kw:ocean-cnbse-gmres-erange]{ocean.cnbse.gmres.erange} (OCEAN)
\item \hyperref[kw:ocean-cnbse-gmres-ffff]{ocean.cnbse.gmres.ffff} (OCEAN)
\item \hyperref[kw:ocean-cnbse-gmres-gprc]{ocean.cnbse.gmres.gprc} (OCEAN)
\item \hyperref[kw:ocean-cnbse-gmres-nloop]{ocean.cnbse.gmres.nloop} (OCEAN)
\item \hyperref[kw:ocean-cnbse-mode]{ocean.cnbse.mode} (OCEAN)
\item \hyperref[kw:ocean-cnbse-niter]{ocean.cnbse.niter} (OCEAN)
\item \hyperref[kw:ocean-cnbse-rad]{ocean.cnbse.rad} (OCEAN)
\item \hyperref[kw:ocean-cnbse-solver]{ocean.cnbse.solver} (OCEAN)
\item \hyperref[kw:ocean-cnbse-spect-range]{ocean.cnbse.spect\_range} (OCEAN)
\item \hyperref[kw:ocean-cnbse-strength]{ocean.cnbse.strength} (OCEAN)
\item \hyperref[kw:ocean-cnbse-xmesh]{ocean.cnbse.xmesh} (OCEAN)
\item \hyperref[kw:ocean-core-offset]{ocean.core\_offset} (OCEAN)
\item \hyperref[kw:ocean-degauss]{ocean.degauss} (OCEAN)
\item \hyperref[kw:ocean-dft]{ocean.dft} (OCEAN)
\item \hyperref[kw:ocean-dft-energy-range]{ocean.dft\_energy\_range} (OCEAN)
\item \hyperref[kw:ocean-diemac]{ocean.diemac} (OCEAN)
\item \hyperref[kw:ocean-ecut]{ocean.ecut} (OCEAN)
\item \hyperref[kw:ocean-edges]{ocean.edges} (OCEAN)
\item \hyperref[kw:ocean-fband]{ocean.fband} (OCEAN)
\item \hyperref[kw:ocean-k0]{ocean.k0} (OCEAN)
\item \hyperref[kw:ocean-ldau]{ocean.ldau} (OCEAN)
\item \hyperref[kw:ocean-metal]{ocean.metal} (OCEAN)
\item \hyperref[kw:ocean-natom]{ocean.natom} (OCEAN)
\item \hyperref[kw:ocean-nbands]{ocean.nbands} (OCEAN)
\item \hyperref[kw:ocean-ngkpt]{ocean.ngkpt} (OCEAN)
\item \hyperref[kw:ocean-nkpt]{ocean.nkpt} (OCEAN)
\item \hyperref[kw:ocean-nspin]{ocean.nspin} (OCEAN)
\item \hyperref[kw:ocean-nstep]{ocean.nstep} (OCEAN)
\item \hyperref[kw:ocean-ntypat]{ocean.ntypat} (OCEAN)
\item \hyperref[kw:ocean-occopt]{ocean.occopt} (OCEAN)
\item \hyperref[kw:ocean-opf-fill]{ocean.opf.fill} (OCEAN)
\item \hyperref[kw:ocean-opf-hfkgrid]{ocean.opf.hfkgrid} (OCEAN)
\item \hyperref[kw:ocean-opf-opts]{ocean.opf.opts} (OCEAN)
\item \hyperref[kw:ocean-opf-program]{ocean.opf.program} (OCEAN)
\item \hyperref[kw:ocean-para-prefix]{ocean.para\_prefix} (OCEAN)
\item \hyperref[kw:ocean-photon-energy]{ocean.photon.energy} (OCEAN)
\item \hyperref[kw:ocean-photon-operator]{ocean.photon.operator} (OCEAN)
\item \hyperref[kw:ocean-photon-polarization]{ocean.photon.polarization} (OCEAN)
\item \hyperref[kw:ocean-photon-qhat]{ocean.photon.qhat} (OCEAN)
\item \hyperref[kw:ocean-photon-q]{ocean.photon\_q} (OCEAN)
\item \hyperref[kw:ocean-pp-list]{ocean.pp\_list} (OCEAN)
\item \hyperref[kw:ocean-rprim]{ocean.rprim} (OCEAN)
\item \hyperref[kw:ocean-scfac]{ocean.scfac} (OCEAN)
\item \hyperref[kw:ocean-scratch]{ocean.scratch} (OCEAN)
\item \hyperref[kw:ocean-screen-nbands]{ocean.screen.nbands} (OCEAN)
\item \hyperref[kw:ocean-screen-nkpt]{ocean.screen.nkpt} (OCEAN)
\item \hyperref[kw:ocean-screen-shells]{ocean.screen.shells} (OCEAN)
\item \hyperref[kw:ocean-screen-energy-range]{ocean.screen\_energy\_range} (OCEAN)
\item \hyperref[kw:ocean-ser-prefix]{ocean.ser\_prefix} (OCEAN)
\item \hyperref[kw:ocean-smag]{ocean.smag} (OCEAN)
\item \hyperref[kw:ocean-smag-ixc]{ocean.smag.ixc} (OCEAN)
\item \hyperref[kw:ocean-spin-orbit]{ocean.spin-orbit} (OCEAN)
\item \hyperref[kw:ocean-toldfe]{ocean.toldfe} (OCEAN)
\item \hyperref[kw:ocean-tolwfr]{ocean.tolwfr} (OCEAN)
\item \hyperref[kw:ocean-typat]{ocean.typat} (OCEAN)
\item \hyperref[kw:ocean-xred]{ocean.xred} (OCEAN)
\item \hyperref[kw:ocean-znucl]{ocean.znucl} (OCEAN)
\item \hyperref[kw:ocean-zymb]{ocean.zymb} (OCEAN)
\item \hyperref[kw:opcons-drude]{opcons.drude} (FEFF)
\item \hyperref[kw:opcons-usesaved]{opcons.usesaved} (FEFF)
\item \hyperref[kw:orca-basis-basis]{orca.basis.basis} (ORCA)
\item \hyperref[kw:orca-basis-decontract]{orca.basis.decontract} (ORCA)
\item \hyperref[kw:orca-ci]{orca.ci} (ORCA)
\item \hyperref[kw:orca-citype]{orca.citype} (ORCA)
\item \hyperref[kw:orca-coords-units]{orca.coords.units} (ORCA)
\item \hyperref[kw:orca-cpcm]{orca.cpcm} (ORCA)
\item \hyperref[kw:orca-expert]{orca.expert} (ORCA)
\item \hyperref[kw:orca-geom-optimizationquality]{orca.geom.optimizationquality} (ORCA)
\item \hyperref[kw:orca-method-allowrhf]{orca.method.allowrhf} (ORCA)
\item \hyperref[kw:orca-method-amass]{orca.method.amass} (ORCA)
\item \hyperref[kw:orca-method-frozencore]{orca.method.frozencore} (ORCA)
\item \hyperref[kw:orca-method-grid]{orca.method.grid} (ORCA)
\item \hyperref[kw:orca-method-gridx]{orca.method.gridx} (ORCA)
\item \hyperref[kw:orca-method-method]{orca.method.method} (ORCA)
\item \hyperref[kw:orca-method-ri]{orca.method.ri} (ORCA)
\item \hyperref[kw:orca-method-runtype]{orca.method.runtype} (ORCA)
\item \hyperref[kw:orca-method-usesymm]{orca.method.usesymm} (ORCA)
\item \hyperref[kw:orca-mp2]{orca.mp2} (ORCA)
\item \hyperref[kw:orca-mp2type]{orca.mp2type} (ORCA)
\item \hyperref[kw:orca-output-pdbfile]{orca.output.pdbfile} (ORCA)
\item \hyperref[kw:orca-output-print]{orca.output.print} (ORCA)
\item \hyperref[kw:orca-output-printlevel]{orca.output.printlevel} (ORCA)
\item \hyperref[kw:orca-output-xyzfile]{orca.output.xyzfile} (ORCA)
\item \hyperref[kw:orca-rel-method]{orca.rel.method} (ORCA)
\item \hyperref[kw:orca-rel-soctype]{orca.rel.soctype} (ORCA)
\item \hyperref[kw:orca-scf-autostart]{orca.scf.autostart} (ORCA)
\item \hyperref[kw:orca-scf-cnvdamp]{orca.scf.cnvdamp} (ORCA)
\item \hyperref[kw:orca-scf-cnvshift]{orca.scf.cnvshift} (ORCA)
\item \hyperref[kw:orca-scf-convergence]{orca.scf.convergence} (ORCA)
\item \hyperref[kw:orca-scf-convergencestrategy]{orca.scf.convergencestrategy} (ORCA)
\item \hyperref[kw:orca-scf-diis]{orca.scf.diis} (ORCA)
\item \hyperref[kw:orca-scf-fracocc]{orca.scf.fracocc} (ORCA)
\item \hyperref[kw:orca-scf-guess]{orca.scf.guess} (ORCA)
\item \hyperref[kw:orca-scf-hftyp]{orca.scf.hftyp} (ORCA)
\item \hyperref[kw:orca-scf-jmatrix]{orca.scf.jmatrix} (ORCA)
\item \hyperref[kw:orca-scf-kdiis]{orca.scf.kdiis} (ORCA)
\item \hyperref[kw:orca-scf-keepdens]{orca.scf.keepdens} (ORCA)
\item \hyperref[kw:orca-scf-keepints]{orca.scf.keepints} (ORCA)
\item \hyperref[kw:orca-scf-kmatrix]{orca.scf.kmatrix} (ORCA)
\item \hyperref[kw:orca-scf-maxiter]{orca.scf.maxiter} (ORCA)
\item \hyperref[kw:orca-scf-nr]{orca.scf.nr} (ORCA)
\item \hyperref[kw:orca-scf-readints]{orca.scf.readints} (ORCA)
\item \hyperref[kw:orca-scf-scfmode]{orca.scf.scfmode} (ORCA)
\item \hyperref[kw:orca-scf-smeartemp]{orca.scf.smeartemp} (ORCA)
\item \hyperref[kw:orca-scf-soscf]{orca.scf.soscf} (ORCA)
\item \hyperref[kw:orca-scf-uno]{orca.scf.uno} (ORCA)
\item \hyperref[kw:orca-scf-usecheapints]{orca.scf.usecheapints} (ORCA)
\item \hyperref[kw:orca-scf-valformat]{orca.scf.valformat} (ORCA)
\item \hyperref[kw:orca-tole]{orca.tole} (ORCA)
\item \hyperref[kw:orca-tolmaxg]{orca.tolmaxg} (ORCA)
\item \hyperref[kw:orca-tolrmaxd]{orca.tolrmaxd} (ORCA)
\item \hyperref[kw:orca-tolrmsd]{orca.tolrmsd} (ORCA)
\item \hyperref[kw:orca-tolrmsg]{orca.tolrmsg} (ORCA)
\item \hyperref[kw:phsf-broadening]{phsf.broadening} (SIESTA)
\item \hyperref[kw:phsf-ekeqp]{phsf.ekeqp} (SIESTA)
\item \hyperref[kw:phsf-numtimepoints]{phsf.numtimepoints} (SIESTA)
\item \hyperref[kw:polarization]{polarization} (Physical Properties)
\item \hyperref[kw:pspfiles]{pspfiles} (General Methodology)
\item \hyperref[kw:pw-encut]{pw\_encut} (General Methodology)
\item \hyperref[kw:run-dmdw]{run\_dmdw} (FEFF)
\item \hyperref[kw:scf-conv]{scf\_conv} (General Methodology)
\item \hyperref[kw:scratch]{scratch} (Corvus Controls)
\item \hyperref[kw:siesta-AtomicCoordinatesFormat]{siesta.AtomicCoordinatesFormat} (SIESTA)
\item \hyperref[kw:siesta-Beta-Broadening]{siesta.Beta.Broadening} (SIESTA)
\item \hyperref[kw:siesta-Block-AtomicCoordinatesAndAtomicSpecies]{siesta.Block.AtomicCoordinatesAndAtomicSpecies} (SIESTA)
\item \hyperref[kw:siesta-Block-ChemicalSpeciesLabel]{siesta.Block.ChemicalSpeciesLabel} (SIESTA)
\item \hyperref[kw:siesta-Block-LatticeParameters]{siesta.Block.LatticeParameters} (SIESTA)
\item \hyperref[kw:siesta-Block-PAO-Basis]{siesta.Block.PAO.Basis} (SIESTA)
\item \hyperref[kw:siesta-Coreresponse-Broadening]{siesta.Coreresponse.Broadening} (SIESTA)
\item \hyperref[kw:siesta-Diag-DivideAndConquer]{siesta.Diag.DivideAndConquer} (SIESTA)
\item \hyperref[kw:siesta-DM-InitSpin-AF]{siesta.DM.InitSpin.AF} (SIESTA)
\item \hyperref[kw:siesta-DM-MixingWeight]{siesta.DM.MixingWeight} (SIESTA)
\item \hyperref[kw:siesta-DM-NumberBroyden]{siesta.DM.NumberBroyden} (SIESTA)
\item \hyperref[kw:siesta-DM-NumberBroyden]{siesta.DM.NumberBroyden} (SIESTA)
\item \hyperref[kw:siesta-DM-NumberPulay]{siesta.DM.NumberPulay} (SIESTA)
\item \hyperref[kw:siesta-DM-Tolerance]{siesta.DM.Tolerance} (SIESTA)
\item \hyperref[kw:siesta-ElectronicTemperature]{siesta.ElectronicTemperature} (SIESTA)
\item \hyperref[kw:siesta-expert]{siesta.expert} (SIESTA)
\item \hyperref[kw:siesta-LatticeConstant]{siesta.LatticeConstant} (SIESTA)
\item \hyperref[kw:siesta-LongOutput]{siesta.LongOutput} (SIESTA)
\item \hyperref[kw:siesta-MaxSCFIterations]{siesta.MaxSCFIterations} (SIESTA)
\item \hyperref[kw:siesta-MeshCutoff]{siesta.MeshCutoff} (SIESTA)
\item \hyperref[kw:siesta-NumberOfAtoms]{siesta.NumberOfAtoms} (SIESTA)
\item \hyperref[kw:siesta-NumberOfSpecies]{siesta.NumberOfSpecies} (SIESTA)
\item \hyperref[kw:siesta-PAO-BasisSize]{siesta.PAO.BasisSize} (SIESTA)
\item \hyperref[kw:siesta-PAO-BasisType]{siesta.PAO.BasisType} (SIESTA)
\item \hyperref[kw:siesta-SolutionMethod]{siesta.SolutionMethod} (SIESTA)
\item \hyperref[kw:siesta-SpinPolarized]{siesta.SpinPolarized} (SIESTA)
\item \hyperref[kw:siesta-TD-CoreExcitedAtom]{siesta.TD.CoreExcitedAtom} (SIESTA)
\item \hyperref[kw:siesta-TD-CorePerturbationCharge]{siesta.TD.CorePerturbationCharge} (SIESTA)
\item \hyperref[kw:siesta-TD-mxPC]{siesta.TD.mxPC} (SIESTA)
\item \hyperref[kw:siesta-TD-NumberOfTimeSteps]{siesta.TD.NumberOfTimeSteps} (SIESTA)
\item \hyperref[kw:siesta-TD-ShapeOfEfield]{siesta.TD.ShapeOfEfield} (SIESTA)
\item \hyperref[kw:siesta-TD-TimeStep]{siesta.TD.TimeStep} (SIESTA)
\item \hyperref[kw:siesta-UseSaveData]{siesta.UseSaveData} (SIESTA)
\item \hyperref[kw:siesta-XC-authors]{siesta.XC.authors} (SIESTA)
\item \hyperref[kw:siesta-XC-functional]{siesta.XC.functional} (SIESTA)
\item \hyperref[kw:species]{species} (Physical Properties)
\item \hyperref[kw:spectral-broadening]{spectral\_broadening} (Physical Properties)
\item \hyperref[kw:spectralFunction-file]{spectralFunction\_file} (Miscellaneous)
\item \hyperref[kw:spin-moment]{spin\_moment} (Physical Properties)
\item \hyperref[kw:supercell-dimensions]{supercell\_dimensions} (Physical Properties)
\item \hyperref[kw:target-list]{target\_list} (Corvus Controls)
\item \hyperref[kw:title]{title} (Corvus Controls)
\item \hyperref[kw:usehandlers]{usehandlers} (Corvus Controls)
\item \hyperref[kw:usesaved]{usesaved} (Corvus Controls)
\item \hyperref[kw:vasp-outcar]{vasp\_outcar} (Physical Properties)
\item \hyperref[kw:vasp-snapshot]{vasp\_snapshot} (Physical Properties)
\item \hyperref[kw:vasp-xdatcar]{vasp\_xdatcar} (Physical Properties)
\item \hyperref[kw:vasp-xml]{vasp\_xml} (Physical Properties)
\item \hyperref[kw:write-input-only]{write\_input\_only} (Corvus Controls)
\item \hyperref[kw:xanes-file]{xanes\_file} (Miscellaneous)
\item \hyperref[kw:xc]{xc} (General Methodology)
\item \hyperref[kw:xyz-input]{xyz\_input} (Physical Properties)
\item \hyperref[kw:xyz-snapshot]{xyz\_snapshot} (Physical Properties)
\end{itemize}

\section{Corvus Controls}

\subsection{multiprocessing\_level}
\label{kw:multiprocessing-level}
\begin{longtable}{|p{1.5in}|p{4.5in}|}
\hline
Keyword & multiprocessing\_level \\
\hline
Kind & scalar \\
\hline
Type & String \\
\hline
Default & cfavg \\
\hline
Deprecated & No \\
\hline
\end{longtable}
 Which level of the calculation do you want to use multiprocessing on.
 Can be "cfavg" or "loop".

\paragraph{Schema}
\begin{Verbatim}
cfavg |X| S
\end{Verbatim}

\subsection{multiprocessing\_ncpu}
\label{kw:multiprocessing-ncpu}
\begin{longtable}{|p{1.5in}|p{4.5in}|}
\hline
Keyword & multiprocessing\_ncpu \\
\hline
Kind & scalar \\
\hline
Type & Integer \\
\hline
Default & 1 \\
\hline
Deprecated & No \\
\hline
\end{longtable}
 Number of processors to use in multiprocessing.
 This should be maximum of the number of cpus on a single
 node.

\paragraph{Schema}
\begin{Verbatim}
1 |X| I
\end{Verbatim}

\subsection{scratch}
\label{kw:scratch}
\begin{longtable}{|p{1.5in}|p{4.5in}|}
\hline
Keyword & scratch \\
\hline
Kind & scalar \\
\hline
Type & String \\
\hline
Default & . \\
\hline
Deprecated & No \\
\hline
\end{longtable}
 Directory for disk scratch for those Handlers that require large amounts of
 disk. If the directory is not present or it can not be created, Corvus reverts
 to the default.
 NOTE: This input variable is not fully implemented yet.

\paragraph{Schema}
\begin{Verbatim}
. |X| S ...
\end{Verbatim}

\subsection{target\_list}
\label{kw:target-list}
\begin{longtable}{|p{1.5in}|p{4.5in}|}
\hline
Keyword & target\_list \\
\hline
Kind & repeating-record \\
\hline
Type & Repeating record with columns: String \\
\hline
Default & Required / No default \\
\hline
Deprecated & No \\
\hline
\end{longtable}
 Space separated list with all the target properties requested for this calculation.


\paragraph{Schema}
\begin{Verbatim}
|X| S ...
|X| ...
\end{Verbatim}

\subsection{title}
\label{kw:title}
\begin{longtable}{|p{1.5in}|p{4.5in}|}
\hline
Keyword & title \\
\hline
Kind & repeating-record \\
\hline
Type & Repeating record with columns: Free-form \\
\hline
Default & Required / No default \\
\hline
Deprecated & No \\
\hline
\end{longtable}
 Title of this calculation.

\paragraph{Schema}
\begin{Verbatim}
This is a Corvus calculation |X| P
                             |X| ...
\end{Verbatim}

\subsection{usehandlers}
\label{kw:usehandlers}
\begin{longtable}{|p{1.5in}|p{4.5in}|}
\hline
Keyword & usehandlers \\
\hline
Kind & scalar \\
\hline
Type & String \\
\hline
Default & Required / No default \\
\hline
Deprecated & No \\
\hline
\end{longtable}
 Explicitely declare what Handlers are to be used in the generation of the
 Workflow. This helps the current simple Workflow generator when a given target
 is provided by more than one Handler.
 Current possible values are:
    Feff, FeffRixs, Dmdw, Abinit, Vasp, Nwchem, Orca

\paragraph{Schema}
\begin{Verbatim}
|X| S ...
\end{Verbatim}

\subsection{usesaved}
\label{kw:usesaved}
\begin{longtable}{|p{1.5in}|p{4.5in}|}
\hline
Keyword & usesaved \\
\hline
Kind & scalar \\
\hline
Type & Logical \\
\hline
Default & False \\
\hline
Deprecated & No \\
\hline
\end{longtable}
 Use previously calculated data rather than recalculating

\paragraph{Schema}
\begin{Verbatim}
False |X| L
\end{Verbatim}

\subsection{write\_input\_only}
\label{kw:write-input-only}
\begin{longtable}{|p{1.5in}|p{4.5in}|}
\hline
Keyword & write\_input\_only \\
\hline
Kind & scalar \\
\hline
Type & Logical \\
\hline
Default & False \\
\hline
Deprecated & No \\
\hline
\end{longtable}
 If True, write input for each of the top level calculations, then exit.
 Not implemented for all calculations.

\paragraph{Schema}
\begin{Verbatim}
False |X| L
\end{Verbatim}

\section{General Methodology}

\subsection{clusterradius}
\label{kw:clusterradius}
\begin{longtable}{|p{1.5in}|p{4.5in}|}
\hline
Keyword & clusterradius \\
\hline
Kind & scalar \\
\hline
Type & Float \\
\hline
Default & 12.0 \\
\hline
Deprecated & No \\
\hline
\end{longtable}
 Radius of clusters created from cif files.

\paragraph{Schema}
\begin{Verbatim}
12.0 |X| F
\end{Verbatim}

\subsection{constant\_volume}
\label{kw:constant-volume}
\begin{longtable}{|p{1.5in}|p{4.5in}|}
\hline
Keyword & constant\_volume \\
\hline
Kind & scalar \\
\hline
Type & Logical \\
\hline
Default & True \\
\hline
Deprecated & No \\
\hline
\end{longtable}
 Keep the simulation cell volume constant in cell optimizations.

\paragraph{Schema}
\begin{Verbatim}
True |X| L
\end{Verbatim}

\subsection{keep\_symm}
\label{kw:keep-symm}
\begin{longtable}{|p{1.5in}|p{4.5in}|}
\hline
Keyword & keep\_symm \\
\hline
Kind & scalar \\
\hline
Type & Logical \\
\hline
Default & True \\
\hline
Deprecated & No \\
\hline
\end{longtable}
 Toggle the preservation of initial symmetry in optimizations.

\paragraph{Schema}
\begin{Verbatim}
True |X| L
\end{Verbatim}

\subsection{method}
\label{kw:method}
\begin{longtable}{|p{1.5in}|p{4.5in}|}
\hline
Keyword & method \\
\hline
Kind & scalar \\
\hline
Type & String \\
\hline
Default & dft \\
\hline
Deprecated & No \\
\hline
\end{longtable}
 Method used to compute the target. This variable is somewhat context-dependent
 and is ccurrently not fully implemented.
 Current possible values are:
    dft, mp2, ccsd

\paragraph{Schema}
\begin{Verbatim}
dft |X| S
\end{Verbatim}

\subsection{molecule\_vacuum\_margin}
\label{kw:molecule-vacuum-margin}
\begin{longtable}{|p{1.5in}|p{4.5in}|}
\hline
Keyword & molecule\_vacuum\_margin \\
\hline
Kind & scalar \\
\hline
Type & Float \\
\hline
Default & Required / No default \\
\hline
Deprecated & No \\
\hline
\end{longtable}
 Added to the largest distance between two atoms in the
 molecule before making a unit cell

\paragraph{Schema}
\begin{Verbatim}
|X| F
\end{Verbatim}

\subsection{mt\_overlap}
\label{kw:mt-overlap}
\begin{longtable}{|p{1.5in}|p{4.5in}|}
\hline
Keyword & mt\_overlap \\
\hline
Kind & repeating-record \\
\hline
Type & Repeating record with columns: String, Float \\
\hline
Default & Required / No default \\
\hline
Deprecated & No \\
\hline
\end{longtable}
 set the muffin-tin overlap for different elements.

\paragraph{Schema}
\begin{Verbatim}
H 0.8  |X| S F
|X| ...
\end{Verbatim}

\subsection{nkpoints}
\label{kw:nkpoints}
\begin{longtable}{|p{1.5in}|p{4.5in}|}
\hline
Keyword & nkpoints \\
\hline
Kind & vector \\
\hline
Type & Vector: Integer, Integer, Integer \\
\hline
Default & 1 1 1 \\
\hline
Deprecated & No \\
\hline
\end{longtable}
 Define the number of k-point in each direction of the grid for reciprocal
 space simulations.

\paragraph{Schema}
\begin{Verbatim}
1 1 1 |X| I I I
\end{Verbatim}

\subsection{nqpoints}
\label{kw:nqpoints}
\begin{longtable}{|p{1.5in}|p{4.5in}|}
\hline
Keyword & nqpoints \\
\hline
Kind & vector \\
\hline
Type & Vector: Integer, Integer, Integer \\
\hline
Default & Required / No default \\
\hline
Deprecated & No \\
\hline
\end{longtable}
 Define the number of q-point perturbations in each direction of the grid for
 density functional perturbation theory (DFPT) simulations.

\paragraph{Schema}
\begin{Verbatim}
 |X| I I I
\end{Verbatim}

\subsection{pspfiles}
\label{kw:pspfiles}
\begin{longtable}{|p{1.5in}|p{4.5in}|}
\hline
Keyword & pspfiles \\
\hline
Kind & repeating-record \\
\hline
Type & Repeating record with columns: String, String \\
\hline
Default & Required / No default \\
\hline
Deprecated & No \\
\hline
\end{longtable}
 List of pseudopotentials for each atom in the system. Each line contains
 the label for the atoms (usually the element % name) and the name of the file
 with the pseudopotential. The required format for the files will depend on the
 Handler used.

\paragraph{Schema}
\begin{Verbatim}
|X| S S
|X| ...
\end{Verbatim}

\subsection{pw\_encut}
\label{kw:pw-encut}
\begin{longtable}{|p{1.5in}|p{4.5in}|}
\hline
Keyword & pw\_encut \\
\hline
Kind & scalar \\
\hline
Type & Float \\
\hline
Default & 15.0 \\
\hline
Deprecated & No \\
\hline
\end{longtable}
 Planwave energy cutoff for reciprocal space simulations, in au.

\paragraph{Schema}
\begin{Verbatim}
15.0 |X| F
\end{Verbatim}

\subsection{scf\_conv}
\label{kw:scf-conv}
\begin{longtable}{|p{1.5in}|p{4.5in}|}
\hline
Keyword & scf\_conv \\
\hline
Kind & scalar \\
\hline
Type & Float \\
\hline
Default & 1.0e-5 \\
\hline
Deprecated & No \\
\hline
\end{longtable}
 Convergence threshold for SCF cycles (HF, DFT, etc) in au. The value set by
 this variable might be internally overridden if the target requires tighter
 convergence settings.

\paragraph{Schema}
\begin{Verbatim}
1.0e-5|X| F
\end{Verbatim}

\subsection{xc}
\label{kw:xc}
\begin{longtable}{|p{1.5in}|p{4.5in}|}
\hline
Keyword & xc \\
\hline
Kind & scalar \\
\hline
Type & String \\
\hline
Default & lda \\
\hline
Deprecated & No \\
\hline
\end{longtable}
 Exchange-correlation functional to use if the "method" selected is "dft".
 Current possible values are:
    lda, pbe, b3lyp

\paragraph{Schema}
\begin{Verbatim}
lda |X| S
\end{Verbatim}

\section{Physical Properties}

\subsection{absorbing\_atom}
\label{kw:absorbing-atom}
\begin{longtable}{|p{1.5in}|p{4.5in}|}
\hline
Keyword & absorbing\_atom \\
\hline
Kind & scalar \\
\hline
Type & Integer \\
\hline
Default & Required / No default \\
\hline
Deprecated & No \\
\hline
\end{longtable}
 Index of the absorbing atom for core spectroscopies such as XANES, EXAFS and
 XES. The index counts from one for the first atom in the associated cluster (see above).

\paragraph{Schema}
\begin{Verbatim}
|X| I ...
\end{Verbatim}

\subsection{absorbing\_atom\_by\_label}
\label{kw:absorbing-atom-by-label}
\begin{longtable}{|p{1.5in}|p{4.5in}|}
\hline
Keyword & absorbing\_atom\_by\_label \\
\hline
Kind & scalar \\
\hline
Type & String \\
\hline
Default & Required / No default \\
\hline
Deprecated & No \\
\hline
\end{longtable}
 Set the absorber by the label associated with it in the cif file.
 Will set all atoms that have a label that starts with this string as absorbers.

\paragraph{Schema}
\begin{Verbatim}
|X| S ...
\end{Verbatim}

\subsection{absorbing\_atom\_type}
\label{kw:absorbing-atom-type}
\begin{longtable}{|p{1.5in}|p{4.5in}|}
\hline
Keyword & absorbing\_atom\_type \\
\hline
Kind & scalar \\
\hline
Type & String \\
\hline
Default & Required / No default \\
\hline
Deprecated & No \\
\hline
\end{longtable}
 Symbol of atom you would like to calculate the XANES for. Will loop over symmetry unique sites defined
 in the CIF file denoted by cif\_input (see below).

\paragraph{Schema}
\begin{Verbatim}
|X| S ...
\end{Verbatim}

\subsection{atomic\_charge}
\label{kw:atomic-charge}
\begin{longtable}{|p{1.5in}|p{4.5in}|}
\hline
Keyword & atomic\_charge \\
\hline
Kind & repeating-record \\
\hline
Type & Repeating record with columns: String, Float \\
\hline
Default & Required / No default \\
\hline
Deprecated & No \\
\hline
\end{longtable}
 charge to be placed on each atoms of a given species.
 atomic\_symbol charge
 ...

\paragraph{Schema}
\begin{Verbatim}
|X| S F
|X| ...
\end{Verbatim}

\subsection{cell\_angles\_abc}
\label{kw:cell-angles-abc}
\begin{longtable}{|p{1.5in}|p{4.5in}|}
\hline
Keyword & cell\_angles\_abc \\
\hline
Kind & vector \\
\hline
Type & Vector: Float, Float, Float \\
\hline
Default & Required / No default \\
\hline
Deprecated & No \\
\hline
\end{longtable}
 Scaling of the a, b an c axes of the simulation cell, in Angstroms.

\paragraph{Schema}
\begin{Verbatim}
|X| F F F
\end{Verbatim}

\subsection{cell\_scaling\_abc}
\label{kw:cell-scaling-abc}
\begin{longtable}{|p{1.5in}|p{4.5in}|}
\hline
Keyword & cell\_scaling\_abc \\
\hline
Kind & vector \\
\hline
Type & Vector: Float, Float, Float \\
\hline
Default & Required / No default \\
\hline
Deprecated & No \\
\hline
\end{longtable}
 Scaling of the a, b an c axes of the simulation cell, in Angstroms.

\paragraph{Schema}
\begin{Verbatim}
|X| F F F
\end{Verbatim}

\subsection{cell\_scaling\_iso}
\label{kw:cell-scaling-iso}
\begin{longtable}{|p{1.5in}|p{4.5in}|}
\hline
Keyword & cell\_scaling\_iso \\
\hline
Kind & scalar \\
\hline
Type & Float \\
\hline
Default & 1.0 \\
\hline
Deprecated & No \\
\hline
\end{longtable}
 Unitless isotropic scaling of the unit cell.

\paragraph{Schema}
\begin{Verbatim}
1.0 |X| F
\end{Verbatim}

\subsection{cell\_struc\_opt\_flags}
\label{kw:cell-struc-opt-flags}
\begin{longtable}{|p{1.5in}|p{4.5in}|}
\hline
Keyword & cell\_struc\_opt\_flags \\
\hline
Kind & repeating-record \\
\hline
Type & Repeating record with columns: Logical, Logical, Logical \\
\hline
Default & Required / No default \\
\hline
Deprecated & No \\
\hline
\end{longtable}
 Toggle the optimization of the x, y and z coordinates of each atom in the
 system.

\paragraph{Schema}
\begin{Verbatim}
|X| L L L
|X| ...
\end{Verbatim}

\subsection{cell\_struc\_xyz\_red}
\label{kw:cell-struc-xyz-red}
\begin{longtable}{|p{1.5in}|p{4.5in}|}
\hline
Keyword & cell\_struc\_xyz\_red \\
\hline
Kind & repeating-record \\
\hline
Type & Repeating record with columns: String, Float, Float, Float \\
\hline
Default & Required / No default \\
\hline
Deprecated & No \\
\hline
\end{longtable}
 Structure of the simulation cell in reduced coordinates if the system is
 extended and has periodic boundary conditions.

\paragraph{Schema}
\begin{Verbatim}
|X| S F F F
|X| ...
\end{Verbatim}

\subsection{cell\_vectors}
\label{kw:cell-vectors}
\begin{longtable}{|p{1.5in}|p{4.5in}|}
\hline
Keyword & cell\_vectors \\
\hline
Kind & matrix \\
\hline
Type & 3 rows of Float, Float, Float \\
\hline
Default & Required / No default \\
\hline
Deprecated & No \\
\hline
\end{longtable}
 Normalized simulation cell vector directions. These vectors are scaled using
 the "cell\_scaling\_iso" and "cell\_scaling\_abc" input variables to generate the
 simulation cell.

\paragraph{Schema}
\begin{Verbatim}
|X| F F F
|X| F F F
|X| F F F
\end{Verbatim}

\subsection{charge}
\label{kw:charge}
\begin{longtable}{|p{1.5in}|p{4.5in}|}
\hline
Keyword & charge \\
\hline
Kind & scalar \\
\hline
Type & Integer \\
\hline
Default & Required / No default \\
\hline
Deprecated & No \\
\hline
\end{longtable}
 Net charge of the system.

\paragraph{Schema}
\begin{Verbatim}
|X| I
\end{Verbatim}

\subsection{cif\_input}
\label{kw:cif-input}
\begin{longtable}{|p{1.5in}|p{4.5in}|}
\hline
Keyword & cif\_input \\
\hline
Kind & scalar \\
\hline
Type & String \\
\hline
Default & Required / No default \\
\hline
Deprecated & No \\
\hline
\end{longtable}
 cif input file name.

\paragraph{Schema}
\begin{Verbatim}
|X| S
\end{Verbatim}

\subsection{cluster}
\label{kw:cluster}
\begin{longtable}{|p{1.5in}|p{4.5in}|}
\hline
Keyword & cluster \\
\hline
Kind & repeating-record \\
\hline
Type & Repeating record with columns: String, Float, Float, Float \\
\hline
Default & Required / No default \\
\hline
Deprecated & No \\
\hline
\end{longtable}
 Structure of the system is the system is not extended (i. e. is a molecule or
 cluster). The coordinates are in Angstroms, and have xyz format, i.e.,
 Atomic\_symbol1 x y z
 Atomic\_symbol2 x y z
 .
 .
 .

\paragraph{Schema}
\begin{Verbatim}
|X| S F F F
|X| ...
\end{Verbatim}

\subsection{fermi\_shift}
\label{kw:fermi-shift}
\begin{longtable}{|p{1.5in}|p{4.5in}|}
\hline
Keyword & fermi\_shift \\
\hline
Kind & scalar \\
\hline
Type & Float \\
\hline
Default & Required / No default \\
\hline
Deprecated & No \\
\hline
\end{longtable}
 Shift of the fermi-energy in eV.

\paragraph{Schema}
\begin{Verbatim}
|X| F
\end{Verbatim}

\subsection{ismetal}
\label{kw:ismetal}
\begin{longtable}{|p{1.5in}|p{4.5in}|}
\hline
Keyword & ismetal \\
\hline
Kind & scalar \\
\hline
Type & Logical \\
\hline
Default & False \\
\hline
Deprecated & No \\
\hline
\end{longtable}
 Toggle whether the system should be treated as a metal or as an insulator.

\paragraph{Schema}
\begin{Verbatim}
False |X| L
\end{Verbatim}

\subsection{multiplicity}
\label{kw:multiplicity}
\begin{longtable}{|p{1.5in}|p{4.5in}|}
\hline
Keyword & multiplicity \\
\hline
Kind & scalar \\
\hline
Type & Integer \\
\hline
Default & 0 \\
\hline
Deprecated & No \\
\hline
\end{longtable}
 Multiplicity of the system.

\paragraph{Schema}
\begin{Verbatim}
0 |X| I
\end{Verbatim}

\subsection{number\_of\_atoms}
\label{kw:number-of-atoms}
\begin{longtable}{|p{1.5in}|p{4.5in}|}
\hline
Keyword & number\_of\_atoms \\
\hline
Kind & scalar \\
\hline
Type & Integer \\
\hline
Default & Required / No default \\
\hline
Deprecated & No \\
\hline
\end{longtable}
 Number of atoms in the system

\paragraph{Schema}
\begin{Verbatim}
|X| I
\end{Verbatim}

\subsection{polarization}
\label{kw:polarization}
\begin{longtable}{|p{1.5in}|p{4.5in}|}
\hline
Keyword & polarization \\
\hline
Kind & repeating-record \\
\hline
Type & Repeating record with columns: Float, Float, Float \\
\hline
Default & Required / No default \\
\hline
Deprecated & No \\
\hline
\end{longtable}
 Set any number of different polarization directions to calculate in cartesian coordinates.
 Example: set to x, y, z, i.e.,
 polarization{
    1.0 0.0 0.0
    0.0 1.0 0.0
    0.0 0.0 1.0

 If set to 0.0 0.0 0.0 (default), only polarization average will be calculated.
}

\paragraph{Schema}
\begin{Verbatim}
    0.0 0.0 0.0 |X| F F F
    |X| ...
\end{Verbatim}

\subsection{species}
\label{kw:species}
\begin{longtable}{|p{1.5in}|p{4.5in}|}
\hline
Keyword & species \\
\hline
Kind & vector \\
\hline
Type & Vector: String, Integer \\
\hline
Default & Required / No default \\
\hline
Deprecated & No \\
\hline
\end{longtable}
 chemical species and number of each species

\paragraph{Schema}
\begin{Verbatim}
|X| S I
\end{Verbatim}

\subsection{spectral\_broadening}
\label{kw:spectral-broadening}
\begin{longtable}{|p{1.5in}|p{4.5in}|}
\hline
Keyword & spectral\_broadening \\
\hline
Kind & scalar \\
\hline
Type & Float \\
\hline
Default & 0.0 \\
\hline
Deprecated & No \\
\hline
\end{longtable}
 Broadening used in some spectroscopic methods.

\paragraph{Schema}
\begin{Verbatim}
0.0 |X| F
\end{Verbatim}

\subsection{spin\_moment}
\label{kw:spin-moment}
\begin{longtable}{|p{1.5in}|p{4.5in}|}
\hline
Keyword & spin\_moment \\
\hline
Kind & repeating-record \\
\hline
Type & Repeating record with columns: Integer, Float \\
\hline
Default & Required / No default \\
\hline
Deprecated & No \\
\hline
\end{longtable}
 spin moments of atomic types
 set spin moment for ferromagnetic systems.
 Only allows one moment per chemical element.
 Z spin\_moment

\paragraph{Schema}
\begin{Verbatim}
|X| I F
|X| ...
\end{Verbatim}

\subsection{supercell\_dimensions}
\label{kw:supercell-dimensions}
\begin{longtable}{|p{1.5in}|p{4.5in}|}
\hline
Keyword & supercell\_dimensions \\
\hline
Kind & vector \\
\hline
Type & Vector: Integer, Integer, Integer \\
\hline
Default & Required / No default \\
\hline
Deprecated & No \\
\hline
\end{longtable}
 supercell dimensions (number of cells in each direction).

\paragraph{Schema}
\begin{Verbatim}
|X| I I I
\end{Verbatim}

\subsection{vasp\_outcar}
\label{kw:vasp-outcar}
\begin{longtable}{|p{1.5in}|p{4.5in}|}
\hline
Keyword & vasp\_outcar \\
\hline
Kind & scalar \\
\hline
Type & String \\
\hline
Default & Required / No default \\
\hline
Deprecated & No \\
\hline
\end{longtable}
 specify outcar file to use as input to corvus.

\paragraph{Schema}
\begin{Verbatim}
|X| S
\end{Verbatim}

\subsection{vasp\_snapshot}
\label{kw:vasp-snapshot}
\begin{longtable}{|p{1.5in}|p{4.5in}|}
\hline
Keyword & vasp\_snapshot \\
\hline
Kind & scalar \\
\hline
Type & Integer \\
\hline
Default & -1 \\
\hline
Deprecated & No \\
\hline
\end{longtable}
 specify index of structure (snapshot if MD run) to use.

\paragraph{Schema}
\begin{Verbatim}
-1 |X| I
\end{Verbatim}

\subsection{vasp\_xdatcar}
\label{kw:vasp-xdatcar}
\begin{longtable}{|p{1.5in}|p{4.5in}|}
\hline
Keyword & vasp\_xdatcar \\
\hline
Kind & scalar \\
\hline
Type & String \\
\hline
Default & Required / No default \\
\hline
Deprecated & No \\
\hline
\end{longtable}
 specify xdatcar file to use as input to corvus.

\paragraph{Schema}
\begin{Verbatim}
|X| S
\end{Verbatim}

\subsection{vasp\_xml}
\label{kw:vasp-xml}
\begin{longtable}{|p{1.5in}|p{4.5in}|}
\hline
Keyword & vasp\_xml \\
\hline
Kind & scalar \\
\hline
Type & String \\
\hline
Default & Required / No default \\
\hline
Deprecated & No \\
\hline
\end{longtable}
 specify vasp output xml file to use as input to corvus.

\paragraph{Schema}
\begin{Verbatim}
|X| S
\end{Verbatim}

\subsection{xyz\_input}
\label{kw:xyz-input}
\begin{longtable}{|p{1.5in}|p{4.5in}|}
\hline
Keyword & xyz\_input \\
\hline
Kind & scalar \\
\hline
Type & String \\
\hline
Default & Required / No default \\
\hline
Deprecated & No \\
\hline
\end{longtable}
 xyz input file name

\paragraph{Schema}
\begin{Verbatim}
|X| S
\end{Verbatim}

\subsection{xyz\_snapshot}
\label{kw:xyz-snapshot}
\begin{longtable}{|p{1.5in}|p{4.5in}|}
\hline
Keyword & xyz\_snapshot \\
\hline
Kind & scalar \\
\hline
Type & Integer \\
\hline
Default & 1 \\
\hline
Deprecated & No \\
\hline
\end{longtable}
 for multiple structure xyz files, use this snapshot.

\paragraph{Schema}
\begin{Verbatim}
1 |X| I
\end{Verbatim}

\section{Abinit}

\subsection{abinit.ng2qpt}
\label{kw:abinit-ng2qpt}
\begin{longtable}{|p{1.5in}|p{4.5in}|}
\hline
Keyword & abinit.ng2qpt \\
\hline
Kind & vector \\
\hline
Type & Vector: Integer, Integer, Integer \\
\hline
Default & Required / No default \\
\hline
Deprecated & No \\
\hline
\end{longtable}
 Equivalent to Abinit's (anaddb) ng2qpt input variable: Defines the number of
 q-point perturbations in each direction of the finer grid for density
 functional perturbation theory (DFPT) simulations. Please refer to the
 Abinit manual for further details.

\paragraph{Schema}
\begin{Verbatim}
|X| I I I
\end{Verbatim}

\subsection{abinit.ngqpt}
\label{kw:abinit-ngqpt}
\begin{longtable}{|p{1.5in}|p{4.5in}|}
\hline
Keyword & abinit.ngqpt \\
\hline
Kind & vector \\
\hline
Type & Vector: Integer, Integer, Integer \\
\hline
Default & Required / No default \\
\hline
Deprecated & No \\
\hline
\end{longtable}
 Equivalent to Abinit's ngqpt input variable: Defines the number of q-point
 perturbations in each direction of the grid for density functional
 perturbation theory (DFPT) simulations, and is that equivalent to the general
 "nqpoints" Corvus input variable. Please refer to the Abinit manual for
 further details.

 NOTE: This code-specific input variable will be replaced by the general
 "nqpoints" variable in the future.

\paragraph{Schema}
\begin{Verbatim}
|X| I I I
\end{Verbatim}

\subsection{abinit.verbatim}
\label{kw:abinit-verbatim}
\begin{longtable}{|p{1.5in}|p{4.5in}|}
\hline
Keyword & abinit.verbatim \\
\hline
Kind & scalar \\
\hline
Type & Free-form \\
\hline
Default & Required / No default \\
\hline
Deprecated & No \\
\hline
\end{longtable}
 The content of this input variable is passed as is (i. e. "verbatim") into the
 input of all Abinit calculations, and it is meant to help with any Abinit
 input that is currently not implemented in Corvus. Thus, it should be used
 carfully to avoid inconsistencies with the automatically generated input.
 Users should refed to the Abinit manual for information on this extra input.

\paragraph{Schema}
\begin{Verbatim}
|X| P
\end{Verbatim}

\section{DMDW}

\subsection{dmdw.ioflag}
\label{kw:dmdw-ioflag}
\begin{longtable}{|p{1.5in}|p{4.5in}|}
\hline
Keyword & dmdw.ioflag \\
\hline
Kind & scalar \\
\hline
Type & Free-form \\
\hline
Default & Required / No default \\
\hline
Deprecated & No \\
\hline
\end{longtable}
 Set the amount of output printed by DMDW. Possible values are:
   0: Terse, prints out only the desired result ($\sigma^2$, $u^2$, etc).
   1: Verbose, prints out the pole frequencies and weights, as well as
      estimates of the Einstein temperatures for each path.
 Please refer to the DMDW section of the Feff manual for further details.

 NOTE: The current format described below is temporary for compatibility with
 previous versions of Corvus. This will be changed to "Integer" in the future.

\paragraph{Schema}
\begin{Verbatim}
|X| P
\end{Verbatim}

\subsection{dmdw.nlanc}
\label{kw:dmdw-nlanc}
\begin{longtable}{|p{1.5in}|p{4.5in}|}
\hline
Keyword & dmdw.nlanc \\
\hline
Kind & scalar \\
\hline
Type & Free-form \\
\hline
Default & Required / No default \\
\hline
Deprecated & No \\
\hline
\end{longtable}
 Set the number of Lanczos poles (i. e. iterations) in DMDW. Larger values
 usually improve convergence of the target quantity. However, the number of
 poles should not exceed the dimensions of the subspace spanned by the
 projection of the path into the appropriate eigenmodes of the Hessian. This
 means that this variable should always be less than 3*N-6, where N is the
 number of atoms in the system. In practice, a value of 6-8 is sufficient to
 obtain converged mean square relative displacements (MSRD, or $\sigma^2$) for EXAFS.
 For crystallographic mean square displacements (MSD, or $u^2$) at least 16 poles
 are usually needed. Please refer to the DMDW section of the Feff manual for
 further details.

 NOTE: The current format described below is temporary for compatibility with
 previous versions of Corvus. This will be changed to "Integer" in the future.

\paragraph{Schema}
\begin{Verbatim}
|X| P
\end{Verbatim}

\subsection{dmdw.paths}
\label{kw:dmdw-paths}
\begin{longtable}{|p{1.5in}|p{4.5in}|}
\hline
Keyword & dmdw.paths \\
\hline
Kind & scalar \\
\hline
Type & Free-form \\
\hline
Default & Required / No default \\
\hline
Deprecated & No \\
\hline
\end{longtable}
 Set up the path descriptors that will generate the list of paths for which
 properties will be calculated. The list has the following form:
   <Number of descriptors>
     <Descriptor 1>
     <Descriptor 2>
    .
    .

 Each descriptor has the form:
   <Number of atoms in path> <Atom index 1> ... <Max. path length (Ang)>

 The atom indices can take the value 0 which acts as a wildcard for all atoms
 in the system.

 Examples:

 2
 2 1 0   3.0
 3 2 0 5 6.0

 This section defines 2 paths descriptors. The first one generates all paths
 with two atoms, starting in atom 1 and going to all other atoms in the
 systems, but subject to a maximum effective path length of 3.0 Ang. The
 second one generates all paths with three atoms, starting in atom 2 and
 ending in atom 5, while passing trhough all other atoms in the system, but
 with a maximum effective length of 6.0 Ang.
 Using the same syntax atoms can be selected to compute their $u^2$. For example
 the paths section

 1
 1 0 0.0

 will produce $u^2$ for all atoms in the system. Please refer to the DMDW
 section of the Feff manual for further details.

 NOTE: The current format described below is temporary for compatibility with
 previous versions of Corvus.

\paragraph{Schema}
\begin{Verbatim}
|X| P
\end{Verbatim}

\subsection{dmdw.tempgrid}
\label{kw:dmdw-tempgrid}
\begin{longtable}{|p{1.5in}|p{4.5in}|}
\hline
Keyword & dmdw.tempgrid \\
\hline
Kind & scalar \\
\hline
Type & Free-form \\
\hline
Default & Required / No default \\
\hline
Deprecated & No \\
\hline
\end{longtable}
 Set up the temperature grid to compute the thermal properties in DMDW. It has
 the following form:

   <Number of temperature> <Min. temp.> <Max. temp.>

 If the number of desired temperatures is just one, the maximum temperature
 input is not needed. Please refer to the DMDW section of the Feff manual for
 further details.

 NOTE: The current format described below is temporary for compatibility with
 previous versions of Corvus.

\paragraph{Schema}
\begin{Verbatim}
|X| P
\end{Verbatim}

\section{NWChem}

\subsection{nwchem.basis}
\label{kw:nwchem-basis}
\begin{longtable}{|p{1.5in}|p{4.5in}|}
\hline
Keyword & nwchem.basis \\
\hline
Kind & repeating-record \\
\hline
Type & Repeating record with columns: String \\
\hline
Default & Required / No default \\
\hline
Deprecated & No \\
\hline
\end{longtable}
 Set the Gaussian basis set to be used in the NWChem calculations. The format
 is the same as in NWChem:

 <Atom label> <Basis set name>
  .
  .

 Please refer to the NWChem manual for further details.

\paragraph{Schema}
\begin{Verbatim}
|X| S ...
|X| ...
\end{Verbatim}

\subsection{nwchem.charge}
\label{kw:nwchem-charge}
\begin{longtable}{|p{1.5in}|p{4.5in}|}
\hline
Keyword & nwchem.charge \\
\hline
Kind & scalar \\
\hline
Type & String \\
\hline
Default & Required / No default \\
\hline
Deprecated & No \\
\hline
\end{longtable}
 Set the net charge of the system for an NWChem DFT simulation.
 NOTE: This input variable will be superseded in the future by the more
 general "charge" variable.

 Please refer to the NWChem manual for further details.

\paragraph{Schema}
\begin{Verbatim}
|X| S
\end{Verbatim}

\subsection{nwchem.iter}
\label{kw:nwchem-iter}
\begin{longtable}{|p{1.5in}|p{4.5in}|}
\hline
Keyword & nwchem.iter \\
\hline
Kind & scalar \\
\hline
Type & String \\
\hline
Default & Required / No default \\
\hline
Deprecated & No \\
\hline
\end{longtable}
\paragraph{Schema}
\begin{Verbatim}
|X| S
\end{Verbatim}

\subsection{nwchem.mult}
\label{kw:nwchem-mult}
\begin{longtable}{|p{1.5in}|p{4.5in}|}
\hline
Keyword & nwchem.mult \\
\hline
Kind & scalar \\
\hline
Type & String \\
\hline
Default & Required / No default \\
\hline
Deprecated & No \\
\hline
\end{longtable}
 Set the multiplicity of the system for an NWChem DFT simulation.
 NOTE: This input variable will be superseded in the future by the more
 general "multiplicity" variable.

 Please refer to the NWChem manual for further details.

\paragraph{Schema}
\begin{Verbatim}
|X| S
\end{Verbatim}

\subsection{nwchem.qmd.dt\_nucl}
\label{kw:nwchem-qmd-dt-nucl}
\begin{longtable}{|p{1.5in}|p{4.5in}|}
\hline
Keyword & nwchem.qmd.dt\_nucl \\
\hline
Kind & scalar \\
\hline
Type & String \\
\hline
Default & Required / No default \\
\hline
Deprecated & No \\
\hline
\end{longtable}
 Set the time step for nuclear motion in a quantum MD simulation.

 Please refer to the NWChem manual for further details.

\paragraph{Schema}
\begin{Verbatim}
|X| S
\end{Verbatim}

\subsection{nwchem.qmd.nstep\_nucl}
\label{kw:nwchem-qmd-nstep-nucl}
\begin{longtable}{|p{1.5in}|p{4.5in}|}
\hline
Keyword & nwchem.qmd.nstep\_nucl \\
\hline
Kind & scalar \\
\hline
Type & String \\
\hline
Default & Required / No default \\
\hline
Deprecated & No \\
\hline
\end{longtable}
 Set the number of nuclear motion steps in a quantum MD simulation.

 Please refer to the NWChem manual for further details.

\paragraph{Schema}
\begin{Verbatim}
|X| S
\end{Verbatim}

\subsection{nwchem.qmd.print\_xyz}
\label{kw:nwchem-qmd-print-xyz}
\begin{longtable}{|p{1.5in}|p{4.5in}|}
\hline
Keyword & nwchem.qmd.print\_xyz \\
\hline
Kind & scalar \\
\hline
Type & String \\
\hline
Default & Required / No default \\
\hline
Deprecated & No \\
\hline
\end{longtable}
 Toggle printing of xyz coordinates in a quantum MD simulation.

 Please refer to the NWChem manual for further details.

\paragraph{Schema}
\begin{Verbatim}
|X| S
\end{Verbatim}

\subsection{nwchem.qmd.snapstep}
\label{kw:nwchem-qmd-snapstep}
\begin{longtable}{|p{1.5in}|p{4.5in}|}
\hline
Keyword & nwchem.qmd.snapstep \\
\hline
Kind & scalar \\
\hline
Type & String \\
\hline
Default & Required / No default \\
\hline
Deprecated & No \\
\hline
\end{longtable}
\paragraph{Schema}
\begin{Verbatim}
|X| S
\end{Verbatim}

\subsection{nwchem.qmd.targ\_temp}
\label{kw:nwchem-qmd-targ-temp}
\begin{longtable}{|p{1.5in}|p{4.5in}|}
\hline
Keyword & nwchem.qmd.targ\_temp \\
\hline
Kind & scalar \\
\hline
Type & String \\
\hline
Default & Required / No default \\
\hline
Deprecated & No \\
\hline
\end{longtable}
 Set the target temperature in a quantum MD simulation.

 Please refer to the NWChem manual for further details.

\paragraph{Schema}
\begin{Verbatim}
|X| S
\end{Verbatim}

\subsection{nwchem.qmd.thermostat}
\label{kw:nwchem-qmd-thermostat}
\begin{longtable}{|p{1.5in}|p{4.5in}|}
\hline
Keyword & nwchem.qmd.thermostat \\
\hline
Kind & scalar \\
\hline
Type & String \\
\hline
Default & Required / No default \\
\hline
Deprecated & No \\
\hline
\end{longtable}
 Set the thermostat type to be used in a quantum MD simulation.

 Please refer to the NWChem manual for further details.

\paragraph{Schema}
\begin{Verbatim}
|X| S
\end{Verbatim}

\subsection{nwchem.xas.alpha}
\label{kw:nwchem-xas-alpha}
\begin{longtable}{|p{1.5in}|p{4.5in}|}
\hline
Keyword & nwchem.xas.alpha \\
\hline
Kind & scalar \\
\hline
Type & String \\
\hline
Default & Required / No default \\
\hline
Deprecated & No \\
\hline
\end{longtable}
\paragraph{Schema}
\begin{Verbatim}
|X| S
\end{Verbatim}

\subsection{nwchem.xas.nroots}
\label{kw:nwchem-xas-nroots}
\begin{longtable}{|p{1.5in}|p{4.5in}|}
\hline
Keyword & nwchem.xas.nroots \\
\hline
Kind & scalar \\
\hline
Type & String \\
\hline
Default & Required / No default \\
\hline
Deprecated & No \\
\hline
\end{longtable}
\paragraph{Schema}
\begin{Verbatim}
|X| S
\end{Verbatim}

\subsection{nwchem.xas.vec}
\label{kw:nwchem-xas-vec}
\begin{longtable}{|p{1.5in}|p{4.5in}|}
\hline
Keyword & nwchem.xas.vec \\
\hline
Kind & scalar \\
\hline
Type & String \\
\hline
Default & Required / No default \\
\hline
Deprecated & No \\
\hline
\end{longtable}
\paragraph{Schema}
\begin{Verbatim}
|X| S
\end{Verbatim}

\subsection{nwchem.xas.xrayenergywin}
\label{kw:nwchem-xas-xrayenergywin}
\begin{longtable}{|p{1.5in}|p{4.5in}|}
\hline
Keyword & nwchem.xas.xrayenergywin \\
\hline
Kind & vector \\
\hline
Type & Vector: Float, Float \\
\hline
Default & Required / No default \\
\hline
Deprecated & No \\
\hline
\end{longtable}
\paragraph{Schema}
\begin{Verbatim}
|X| F F
\end{Verbatim}

\subsection{nwchem.xaselem}
\label{kw:nwchem-xaselem}
\begin{longtable}{|p{1.5in}|p{4.5in}|}
\hline
Keyword & nwchem.xaselem \\
\hline
Kind & scalar \\
\hline
Type & String \\
\hline
Default & Required / No default \\
\hline
Deprecated & No \\
\hline
\end{longtable}
\paragraph{Schema}
\begin{Verbatim}
|X| S
\end{Verbatim}

\subsection{nwchem.xc}
\label{kw:nwchem-xc}
\begin{longtable}{|p{1.5in}|p{4.5in}|}
\hline
Keyword & nwchem.xc \\
\hline
Kind & scalar \\
\hline
Type & String \\
\hline
Default & Required / No default \\
\hline
Deprecated & No \\
\hline
\end{longtable}
 Set the exchange correlation potential for an NWChem DFT simulation.
 NOTE: This input variable will be superseded in the future by the more
 general "xc" variable.

 Please refer to the NWChem manual for further details.

\paragraph{Schema}
\begin{Verbatim}
|X| S
\end{Verbatim}

\section{FEFF}

\subsection{debyetemp}
\label{kw:debyetemp}
\begin{longtable}{|p{1.5in}|p{4.5in}|}
\hline
Keyword & debyetemp \\
\hline
Kind & scalar \\
\hline
Type & Float \\
\hline
Default & Required / No default \\
\hline
Deprecated & No \\
\hline
\end{longtable}

 debyetemp
 {
 temp

 Specify the Debye Model temperature (in K) for the calculation of EXAFS DW
 factors.

\paragraph{Schema}
\begin{Verbatim}
|X| F
\end{Verbatim}

\subsection{dmdw\_nlanczos}
\label{kw:dmdw-nlanczos}
\begin{longtable}{|p{1.5in}|p{4.5in}|}
\hline
Keyword & dmdw\_nlanczos \\
\hline
Kind & scalar \\
\hline
Type & Integer \\
\hline
Default & 6 \\
\hline
Deprecated & No \\
\hline
\end{longtable}

 dmdw\_nlanczos
 {
 nLanczos\_Interations

 Specify the number of Lanczos iterations to be done for the calculation of
 EXAFS DW factors with the dynamical matrix method.

\paragraph{Schema}
\begin{Verbatim}
6 |X| I
\end{Verbatim}

\subsection{feff.absolute}
\label{kw:feff-absolute}
\begin{longtable}{|p{1.5in}|p{4.5in}|}
\hline
Keyword & feff.absolute \\
\hline
Kind & scalar \\
\hline
Type & Logical \\
\hline
Default & Required / No default \\
\hline
Deprecated & No \\
\hline
\end{longtable}
 feff.absolue{ UseAbsoluteUnits }
 Print spectrum in absolute units instead of normalizing.

\paragraph{Schema}
\begin{Verbatim}
|X| L
\end{Verbatim}

\subsection{feff.afolp}
\label{kw:feff-afolp}
\begin{longtable}{|p{1.5in}|p{4.5in}|}
\hline
Keyword & feff.afolp \\
\hline
Kind & scalar \\
\hline
Type & Float \\
\hline
Default & Required / No default \\
\hline
Deprecated & No \\
\hline
\end{longtable}
 feff.afolp{ folpx }
 Set maximum overlap for automatic overlap search.
 folpx - maximum overlap allowed.

\paragraph{Schema}
\begin{Verbatim}
|X| F
\end{Verbatim}

\subsection{feff.atoms}
\label{kw:feff-atoms}
\begin{longtable}{|p{1.5in}|p{4.5in}|}
\hline
Keyword & feff.atoms \\
\hline
Kind & repeating-record \\
\hline
Type & Repeating record with columns: Float, Float, Float, Integer \\
\hline
Default & Required / No default \\
\hline
Deprecated & No \\
\hline
\end{longtable}

 feff.atoms
 {
 x1  y1  z1 ipot1
 x2  y2  z2 ipot2
 .
 .
 .

 Specify atomic positions in cartesian coordinates (in Angstroms) and
 unique potential indices of each atom in the cluster, one atom per line.

\paragraph{Schema}
\begin{Verbatim}
|X| F F F I
|X| ...
\end{Verbatim}

\subsection{feff.cfaverage}
\label{kw:feff-cfaverage}
\begin{longtable}{|p{1.5in}|p{4.5in}|}
\hline
Keyword & feff.cfaverage \\
\hline
Kind & vector \\
\hline
Type & Vector: Integer, Integer, Float \\
\hline
Default & Required / No default \\
\hline
Deprecated & No \\
\hline
\end{longtable}
 Obsolete: Do not use.

\paragraph{Schema}
\begin{Verbatim}
|X| I I F
\end{Verbatim}

\subsection{feff.cgrid}
\label{kw:feff-cgrid}
\begin{longtable}{|p{1.5in}|p{4.5in}|}
\hline
Keyword & feff.cgrid \\
\hline
Kind & vector \\
\hline
Type & Vector: Float, Integer, Integer, Integer, Integer \\
\hline
Default & Required / No default \\
\hline
Deprecated & No \\
\hline
\end{longtable}
 Not implemented yet.

\paragraph{Schema}
\begin{Verbatim}
|X| F I I I I
\end{Verbatim}

\subsection{feff.chbroad}
\label{kw:feff-chbroad}
\begin{longtable}{|p{1.5in}|p{4.5in}|}
\hline
Keyword & feff.chbroad \\
\hline
Kind & scalar \\
\hline
Type & Integer \\
\hline
Default & Required / No default \\
\hline
Deprecated & No \\
\hline
\end{longtable}
 feff.chbroad{ ichbroad }
 Set method of calculating core-hole lifetime broadening.
 ichbroad -
            0: Calculate Green's function at energy with imaginary part equal
               to Gamma\_CH/2.
            1: Convolve final spectrum with Lorenzian of full width at
               half-max Gamma\_CH.

\paragraph{Schema}
\begin{Verbatim}
|X| I
\end{Verbatim}

\subsection{feff.chsh}
\label{kw:feff-chsh}
\begin{longtable}{|p{1.5in}|p{4.5in}|}
\hline
Keyword & feff.chsh \\
\hline
Kind & scalar \\
\hline
Type & Integer \\
\hline
Default & Required / No default \\
\hline
Deprecated & No \\
\hline
\end{longtable}
 feff.chsh{ ich }
 Correct chemical shift.
 ich = 0 or 1.

\paragraph{Schema}
\begin{Verbatim}
|X| I
\end{Verbatim}

\subsection{feff.chwidth}
\label{kw:feff-chwidth}
\begin{longtable}{|p{1.5in}|p{4.5in}|}
\hline
Keyword & feff.chwidth \\
\hline
Kind & scalar \\
\hline
Type & Float \\
\hline
Default & Required / No default \\
\hline
Deprecated & No \\
\hline
\end{longtable}
 feff.chwidth{ gamma\_ch }
 Set core-hole width instead of using FEFF's internal table of widths.

\paragraph{Schema}
\begin{Verbatim}
|X| F
\end{Verbatim}

\subsection{feff.cif}
\label{kw:feff-cif}
\begin{longtable}{|p{1.5in}|p{4.5in}|}
\hline
Keyword & feff.cif \\
\hline
Kind & scalar \\
\hline
Type & String \\
\hline
Default & Required / No default \\
\hline
Deprecated & No \\
\hline
\end{longtable}
 Specify a cif input file. Dangerous. Use cif\_input instead.

\paragraph{Schema}
\begin{Verbatim}
|X| S
\end{Verbatim}

\subsection{feff.compton}
\label{kw:feff-compton}
\begin{longtable}{|p{1.5in}|p{4.5in}|}
\hline
Keyword & feff.compton \\
\hline
Kind & vector \\
\hline
Type & Vector: Float, Integer, Integer \\
\hline
Default & Required / No default \\
\hline
Deprecated & No \\
\hline
\end{longtable}
 Not implemented yet.

\paragraph{Schema}
\begin{Verbatim}
|X| F I I
\end{Verbatim}

\subsection{feff.config}
\label{kw:feff-config}
\begin{longtable}{|p{1.5in}|p{4.5in}|}
\hline
Keyword & feff.config \\
\hline
Kind & repeating-record \\
\hline
Type & Repeating record with columns: String \\
\hline
Default & Required / No default \\
\hline
Deprecated & No \\
\hline
\end{longtable}
 Not implemented yet.

\paragraph{Schema}
\begin{Verbatim}
|X| S
|X| ...
\end{Verbatim}

\subsection{feff.control}
\label{kw:feff-control}
\begin{longtable}{|p{1.5in}|p{4.5in}|}
\hline
Keyword & feff.control \\
\hline
Kind & vector \\
\hline
Type & Vector: Integer, Integer, Integer, Integer, Integer, Integer \\
\hline
Default & Required / No default \\
\hline
Deprecated & No \\
\hline
\end{longtable}

 feff.control{ ipot ixsph ifms ipaths igenfmt iff2x }

 feff.control lets you run one or more of the feff program modules
 separately. There is a switch for each of six parts of feff:
 0 means not to run that module, 1 means to run it.

\paragraph{Schema}
\begin{Verbatim}
|X| I I I I I I
\end{Verbatim}

\subsection{feff.coordinates}
\label{kw:feff-coordinates}
\begin{longtable}{|p{1.5in}|p{4.5in}|}
\hline
Keyword & feff.coordinates \\
\hline
Kind & scalar \\
\hline
Type & Integer \\
\hline
Default & Required / No default \\
\hline
Deprecated & No \\
\hline
\end{longtable}
 feff.coordinates{ i }
 i must be an integer from 1 through 6. It specifies the units of the atoms of the unit cell given
 in the ATOMS card for reciprocal space calculations. If the card is omitted, the default value
 icoord = 3 is assumed. FIX check this
 1. Cartesian coordinates, Angstrom units. Like feff - you can copy from a real-space
 feff.inp file if your lattice vectors coincide with atoms in that feff.inp file.
 2. Cartesian coordinates, fractional units (i.e., fractions of the lattice vectors ; should be
 numbers between 0 and 1). Similar to feff.
 3. Cartesian coordinates, units are fractional with respect to FIRST lattice vector. Like
 SPRKKR. (default)
 4. Given in lattice coordinates, in fractional units. Like WIEN2k (but beware of some
 ‘funny’ lattice types, e.g. rhombohedral, in WIEN2k case.struct if you’re copy-pasting )
 5. Given in lattice coordinates, units are fractional with respect to FIRST lattice vector.
 6. Given in lattice coordinates, Angstrom units.

\paragraph{Schema}
\begin{Verbatim}
|X| I
\end{Verbatim}

\subsection{feff.corehole}
\label{kw:feff-corehole}
\begin{longtable}{|p{1.5in}|p{4.5in}|}
\hline
Keyword & feff.corehole \\
\hline
Kind & scalar \\
\hline
Type & String \\
\hline
Default & Required / No default \\
\hline
Deprecated & No \\
\hline
\end{longtable}
 feff.corehole{ CoreHoleApproximation }
 Set method for calculating corehole interaction.
 CoreHoleApproximation -
                         None: Do not use a corehole.
                         FSR : Final-state rule
                         RPA : Use RPA dielectric function to calculate
                               screening of core-hole.

\paragraph{Schema}
\begin{Verbatim}
|X| S
\end{Verbatim}

\subsection{feff.corrections}
\label{kw:feff-corrections}
\begin{longtable}{|p{1.5in}|p{4.5in}|}
\hline
Keyword & feff.corrections \\
\hline
Kind & vector \\
\hline
Type & Vector: Float, Float \\
\hline
Default & Required / No default \\
\hline
Deprecated & No \\
\hline
\end{longtable}
 feff.corrections{ vrcorr vicorr }
 Correct the Fermi cutoff and broadening in the final spectrum calculation.
 vrcorr - Shift Fermi cutoff by -vrcorr.
 vicorr - Add extra Lorenzian broadening with half width at half max of vicorr.

\paragraph{Schema}
\begin{Verbatim}
|X| F F
\end{Verbatim}

\subsection{feff.corval}
\label{kw:feff-corval}
\begin{longtable}{|p{1.5in}|p{4.5in}|}
\hline
Keyword & feff.corval \\
\hline
Kind & vector \\
\hline
Type & Vector: Float, Float \\
\hline
Default & Required / No default \\
\hline
Deprecated & No \\
\hline
\end{longtable}
 feff.corval{ emin emax }
 The core-valence separation energy is found by scanning the DOS within an energy window.
 The emin parameter sets the lower bound (in eV) of this energy window and emax the upper.
 They are -70 eV and -20 eV by
 default. For some materials it is necessary to change these bounds, e.g. to -100 eV and
 -30 or -10 eV. For example,
 when SCF convergence is elusive because occupation numbers for one or more l-values are
 changing drastically between SCF iterations due to states moving above and below a poor
 estimate of the core-valence separation energy. We plan to replace the current mechanism by a
 more robust and automated algorithm, but in the meantime users can use the CORVAL card
 to handle some of these difficult cases.

\paragraph{Schema}
\begin{Verbatim}
|X| F F
\end{Verbatim}

\subsection{feff.criteria}
\label{kw:feff-criteria}
\begin{longtable}{|p{1.5in}|p{4.5in}|}
\hline
Keyword & feff.criteria \\
\hline
Kind & vector \\
\hline
Type & Vector: Float, Float \\
\hline
Default & Required / No default \\
\hline
Deprecated & No \\
\hline
\end{longtable}
 feff.criteria{ critcw critpw }
 Cutoff criteria for the path expansion filtering.
 critcw - tolerance for curved wave expansion
 critpw - tolerance for initial plane wave approximation

\paragraph{Schema}
\begin{Verbatim}
|X| F F
\end{Verbatim}

\subsection{feff.crpa}
\label{kw:feff-crpa}
\begin{longtable}{|p{1.5in}|p{4.5in}|}
\hline
Keyword & feff.crpa \\
\hline
Kind & vector \\
\hline
Type & Vector: Integer, Float \\
\hline
Default & Required / No default \\
\hline
Deprecated & No \\
\hline
\end{longtable}
 Not currently working.

\paragraph{Schema}
\begin{Verbatim}
|X| I F
\end{Verbatim}

\subsection{feff.danes}
\label{kw:feff-danes}
\begin{longtable}{|p{1.5in}|p{4.5in}|}
\hline
Keyword & feff.danes \\
\hline
Kind & vector \\
\hline
Type & Vector: Float, Float, Float \\
\hline
Default & Required / No default \\
\hline
Deprecated & No \\
\hline
\end{longtable}
 feff.danes{ xkmax xkstep estep }
 Calculate diffraction anomalous near edge structure.
 xkmax  - calculate up to k = xkmax
 xkstep - use steps of size xkstep
 estep  - near the edge, use steps calculated from estep

\paragraph{Schema}
\begin{Verbatim}
|X| F F F
\end{Verbatim}

\subsection{feff.debye}
\label{kw:feff-debye}
\begin{longtable}{|p{1.5in}|p{4.5in}|}
\hline
Keyword & feff.debye \\
\hline
Kind & scalar \\
\hline
Type & String \\
\hline
Default & Required / No default \\
\hline
Deprecated & No \\
\hline
\end{longtable}
 feff.debye{ Temperature Debye-Temperature idwopt }
 Set temperature and Debye temperature for calculations of EXAFS Debye-Waller
 factors.
   idwopt - set method for calculating DW factors.
 NOTE: This input variable is not fully implemented yet.

\paragraph{Schema}
\begin{Verbatim}
|X| S ...
\end{Verbatim}

\subsection{feff.density}
\label{kw:feff-density}
\begin{longtable}{|p{1.5in}|p{4.5in}|}
\hline
Keyword & feff.density \\
\hline
Kind & block \\
\hline
Type & Multi-line heterogeneous block \\
\hline
Default & Required / No default \\
\hline
Deprecated & No \\
\hline
\end{longtable}
 Not implemented yet.

\paragraph{Schema}
\begin{Verbatim}
# den_type, grid_type, outfile
|X| S S S
# origin
|X| F F F
# number of points in each direction
|X| I I I
# vectors along which to step
|X| F F F
|X| F F F
|X| F F F
\end{Verbatim}

\subsection{feff.dims}
\label{kw:feff-dims}
\begin{longtable}{|p{1.5in}|p{4.5in}|}
\hline
Keyword & feff.dims \\
\hline
Kind & vector \\
\hline
Type & Vector: Integer, Integer \\
\hline
Default & Required / No default \\
\hline
Deprecated & No \\
\hline
\end{longtable}
 feff.dims{ nmax lmax }
 Set maximum dimensions for fms calculations.
 nmax - maximum number of atoms in fms matrix.
 lmax - maximum angular momentum of fms matrix.

\paragraph{Schema}
\begin{Verbatim}
|X| I I
\end{Verbatim}

\subsection{feff.dmdw}
\label{kw:feff-dmdw}
\begin{longtable}{|p{1.5in}|p{4.5in}|}
\hline
Keyword & feff.dmdw \\
\hline
Kind & vector \\
\hline
Type & Vector: String, Integer, Integer, Integer \\
\hline
Default & Required / No default \\
\hline
Deprecated & No \\
\hline
\end{longtable}
 NOTE: This input variable is not fully implemented yet.

\paragraph{Schema}
\begin{Verbatim}
|X| S I I I
\end{Verbatim}

\subsection{feff.edge}
\label{kw:feff-edge}
\begin{longtable}{|p{1.5in}|p{4.5in}|}
\hline
Keyword & feff.edge \\
\hline
Kind & scalar \\
\hline
Type & String \\
\hline
Default & Required / No default \\
\hline
Deprecated & No \\
\hline
\end{longtable}
 feff.edge{ edgeLabel }
 Set edge to calculate for XAS, XES, EELS, and other related spectra.

\paragraph{Schema}
\begin{Verbatim}
|X| S ...
\end{Verbatim}

\subsection{feff.egapse}
\label{kw:feff-egapse}
\begin{longtable}{|p{1.5in}|p{4.5in}|}
\hline
Keyword & feff.egapse \\
\hline
Kind & scalar \\
\hline
Type & Float \\
\hline
Default & Required / No default \\
\hline
Deprecated & No \\
\hline
\end{longtable}
 feff.egapse{ egap }
 Use this gap energy when applying the many-pole self-energy.

\paragraph{Schema}
\begin{Verbatim}
|X| F
\end{Verbatim}

\subsection{feff.egrid}
\label{kw:feff-egrid}
\begin{longtable}{|p{1.5in}|p{4.5in}|}
\hline
Keyword & feff.egrid \\
\hline
Kind & repeating-record \\
\hline
Type & Repeating record with columns: String, String, Float, Float \\
\hline
Default & Required / No default \\
\hline
Deprecated & No \\
\hline
\end{longtable}
 This card can be used to customize the energy grid. The EGRID card is followed by lines
 specifying the type of grid, minimum and maximum values for the grid, and the grid step, i.e.

 grid\_type grid\_min grid\_max grid\_step

 The grid type parameter is a string that can take the values e\_grid, k\_grid, or exp\_grid. When
 using the e\_grid or k\_grid grid types, grid\_min, grid\_max, and grid\_step are given in eV or $\AA^{-1}$
 respectively. For the exp\_grid type, grid\_min and grid\_max are the minimum and maximum
 grid values in eV, and grid\_step is the exponential, i.e.

  $$E_i = grid{\_}min + (grid{\_}max-grid{\_}min)[exp(grid{\_}step \times i) - 1].$$

 A fourth grid type user grid is also available for feff. user grid is followed by an arbitrary
 number of lines, each specifying an energy point in eV , i.e.,
 user\_grid
 0.1
 1.5
 3.45
 6.0
 .
 .
 .

\paragraph{Schema}
\begin{Verbatim}
|X| S S F F
|X| ...
\end{Verbatim}

\subsection{feff.ellipticity}
\label{kw:feff-ellipticity}
\begin{longtable}{|p{1.5in}|p{4.5in}|}
\hline
Keyword & feff.ellipticity \\
\hline
Kind & vector \\
\hline
Type & Vector: Float, Float, Float, Float \\
\hline
Default & Required / No default \\
\hline
Deprecated & No \\
\hline
\end{longtable}
 feff.ellipticity{ elpty x y z }
 Set the ellipticity of the x-rays.
 This card is used with the feff.polarization.
   elpty   - ratio of amplitudes of electric field in the two orthogonal
             directions of elliptically polarized light. Only the absolute
             value of the ratio is important for nonmagnetic materials. The
             present code can distinguish left- and right-circular polarization
             only with the feff.xmcd or feff.xncd. A zero value of the
             ellipticity corresponds to linear polarization, and unity to
             circular polarization. The default value is zero.
   x, y, z - coordinates of any nonzero vector in the direction of the
             incident beam. This vector should be approximately normal to the
             polarization vector

\paragraph{Schema}
\begin{Verbatim}
|X| F F F F
\end{Verbatim}

\subsection{feff.elnes}
\label{kw:feff-elnes}
\begin{longtable}{|p{1.5in}|p{4.5in}|}
\hline
Keyword & feff.elnes \\
\hline
Kind & repeating-record \\
\hline
Type & Repeating record with columns: Float, Float, Float \\
\hline
Default & Required / No default \\
\hline
Deprecated & No \\
\hline
\end{longtable}
 This card is not implements yet.

\paragraph{Schema}
\begin{Verbatim}
|X| F F F
|X| S ...
|X| ...
# In future, fix so that we actually have these params, not just strings.
\end{Verbatim}

\subsection{feff.eps0}
\label{kw:feff-eps0}
\begin{longtable}{|p{1.5in}|p{4.5in}|}
\hline
Keyword & feff.eps0 \\
\hline
Kind & scalar \\
\hline
Type & Float \\
\hline
Default & Required / No default \\
\hline
Deprecated & No \\
\hline
\end{longtable}
 feff.eps0{ eps0 }
 Set dielectric constant used for many-pole self-energy calculations.

\paragraph{Schema}
\begin{Verbatim}
|X| F
\end{Verbatim}

\subsection{feff.equivalence}
\label{kw:feff-equivalence}
\begin{longtable}{|p{1.5in}|p{4.5in}|}
\hline
Keyword & feff.equivalence \\
\hline
Kind & scalar \\
\hline
Type & Integer \\
\hline
Default & 0 \\
\hline
Deprecated & No \\
\hline
\end{longtable}
 feff.equivalence{ ieq }
 This optional card is only active in combination with the CIF card. It tells feff how to
 generate potential types from the list of atom positions in the ‘cif’ file.
 If ieq = 1, the crystallographic equivalence as expressed in the ‘cif’ file is respected; that
 is, every separate line containing a generating atom position will lead to a separate potential
 type. This means that, e.g., in HOPG graphite, the two generating positions will give rise to
 two independent C potentials. This is also the default behavior if the EQUIVALENCE card is
 not specified.
 If ieq = 2, unique potentials are assigned based on atomic number Z only. That is, all
 C atoms will share a C potential and so on. This is how most feff calculations are run.
 Whether it is sensible or not to do this depends on the system and on the property one wishes
 to calculate. Keep in mind that feff is a muffin tin code, and may therefore be indifferent
 to certain differences between crystallographically inequivalent sites. On the other hand, if
 an element occurs in the crystal with different oxidation states, it may be necessary to assign
 separate potentials to these different types in order to describe the crystal properly and get
 accurate spectra.
 If ieq = 3, unique potentials are assigned based on atomic number Z and the first shell.
 This can be useful e.g. to treat larger systems with crystal defects, where only first neighbors
 of the defect need to be treated differently from all more distant atoms of a certain Z. (To be
 implemented.)
 If ieq = 4, a hybrid of methods 1 and 2 is used. That is, if the number of unique crystallo-
 graphic positions does not exceed a hard-coded limit (nphx=9 in the current version), they are
 treated with the correct crystallographic equivalence. If the number of unique crystallographi-
 cally inequivalent sites is larger, they get combined by atomic number Z. This ad hoc approach
 is a practical way of simply limiting the number of unique potentials. This makes sense be-
 cause, first of all, there are certain hardcoded limits that would require recompilation of the
 code, requiring more RAM memory and more work than a user may want to do.

\paragraph{Schema}
\begin{Verbatim}
0 |X| I
\end{Verbatim}

\subsection{feff.equivalence\_nmax}
\label{kw:feff-equivalence-nmax}
\begin{longtable}{|p{1.5in}|p{4.5in}|}
\hline
Keyword & feff.equivalence\_nmax \\
\hline
Kind & scalar \\
\hline
Type & Integer \\
\hline
Default & Required / No default \\
\hline
Deprecated & No \\
\hline
\end{longtable}
\paragraph{Schema}
\begin{Verbatim}
|X| I
\end{Verbatim}

\subsection{feff.exafs}
\label{kw:feff-exafs}
\begin{longtable}{|p{1.5in}|p{4.5in}|}
\hline
Keyword & feff.exafs \\
\hline
Kind & vector \\
\hline
Type & Vector: Float, Float \\
\hline
Default & Required / No default \\
\hline
Deprecated & No \\
\hline
\end{longtable}
 feff.exafs{ xkmax xkstep }
 Calculate EXAFS spectrum.
 xkmax  - calculate spectrum up to k = xkmax
 xkstep - calculate in steps of xkstep

\paragraph{Schema}
\begin{Verbatim}
|X| F F
\end{Verbatim}

\subsection{feff.exchange}
\label{kw:feff-exchange}
\begin{longtable}{|p{1.5in}|p{4.5in}|}
\hline
Keyword & feff.exchange \\
\hline
Kind & vector \\
\hline
Type & Vector: Integer, Float, Float, Integer \\
\hline
Default & Required / No default \\
\hline
Deprecated & No \\
\hline
\end{longtable}
 feff.exchange{ iexc vr vi iexc0 }
 Use feff.exchange to change the self-energy used in x-ray absorption
 calculations. ixc is an index specifying the potential model to use for the
 fine structure, and the optional ixc0 is the index of the model to use for
 the background function. The calculated potential can be corrected by adding
 a constant shift to the Fermi level given by vr0 and to a pure imaginary
 "optical" potential (i.e., uniform decay) given by vi0. Typical errors in
 Feff's self-consistent Fermi level estimate are about 1 eV.
 (The feff.corrections input is similar but allows the user to make small
 changes in vi0 and vr0 after the rest of the calculation is completed, for
 example in a fitting process.)

 Indices for the available exchange models:
   0 Hedin-Lundqvist + a constant imaginary part
   1 Dirac-Hara + a constant imaginary part
   2 ground state + a constant imaginary part
   3 Dirac-Hara + HL imag part + a constant imaginary part
   5 Partially nonlocal: Dirac-Fock for core + HL for valence electrons + a
     constant imaginary part

\paragraph{Schema}
\begin{Verbatim}
|X| I F F I
\end{Verbatim}

\subsection{feff.exelfs}
\label{kw:feff-exelfs}
\begin{longtable}{|p{1.5in}|p{4.5in}|}
\hline
Keyword & feff.exelfs \\
\hline
Kind & repeating-record \\
\hline
Type & Repeating record with columns: Float \\
\hline
Default & Required / No default \\
\hline
Deprecated & No \\
\hline
\end{longtable}
 This card is not implemented yet.

\paragraph{Schema}
\begin{Verbatim}
|X| F
|X| S
|X| ...
# In future, fix so that we actually have these params, not just strings.
\end{Verbatim}

\subsection{feff.extpot}
\label{kw:feff-extpot}
\begin{longtable}{|p{1.5in}|p{4.5in}|}
\hline
Keyword & feff.extpot \\
\hline
Kind & scalar \\
\hline
Type & Logical \\
\hline
Default & Required / No default \\
\hline
Deprecated & No \\
\hline
\end{longtable}
 Not currently workding. Do not use.

\paragraph{Schema}
\begin{Verbatim}
|X| L
\end{Verbatim}

\subsection{feff.fms}
\label{kw:feff-fms}
\begin{longtable}{|p{1.5in}|p{4.5in}|}
\hline
Keyword & feff.fms \\
\hline
Kind & vector \\
\hline
Type & Vector: Float, Integer, Integer, Float, Float, Float \\
\hline
Default & Required / No default \\
\hline
Deprecated & No \\
\hline
\end{longtable}
 feff.fms{ rfms lfms2 minv toler1 toler2 rdirec }
 Use full multiple-scattering.
 rfms   - radius of cluster to use for fms calculation.
 lfms1  - 0 for solids, 1 for molecules
 minv   - set algorithm for matrix inversion
          0: LU decomposition
          2: Lanczos
          3: Broyden (less reliable)
 toler1 - tolerance to stop recursion and Broyden algorithm
 toler2 - sets the matrix element of the Gt matrix to zero if its value is
          less than toler2
 rdirec - sets the matrix element of the Gt matrix to zero if the distance
          between atoms is larger than rdirec

\paragraph{Schema}
\begin{Verbatim}
|X| F I I F F F
\end{Verbatim}

\subsection{feff.folp}
\label{kw:feff-folp}
\begin{longtable}{|p{1.5in}|p{4.5in}|}
\hline
Keyword & feff.folp \\
\hline
Kind & repeating-record \\
\hline
Type & Repeating record with columns: Integer, Float \\
\hline
Default & Required / No default \\
\hline
Deprecated & No \\
\hline
\end{longtable}
 feff.folp{ ipot overlap }
 Set the overlap for unique potential ipot.

\paragraph{Schema}
\begin{Verbatim}
|X| I F
|X| ...
\end{Verbatim}

\subsection{feff.fprime}
\label{kw:feff-fprime}
\begin{longtable}{|p{1.5in}|p{4.5in}|}
\hline
Keyword & feff.fprime \\
\hline
Kind & vector \\
\hline
Type & Vector: Float, Float, Float \\
\hline
Default & Required / No default \\
\hline
Deprecated & No \\
\hline
\end{longtable}
 feff.fprime{ emin emax estep }
 emin  - minimum energy of calculation.
 emax  - maximum energy of calculation
 estep - energy step of calculation
 Calculate atomic scattering factor f'

\paragraph{Schema}
\begin{Verbatim}
|X| F F F
\end{Verbatim}

\subsection{feff.hole}
\label{kw:feff-hole}
\begin{longtable}{|p{1.5in}|p{4.5in}|}
\hline
Keyword & feff.hole \\
\hline
Kind & vector \\
\hline
Type & Vector: Integer, Float \\
\hline
Default & Required / No default \\
\hline
Deprecated & Yes \\
\hline
\end{longtable}
 DEPRECATED: Use feff.edge instead.

 Specify the edge using the hole number ihole, e.g.,
 K-edge : ihole = 1
 L1-edge: ihole = 2
 s02 specifies the EXAFS amplitude reduction factor, and should be set to 1.

\paragraph{Schema}
\begin{Verbatim}
|X| I F
\end{Verbatim}

\subsection{feff.hubbard}
\label{kw:feff-hubbard}
\begin{longtable}{|p{1.5in}|p{4.5in}|}
\hline
Keyword & feff.hubbard \\
\hline
Kind & vector \\
\hline
Type & Vector: Float, Float, Float, Integer \\
\hline
Default & Required / No default \\
\hline
Deprecated & No \\
\hline
\end{longtable}
 Not currently working.

\paragraph{Schema}
\begin{Verbatim}
|X| F F F I
\end{Verbatim}

\subsection{feff.interstitial}
\label{kw:feff-interstitial}
\begin{longtable}{|p{1.5in}|p{4.5in}|}
\hline
Keyword & feff.interstitial \\
\hline
Kind & vector \\
\hline
Type & Vector: Integer, Float, Float \\
\hline
Default & Required / No default \\
\hline
Deprecated & No \\
\hline
\end{longtable}
 feff.interstitial{ inters vtot vint\_scl }
 The construction of the interstitial potential and density may be changed by
 using this card. inters = 1 might be useful when only the surroundings of the
 absorbing atom are specified in 'feff.inp'. inters defines how to find the
 interstitial potential.
    inters=0: The interstitial potential is found by averaging over the
              entire extended cluster in 'feff.inp'. (Default)
    inters=1: the interstitial potential is found locally around the
              absorbing atom.
    vtot:     the volume per atom normalized by ratmin3
              (vtot=(volume per atom)/ratmin3), where ratmin is the shortest
              bond for the absorbing atom. This quantity defines the total
              volume (needed to calculate the interstitial density) of the
              extended cluster specified in 'feff.inp'. If vtot <= 0 then the
              total volume is calculated as a sum of Norman sphere volumes.
              Otherwise, total volume = nat * (vtot * ratmin3), where nat is
              the number of atoms in an extended cluster.

\paragraph{Schema}
\begin{Verbatim}
|X| I F F
\end{Verbatim}

\subsection{feff.ion}
\label{kw:feff-ion}
\begin{longtable}{|p{1.5in}|p{4.5in}|}
\hline
Keyword & feff.ion \\
\hline
Kind & repeating-record \\
\hline
Type & Repeating record with columns: Integer, Float \\
\hline
Default & Required / No default \\
\hline
Deprecated & No \\
\hline
\end{longtable}
 feff.ion{ ipot ionization }
 feff.ion ionizes all atoms with unique potential index ipot.
 NOTE: This input variable is not fully implemented yet.

\paragraph{Schema}
\begin{Verbatim}
|X| I F
|X| ...
\end{Verbatim}

\subsection{feff.iorder}
\label{kw:feff-iorder}
\begin{longtable}{|p{1.5in}|p{4.5in}|}
\hline
Keyword & feff.iorder \\
\hline
Kind & scalar \\
\hline
Type & Integer \\
\hline
Default & Required / No default \\
\hline
Deprecated & No \\
\hline
\end{longtable}
 feff.iorder{ iorder }
 Set order of approximation when calculating effective scattering matrices.

\paragraph{Schema}
\begin{Verbatim}
|X| I
\end{Verbatim}

\subsection{feff.jumprm}
\label{kw:feff-jumprm}
\begin{longtable}{|p{1.5in}|p{4.5in}|}
\hline
Keyword & feff.jumprm \\
\hline
Kind & scalar \\
\hline
Type & Logical \\
\hline
Default & Required / No default \\
\hline
Deprecated & No \\
\hline
\end{longtable}
 feff.jumprm{ UseRemoveJump }
 If true, smooth the jump between muffin tin potentials and interstitial.

\paragraph{Schema}
\begin{Verbatim}
|X| L
\end{Verbatim}

\subsection{feff.kmesh}
\label{kw:feff-kmesh}
\begin{longtable}{|p{1.5in}|p{4.5in}|}
\hline
Keyword & feff.kmesh \\
\hline
Kind & scalar \\
\hline
Type & Integer \\
\hline
Default & Required / No default \\
\hline
Deprecated & No \\
\hline
\end{longtable}
 Specify the kmesh.
 feff.kmesh nkp(x) [nkpy nkpz [ktype [usesym] ] ]
 This card specifies the mesh of k-vectors used to sample the full Brillouin Zone for the evaluation
 of Brillouin Zone integrals. Nkp is the number of points used in the full zone. It can be specified
 either as ”nkpx nkpy nkpz”, ”nkp”, or ”nkp 0 0”. If usesym = 1, the zone is reduced to its
 irreducible wedge using the symmetry options specified in file symfile, which must be present
 in the working directory. The k-mesh is constructed using the tetrahedron method of Bloechl
 et al., Phys. Rev. B, 1990. The parameter ktype is meant for time-saving only and means:
  ktype=1 : regular mesh of nkp points for all modules
  ktype=2 : use nkp points for ldos/fms and nkp/5 points for pot (significant time savings)
  ktype=3 : use nkp points for ldos/fms and nkp/5 points for pot (near edge) ; reduce nkp
 for all modules as we get away from near-edge (somewhat experimental)
 * use a k-mesh of 1000 points in the full BZ.
 feff.kmesh{ 1000 }
 * use a k-mesh of 10x5x3 points for a large, irregular cell.
 feff.kmesh{ 10 5 3 }
 * use a k-mesh of 1000 points and try to save time:
 feff.kmesh{ 1000 0 0 3 }

\paragraph{Schema}
\begin{Verbatim}
|X| I ...
\end{Verbatim}

\subsection{feff.lattice}
\label{kw:feff-lattice}
\begin{longtable}{|p{1.5in}|p{4.5in}|}
\hline
Keyword & feff.lattice \\
\hline
Kind & repeating-record \\
\hline
Type & Repeating record with columns: String \\
\hline
Default & Required / No default \\
\hline
Deprecated & No \\
\hline
\end{longtable}
 This card specifies the lattice. First, its type must be specified using a single letter : P for
 primitive, F for face centered cubic, I for body centered cubic, H for hexagonal. The following
 three lines give the three basis vectors in Carthesian Angstrom coordinates. They are multiplied
 by scale (e.g., 0.529177 to convert from bohr to Angstrom).
 feff.lattice{ P 3.18800
 ax ay az
 bx by bz
 cx cy cz

\paragraph{Schema}
\begin{Verbatim}
# Need to fix this so that all parameters are set in separate fields.
|X| S
|X| ...
\end{Verbatim}

\subsection{feff.ldec}
\label{kw:feff-ldec}
\begin{longtable}{|p{1.5in}|p{4.5in}|}
\hline
Keyword & feff.ldec \\
\hline
Kind & scalar \\
\hline
Type & Integer \\
\hline
Default & Required / No default \\
\hline
Deprecated & No \\
\hline
\end{longtable}
 For use in calculations of NRIXS. Not implemented yet.

\paragraph{Schema}
\begin{Verbatim}
|X| I
\end{Verbatim}

\subsection{feff.ldos}
\label{kw:feff-ldos}
\begin{longtable}{|p{1.5in}|p{4.5in}|}
\hline
Keyword & feff.ldos \\
\hline
Kind & vector \\
\hline
Type & Vector: Float, Float, Float, Integer \\
\hline
Default & Required / No default \\
\hline
Deprecated & No \\
\hline
\end{longtable}
 feff.ldos{ emin emax broadening npoints }
 Run angular momentum project density of states calculations.
 emin       - minimum energy of energy grid.
 emax       - maximum energy of energy grid.
 broadening - Lorenzian broadening to use in calculation.
 npoints    - number of energy points in grid.

\paragraph{Schema}
\begin{Verbatim}
|X| F F F I
\end{Verbatim}

\subsection{feff.lfms1}
\label{kw:feff-lfms1}
\begin{longtable}{|p{1.5in}|p{4.5in}|}
\hline
Keyword & feff.lfms1 \\
\hline
Kind & scalar \\
\hline
Type & Integer \\
\hline
Default & -1 \\
\hline
Deprecated & No \\
\hline
\end{longtable}
 set maximum angular momentum to use in potentials

\paragraph{Schema}
\begin{Verbatim}
-1 |X| I
\end{Verbatim}

\subsection{feff.lfms2}
\label{kw:feff-lfms2}
\begin{longtable}{|p{1.5in}|p{4.5in}|}
\hline
Keyword & feff.lfms2 \\
\hline
Kind & scalar \\
\hline
Type & Integer \\
\hline
Default & -1 \\
\hline
Deprecated & No \\
\hline
\end{longtable}
 set maximum angular momentum to use in FMS

\paragraph{Schema}
\begin{Verbatim}
-1 |X| I
\end{Verbatim}

\subsection{feff.ljmax}
\label{kw:feff-ljmax}
\begin{longtable}{|p{1.5in}|p{4.5in}|}
\hline
Keyword & feff.ljmax \\
\hline
Kind & scalar \\
\hline
Type & Integer \\
\hline
Default & Required / No default \\
\hline
Deprecated & No \\
\hline
\end{longtable}
 For use in calculations of NRIXS. Not implemented yet.

\paragraph{Schema}
\begin{Verbatim}
|X| I
\end{Verbatim}

\subsection{feff.magic}
\label{kw:feff-magic}
\begin{longtable}{|p{1.5in}|p{4.5in}|}
\hline
Keyword & feff.magic \\
\hline
Kind & scalar \\
\hline
Type & Float \\
\hline
Default & Required / No default \\
\hline
Deprecated & No \\
\hline
\end{longtable}
 This card is not implemented yet.

\paragraph{Schema}
\begin{Verbatim}
|X| F
\end{Verbatim}

\subsection{feff.mbconv}
\label{kw:feff-mbconv}
\begin{longtable}{|p{1.5in}|p{4.5in}|}
\hline
Keyword & feff.mbconv \\
\hline
Kind & scalar \\
\hline
Type & Logical \\
\hline
Default & Required / No default \\
\hline
Deprecated & Yes \\
\hline
\end{longtable}
 Deprecated: Should not be used.

\paragraph{Schema}
\begin{Verbatim}
|X| L
\end{Verbatim}

\subsection{feff.mpse}
\label{kw:feff-mpse}
\begin{longtable}{|p{1.5in}|p{4.5in}|}
\hline
Keyword & feff.mpse \\
\hline
Kind & vector \\
\hline
Type & Vector: Integer, Integer \\
\hline
Default & Required / No default \\
\hline
Deprecated & No \\
\hline
\end{longtable}
 feff.mpse{ ipl npl }
 Use many-pole model self-energy. This requires the loss function, which can
 be provided externally, or approximated using feff.opcons.
   ipl - Set method for self-energy inclusion:
          1: use an "average" self-energy which is applied to the whole
             system (default).
          2: use a density dependendent self-energy which is
             different at every point inside the muffin-tin radius.

\paragraph{Schema}
\begin{Verbatim}
|X| I I
\end{Verbatim}

\subsection{feff.multipole}
\label{kw:feff-multipole}
\begin{longtable}{|p{1.5in}|p{4.5in}|}
\hline
Keyword & feff.multipole \\
\hline
Kind & vector \\
\hline
Type & Vector: Integer, Integer \\
\hline
Default & Required / No default \\
\hline
Deprecated & No \\
\hline
\end{longtable}
 feff.multipole{ le2 l2lp }
 Set multipole expansion approximation. Specifies which multipole transitions
 to include in the calculations. The options are: only dipole (le2 = 0,
 default), dipole and magnetic dipole (le2 = 1), dipole and quadrupole
 (le2 = 2). This cannot be used with NRIXS and is not supported with EXELFS
 and ELNES. The additional field l2lp can be used to calculate individual
 dipolar contributions coming from L -> L + 1 (l2lp = 1) and from L -> L - 1
 (l2lp = -1). Notice that in polarization dependent data there is also a
 cross term, which is calculated only when l2lp = 0

\paragraph{Schema}
\begin{Verbatim}
|X| I I
\end{Verbatim}

\subsection{feff.nleg}
\label{kw:feff-nleg}
\begin{longtable}{|p{1.5in}|p{4.5in}|}
\hline
Keyword & feff.nleg \\
\hline
Kind & scalar \\
\hline
Type & Integer \\
\hline
Default & Required / No default \\
\hline
Deprecated & No \\
\hline
\end{longtable}
 feff.nleg{ nleg }
 Set maximum number of legs to use in path expansion. A Single scattering
 path has 2 legs.

\paragraph{Schema}
\begin{Verbatim}
|X| I
\end{Verbatim}

\subsection{feff.nohole}
\label{kw:feff-nohole}
\begin{longtable}{|p{1.5in}|p{4.5in}|}
\hline
Keyword & feff.nohole \\
\hline
Kind & scalar \\
\hline
Type & Logical \\
\hline
Default & Required / No default \\
\hline
Deprecated & Yes \\
\hline
\end{longtable}
 feffnohole{ UseNoHole }
 Do not use a hole in the calculation.
 DEPRECATED - Use feff.corehole instead.

\paragraph{Schema}
\begin{Verbatim}
|X| L
\end{Verbatim}

\subsection{feff.nrixs}
\label{kw:feff-nrixs}
\begin{longtable}{|p{1.5in}|p{4.5in}|}
\hline
Keyword & feff.nrixs \\
\hline
Kind & repeating-record \\
\hline
Type & Repeating record with columns: Integer \\
\hline
Default & Required / No default \\
\hline
Deprecated & No \\
\hline
\end{longtable}
\paragraph{Schema}
\begin{Verbatim}
|X| I
|X| F ...
|X| ...
\end{Verbatim}

\subsection{feff.nstar}
\label{kw:feff-nstar}
\begin{longtable}{|p{1.5in}|p{4.5in}|}
\hline
Keyword & feff.nstar \\
\hline
Kind & scalar \\
\hline
Type & Logical \\
\hline
Default & Required / No default \\
\hline
Deprecated & No \\
\hline
\end{longtable}
 feff.nstar{ WriteNStar }
 Write nstar.dat with effective coordination number N*.

\paragraph{Schema}
\begin{Verbatim}
|X| L
\end{Verbatim}

\subsection{feff.numdens}
\label{kw:feff-numdens}
\begin{longtable}{|p{1.5in}|p{4.5in}|}
\hline
Keyword & feff.numdens \\
\hline
Kind & repeating-record \\
\hline
Type & Repeating record with columns: Integer, Float \\
\hline
Default & Required / No default \\
\hline
Deprecated & No \\
\hline
\end{longtable}
\paragraph{Schema}
\begin{Verbatim}
|X| I F
|X| ...
\end{Verbatim}

\subsection{feff.opcons}
\label{kw:feff-opcons}
\begin{longtable}{|p{1.5in}|p{4.5in}|}
\hline
Keyword & feff.opcons \\
\hline
Kind & scalar \\
\hline
Type & Logical \\
\hline
Default & Required / No default \\
\hline
Deprecated & No \\
\hline
\end{longtable}
 Calculate the dielectric function of the material within the
 atomic approximation. Very fast. Used to calculate loss.dat for
 use with feff.mpse self-energy calculations.

\paragraph{Schema}
\begin{Verbatim}
|X| L
\end{Verbatim}

\subsection{feff.overlap}
\label{kw:feff-overlap}
\begin{longtable}{|p{1.5in}|p{4.5in}|}
\hline
Keyword & feff.overlap \\
\hline
Kind & scalar \\
\hline
Type & Integer \\
\hline
Default & Required / No default \\
\hline
Deprecated & No \\
\hline
\end{longtable}
 feff.overlap can be used to construct approximate overlapped atom potentials
 when atomic coordinates are not known or specified.
 NOTE: This input variable is not fully implemented yet.

\paragraph{Schema}
\begin{Verbatim}
|X| I
\end{Verbatim}

\subsection{feff.pcriteria}
\label{kw:feff-pcriteria}
\begin{longtable}{|p{1.5in}|p{4.5in}|}
\hline
Keyword & feff.pcriteria \\
\hline
Kind & vector \\
\hline
Type & Vector: Float, Float \\
\hline
Default & Required / No default \\
\hline
Deprecated & No \\
\hline
\end{longtable}
 feff.pcriteria{ pcritk pcrith }
 Set criteria for filtering in pathfinder. The keep-criterion pcritk looks at
 the amplitude of chi (in the plane wave approximation) for the current path
 and compares it to a single scattering path of the same effective length.
 To set this value, consider the maximum degeneracy you expect and divide
 your plane wave criterion by this number. For example, in fcc Cu, typical
 degeneracies are 196 for paths with large r, and the minimum degeneracy is 6.
 So a keep criterion of 0.08% is appropriate for a pw criteria of 2.5%. The
 heap-criterion pcrith filters paths as the pathfinder puts all paths into a
 heap (a partially ordered data structure), then removes them in order of
 increasing total path length. Each path that is removed from the heap is
 modified and then considered again as part of the search algorithm. The heap
 filter is used to decide if a path has enough amplitude in it to be worth
 further consideration. If a path can be eliminated at this point, entire
 trees of derivative paths can be neglected, leading to enormous time savings.
 This test does not come into play until paths with at least 4 legs are being
 considered, so single scattering and triangular (2 and 3 legged) paths will
 always pass this test. Because only a small part of a path is used for
 this criterion, it is difficult to predict what appropriate values will be.

\paragraph{Schema}
\begin{Verbatim}
|X| F F
\end{Verbatim}

\subsection{feff.pmbse}
\label{kw:feff-pmbse}
\begin{longtable}{|p{1.5in}|p{4.5in}|}
\hline
Keyword & feff.pmbse \\
\hline
Kind & vector \\
\hline
Type & Vector: Integer, Integer, Integer, Integer \\
\hline
Default & Required / No default \\
\hline
Deprecated & No \\
\hline
\end{longtable}
\paragraph{Schema}
\begin{Verbatim}
|X| I I I I
\end{Verbatim}

\subsection{feff.polarization}
\label{kw:feff-polarization}
\begin{longtable}{|p{1.5in}|p{4.5in}|}
\hline
Keyword & feff.polarization \\
\hline
Kind & vector \\
\hline
Type & Vector: Float, Float, Float \\
\hline
Default & Required / No default \\
\hline
Deprecated & No \\
\hline
\end{longtable}
 feff.polarization{ x y z }
 Set polarization vector of x-rays in cartesian coordinates of input file.

\paragraph{Schema}
\begin{Verbatim}
|X| F F F
\end{Verbatim}

\subsection{feff.potentials}
\label{kw:feff-potentials}
\begin{longtable}{|p{1.5in}|p{4.5in}|}
\hline
Keyword & feff.potentials \\
\hline
Kind & repeating-record \\
\hline
Type & Repeating record with columns: Integer, Integer, String, Integer, Integer, Float \\
\hline
Default & Required / No default \\
\hline
Deprecated & No \\
\hline
\end{longtable}
 feff.potentials{
   ipot iz pot\_label lfms1 lfms2 stoichiometry
   ...
 }
 Set unique potentials for FEFF calculation.


\paragraph{Schema}
\begin{Verbatim}
|X| I I S I I F
|X| ...
\end{Verbatim}

\subsection{feff.potentials.spin}
\label{kw:feff-potentials-spin}
\begin{longtable}{|p{1.5in}|p{4.5in}|}
\hline
Keyword & feff.potentials.spin \\
\hline
Kind & repeating-record \\
\hline
Type & Repeating record with columns: Integer, Float \\
\hline
Default & Required / No default \\
\hline
Deprecated & No \\
\hline
\end{longtable}
 feff.potentials.spin{
  ipot spin\_moment
  ...
 }
 Set spin moments of atoms defined in potentials card

\paragraph{Schema}
\begin{Verbatim}
|X| I F
|X| ...
\end{Verbatim}

\subsection{feff.preps}
\label{kw:feff-preps}
\begin{longtable}{|p{1.5in}|p{4.5in}|}
\hline
Keyword & feff.preps \\
\hline
Kind & scalar \\
\hline
Type & Logical \\
\hline
Default & Required / No default \\
\hline
Deprecated & No \\
\hline
\end{longtable}
 feff.preps{ PrintEpsilon }
 Print dielectric function as calculated by feff.opcons.

\paragraph{Schema}
\begin{Verbatim}
|X| L
\end{Verbatim}

\subsection{feff.print}
\label{kw:feff-print}
\begin{longtable}{|p{1.5in}|p{4.5in}|}
\hline
Keyword & feff.print \\
\hline
Kind & vector \\
\hline
Type & Vector: Integer, Integer, Integer, Integer, Integer, Integer \\
\hline
Default & Required / No default \\
\hline
Deprecated & No \\
\hline
\end{longtable}
 feff.pring{ ipr\_pot ipr\_xsph ipr\_fms ipr\_path ipr\_genfmt ipr\_ff2x }
 Set print level for each module of FEFF:
   pot:
     0 write 'pot.bin' only
     1 add 'misc.dat'
     2 add 'potNN.dat'
     3 add 'atomNN.dat'

   xsph:
     0 write 'phase.bin' and 'xsect.bin' only
     1 add 'axafs.dat' and 'phase.dat'
     2 add 'phaseNN.dat' and 'phminNN.dat'
     3 add 'ratio.dat' (for XMCD normalization) and 'emesh.dat'.

   fms:
     0 write 'gg.bin'
     1 write 'gg.dat'

   path:
     0 write 'paths.dat' only
     1 add 'crit.dat'
     3 add 'fbeta' files (plane wave $|f(\beta)|$ approximations)
     5 Write only 'crit.dat' and do not write 'paths.dat'. (This is useful
       genfmt 0 write 'list.dat', and write 'feff.bin' with all paths with
       importance greater than or equal to two thirds of the curved wave
       importance criterion

   genfmt:
     0 write 'list.dat', and write 'feff.bin' with all paths with importance
       greater than or equal to two thirds of the curved wave importance
       criterion
     1 write all paths to 'feff.bin'

   ff2x:
     0 write 'chi.dat' and 'xmu.dat'
     2 add 'chipNNNN.dat' (chi(k) for each path individually)
     3 add 'feffNNNN.dat' and 'files.dat', and do not add 'chipNNNN.dat'
       files

\paragraph{Schema}
\begin{Verbatim}
|X| I I I I I I
\end{Verbatim}

\subsection{feff.rconv}
\label{kw:feff-rconv}
\begin{longtable}{|p{1.5in}|p{4.5in}|}
\hline
Keyword & feff.rconv \\
\hline
Kind & vector \\
\hline
Type & Vector: Float, String \\
\hline
Default & Required / No default \\
\hline
Deprecated & No \\
\hline
\end{longtable}
 Advanced: Print running convolution with the spectral function.

\paragraph{Schema}
\begin{Verbatim}
|X| F S
\end{Verbatim}

\subsection{feff.real}
\label{kw:feff-real}
\begin{longtable}{|p{1.5in}|p{4.5in}|}
\hline
Keyword & feff.real \\
\hline
Kind & scalar \\
\hline
Type & Logical \\
\hline
Default & Required / No default \\
\hline
Deprecated & No \\
\hline
\end{longtable}
 Advanced: Use real phase shifts.

\paragraph{Schema}
\begin{Verbatim}
|X| L
\end{Verbatim}

\subsection{feff.reciprocal}
\label{kw:feff-reciprocal}
\begin{longtable}{|p{1.5in}|p{4.5in}|}
\hline
Keyword & feff.reciprocal \\
\hline
Kind & scalar \\
\hline
Type & Logical \\
\hline
Default & Required / No default \\
\hline
Deprecated & No \\
\hline
\end{longtable}
 Work in reciprocal space (k-space calculation).

\paragraph{Schema}
\begin{Verbatim}
|X| L
\end{Verbatim}

\subsection{feff.restart}
\label{kw:feff-restart}
\begin{longtable}{|p{1.5in}|p{4.5in}|}
\hline
Keyword & feff.restart \\
\hline
Kind & scalar \\
\hline
Type & Logical \\
\hline
Default & Required / No default \\
\hline
Deprecated & No \\
\hline
\end{longtable}
 Not implemented yet.

\paragraph{Schema}
\begin{Verbatim}
|X| L
\end{Verbatim}

\subsection{feff.rgrid}
\label{kw:feff-rgrid}
\begin{longtable}{|p{1.5in}|p{4.5in}|}
\hline
Keyword & feff.rgrid \\
\hline
Kind & scalar \\
\hline
Type & Float \\
\hline
Default & Required / No default \\
\hline
Deprecated & No \\
\hline
\end{longtable}
\paragraph{Schema}
\begin{Verbatim}
|X| F
\end{Verbatim}

\subsection{feff.rhozzp}
\label{kw:feff-rhozzp}
\begin{longtable}{|p{1.5in}|p{4.5in}|}
\hline
Keyword & feff.rhozzp \\
\hline
Kind & scalar \\
\hline
Type & Logical \\
\hline
Default & Required / No default \\
\hline
Deprecated & No \\
\hline
\end{longtable}
 Not implemented yet.

\paragraph{Schema}
\begin{Verbatim}
|X| L
\end{Verbatim}

\subsection{feff.rixs}
\label{kw:feff-rixs}
\begin{longtable}{|p{1.5in}|p{4.5in}|}
\hline
Keyword & feff.rixs \\
\hline
Kind & vector \\
\hline
Type & Vector: Float, Float, Float \\
\hline
Default & Required / No default \\
\hline
Deprecated & No \\
\hline
\end{longtable}
 feff.RIXS {gam\_Ei gam\_El xmu}
 The RIXS card sets the parameters for the RIXS calculation. It must be used with a workflow
 or script to produce RIXS.
 gam\_Ei
    Broadening to use in the incident energy direction.

 gam\_El
    Broadening to use in the energy loss direction.

 xmu
    Fermi energy to use for the RIXS calculation, if not present, the fermi level calculated in
    the SCF step will be used.

\paragraph{Schema}
\begin{Verbatim}
|X| F F F
\end{Verbatim}

\subsection{feff.rlprint}
\label{kw:feff-rlprint}
\begin{longtable}{|p{1.5in}|p{4.5in}|}
\hline
Keyword & feff.rlprint \\
\hline
Kind & scalar \\
\hline
Type & Logical \\
\hline
Default & Required / No default \\
\hline
Deprecated & No \\
\hline
\end{longtable}
 Advanced used in RIXS calculations, but automated by corvus. Prints real scattering wavefunction $R_l$.

\paragraph{Schema}
\begin{Verbatim}
|X| L
\end{Verbatim}

\subsection{feff.rmultiplier}
\label{kw:feff-rmultiplier}
\begin{longtable}{|p{1.5in}|p{4.5in}|}
\hline
Keyword & feff.rmultiplier \\
\hline
Kind & scalar \\
\hline
Type & Float \\
\hline
Default & Required / No default \\
\hline
Deprecated & No \\
\hline
\end{longtable}
 feff.rmultiplyer{ rmult }
 Multiply coordinates of all atoms by rmult, expanding or contracting the
 system.

\paragraph{Schema}
\begin{Verbatim}
|X| F
\end{Verbatim}

\subsection{feff.rpath}
\label{kw:feff-rpath}
\begin{longtable}{|p{1.5in}|p{4.5in}|}
\hline
Keyword & feff.rpath \\
\hline
Kind & scalar \\
\hline
Type & Float \\
\hline
Default & Required / No default \\
\hline
Deprecated & No \\
\hline
\end{longtable}
 feff.rpath{ rmax }
 Set maximum path length for path expansion calculations of EXAFS, EXELFS,
 DAFS, etc.

\paragraph{Schema}
\begin{Verbatim}
|X| F
\end{Verbatim}

\subsection{feff.rphases}
\label{kw:feff-rphases}
\begin{longtable}{|p{1.5in}|p{4.5in}|}
\hline
Keyword & feff.rphases \\
\hline
Kind & scalar \\
\hline
Type & Logical \\
\hline
Default & Required / No default \\
\hline
Deprecated & No \\
\hline
\end{longtable}
 feff.rphases{ UseRealPhases }
 Set real phase shift approximation.
 UseRealPhases - use real phase shift if true.

\paragraph{Schema}
\begin{Verbatim}
|X| L
\end{Verbatim}

\subsection{feff.rsigma}
\label{kw:feff-rsigma}
\begin{longtable}{|p{1.5in}|p{4.5in}|}
\hline
Keyword & feff.rsigma \\
\hline
Kind & scalar \\
\hline
Type & Logical \\
\hline
Default & Required / No default \\
\hline
Deprecated & No \\
\hline
\end{longtable}
 feff.rsigma{ UseRealSelfEnergy }
 If true use only real part of self-energy.

\paragraph{Schema}
\begin{Verbatim}
|X| L
\end{Verbatim}

\subsection{feff.s02}
\label{kw:feff-s02}
\begin{longtable}{|p{1.5in}|p{4.5in}|}
\hline
Keyword & feff.s02 \\
\hline
Kind & scalar \\
\hline
Type & Float \\
\hline
Default & Required / No default \\
\hline
Deprecated & No \\
\hline
\end{longtable}
 feff.s02{ s02 }
 Set EXAFS amplitude reduction factor.

\paragraph{Schema}
\begin{Verbatim}
|X| F
\end{Verbatim}

\subsection{feff.scf}
\label{kw:feff-scf}
\begin{longtable}{|p{1.5in}|p{4.5in}|}
\hline
Keyword & feff.scf \\
\hline
Kind & vector \\
\hline
Type & Vector: Float, Integer, Integer, Float, Integer \\
\hline
Default & Required / No default \\
\hline
Deprecated & No \\
\hline
\end{longtable}
 feff.scf{ rfms1 lfms1 nscmt ca nmix }
 Use self-consistent field calculation of potentials.
   rfms1 - radius of cluster to use for scf calculations.
   lfms1 - 0 for solids, 1 for molecules
   nscmt - maximum number of iterations in SCF calculation.
   ca    - convergence factor for Broyden algorithm
   nmix  - number of initial iterations of mixing algorithm to use.

\paragraph{Schema}
\begin{Verbatim}
|X| F I I F I
\end{Verbatim}

\subsection{feff.scframp}
\label{kw:feff-scframp}
\begin{longtable}{|p{1.5in}|p{4.5in}|}
\hline
Keyword & feff.scframp \\
\hline
Kind & vector \\
\hline
Type & Vector: Float, Integer \\
\hline
Default & Required / No default \\
\hline
Deprecated & No \\
\hline
\end{longtable}
 Ramp scf radius up to final radius, starting from start\_radius,
 and taking n\_ramp steps.
 SCFRAMP start\_radius n\_ramp

\paragraph{Schema}
\begin{Verbatim}
|X| F I
\end{Verbatim}

\subsection{feff.screen}
\label{kw:feff-screen}
\begin{longtable}{|p{1.5in}|p{4.5in}|}
\hline
Keyword & feff.screen \\
\hline
Kind & repeating-record \\
\hline
Type & Repeating record with columns: String, String \\
\hline
Default & Required / No default \\
\hline
Deprecated & No \\
\hline
\end{longtable}
 The screen module, which calculates the RPA core hole potential, is a ’silent’ module: it has
 no obvious input but instead runs entirely on default values. Using the feff.SCREEN card you
 can change these default values. They will be written to an optional ‘screen.inp’ file (which
 you can also edit manually). The feff.SCREEN card can occur more than once in ‘feff.inp’; all
 entries will be applied to the calculation. parameter must be one of : ner (40), nei (20), maxl
 (4), irrh (1), iend (0), lfxc (0), emin (-40 eV), emax (0 eV), eimax (2 eV), ermin (0.001 eV),
 rfms (4.0), nrptx0 (251). For most calculations the default values given (between brackets) are
 fine. Occasionally we’ve changed rfms, maxl, or emin. Note that the screen is only active with
 feff.COREHOLE{ RPA }.
 * Set the cluster radius for the RPA potential calculation higher than the default of 4.0
 feff.COREHOLE{ RPA }
 feff.SCREEN{ rfms 5.0 }

\paragraph{Schema}
\begin{Verbatim}
# We will also want to improve this interface.
|X| S S
|X| ...
\end{Verbatim}

\subsection{feff.scxc}
\label{kw:feff-scxc}
\begin{longtable}{|p{1.5in}|p{4.5in}|}
\hline
Keyword & feff.scxc \\
\hline
Kind & scalar \\
\hline
Type & Integer \\
\hline
Default & Required / No default \\
\hline
Deprecated & No \\
\hline
\end{longtable}
\paragraph{Schema}
\begin{Verbatim}
|X| I
\end{Verbatim}

\subsection{feff.self}
\label{kw:feff-self}
\begin{longtable}{|p{1.5in}|p{4.5in}|}
\hline
Keyword & feff.self \\
\hline
Kind & scalar \\
\hline
Type & Logical \\
\hline
Default & Required / No default \\
\hline
Deprecated & No \\
\hline
\end{longtable}
 feff.self{ PrintSelfEnergy }
 Print out quasiparticle self-energy calculated during spectral-function
 convolution calculations.

\paragraph{Schema}
\begin{Verbatim}
|X| L
\end{Verbatim}

\subsection{feff.setedge}
\label{kw:feff-setedge}
\begin{longtable}{|p{1.5in}|p{4.5in}|}
\hline
Keyword & feff.setedge \\
\hline
Kind & scalar \\
\hline
Type & Logical \\
\hline
Default & Required / No default \\
\hline
Deprecated & No \\
\hline
\end{longtable}
 feff.setedge{ UseEdgeTable }
 Use table of experimental edge energies.

\paragraph{Schema}
\begin{Verbatim}
|X| L
\end{Verbatim}

\subsection{feff.sfconv}
\label{kw:feff-sfconv}
\begin{longtable}{|p{1.5in}|p{4.5in}|}
\hline
Keyword & feff.sfconv \\
\hline
Kind & scalar \\
\hline
Type & Logical \\
\hline
Default & Required / No default \\
\hline
Deprecated & No \\
\hline
\end{longtable}
 feff.sfconv{ UseSpectralFunctionConvolution }
 If true, convolve spectrum with spectral function to account for
 multi-electron excitations.

\paragraph{Schema}
\begin{Verbatim}
|X| L
\end{Verbatim}

\subsection{feff.sfse}
\label{kw:feff-sfse}
\begin{longtable}{|p{1.5in}|p{4.5in}|}
\hline
Keyword & feff.sfse \\
\hline
Kind & scalar \\
\hline
Type & Float \\
\hline
Default & Required / No default \\
\hline
Deprecated & No \\
\hline
\end{longtable}
 feff.sfse{ k }
 Print out off shell self-energy Sigma(k,E) calculated during
 spectral-function convolution calculations.

\paragraph{Schema}
\begin{Verbatim}
|X| F
\end{Verbatim}

\subsection{feff.sgroup}
\label{kw:feff-sgroup}
\begin{longtable}{|p{1.5in}|p{4.5in}|}
\hline
Keyword & feff.sgroup \\
\hline
Kind & scalar \\
\hline
Type & Logical \\
\hline
Default & Required / No default \\
\hline
Deprecated & No \\
\hline
\end{longtable}
 Specify space group by number (1 - 230).

\paragraph{Schema}
\begin{Verbatim}
|X| L
\end{Verbatim}

\subsection{feff.sig2}
\label{kw:feff-sig2}
\begin{longtable}{|p{1.5in}|p{4.5in}|}
\hline
Keyword & feff.sig2 \\
\hline
Kind & scalar \\
\hline
Type & Float \\
\hline
Default & Required / No default \\
\hline
Deprecated & No \\
\hline
\end{longtable}
 feff.sig2{ sig2 }
 Set a single Debye-Waller factor for all paths, exp{-2*sig2*k**2}.

\paragraph{Schema}
\begin{Verbatim}
|X| F
\end{Verbatim}

\subsection{feff.sig3}
\label{kw:feff-sig3}
\begin{longtable}{|p{1.5in}|p{4.5in}|}
\hline
Keyword & feff.sig3 \\
\hline
Kind & vector \\
\hline
Type & Vector: Float, Float \\
\hline
Default & Required / No default \\
\hline
Deprecated & No \\
\hline
\end{longtable}
 feff.sig3{ alphat thetae }
 Set first and third cumulants for single scattering paths:
   alphat - first cumulant
   thetae - third cumulant

\paragraph{Schema}
\begin{Verbatim}
|X| F F
\end{Verbatim}

\subsection{feff.siggk}
\label{kw:feff-siggk}
\begin{longtable}{|p{1.5in}|p{4.5in}|}
\hline
Keyword & feff.siggk \\
\hline
Kind & scalar \\
\hline
Type & Float \\
\hline
Default & Required / No default \\
\hline
Deprecated & No \\
\hline
\end{longtable}
 Not implemented yet.

\paragraph{Schema}
\begin{Verbatim}
|X| F
\end{Verbatim}

\subsection{feff.spin}
\label{kw:feff-spin}
\begin{longtable}{|p{1.5in}|p{4.5in}|}
\hline
Keyword & feff.spin \\
\hline
Kind & vector \\
\hline
Type & Vector: Integer, Float, Float, Float \\
\hline
Default & Required / No default \\
\hline
Deprecated & No \\
\hline
\end{longtable}
 feff.spin{ ispin sx sy sz }
 Specify the type of spin-dependent calculation for spin along the (x, y, z)
 direction, along the z-axis by default. The SPIN card is required for the
 calculation of all spin-dependent effects, including XMCD and SPXAS.
   ispin:
      2 spin-up SPXAS and LDOS
     -2 spin-down SPXAS and LDOS
      1 spin-up portion of XMCD calculations
     -1 spin-down portion of XMCD calculations

\paragraph{Schema}
\begin{Verbatim}
|X| I F F F
\end{Verbatim}

\subsection{feff.ss}
\label{kw:feff-ss}
\begin{longtable}{|p{1.5in}|p{4.5in}|}
\hline
Keyword & feff.ss \\
\hline
Kind & repeating-record \\
\hline
Type & Repeating record with columns: Integer, Integer, Float, Float \\
\hline
Default & Required / No default \\
\hline
Deprecated & No \\
\hline
\end{longtable}
 Set an

\paragraph{Schema}
\begin{Verbatim}
|X| I I F F
|X| ...
\end{Verbatim}

\subsection{feff.strfac}
\label{kw:feff-strfac}
\begin{longtable}{|p{1.5in}|p{4.5in}|}
\hline
Keyword & feff.strfac \\
\hline
Kind & vector \\
\hline
Type & Vector: Float, Float, Float \\
\hline
Default & Required / No default \\
\hline
Deprecated & No \\
\hline
\end{longtable}
 feff.strfac{ eta gmax rmax }
 This card gives three non-physical internal parameters for the calculation of the KKR structure
 factors : the Ewald parameter and a multiplicative cutoff factor for sums over reciprocal (gmax)
 and real space (rmax) sums. Multiplicative means the code makes a ’smart’ guess of a cutoff
 radius, but if one suspects something fishy is going on, one can here e.g. use gmax=2 to multiply
 this guess by 2. Eta is an absolute number. Given the stability of the Ewald algorithm, it
 shouldn’t be necessary to use this card. Its use is not recommended. Only active in combination
 with the feff.reciprocal card.

\paragraph{Schema}
\begin{Verbatim}
|X| F F F
\end{Verbatim}

\subsection{feff.symmetry}
\label{kw:feff-symmetry}
\begin{longtable}{|p{1.5in}|p{4.5in}|}
\hline
Keyword & feff.symmetry \\
\hline
Kind & scalar \\
\hline
Type & Logical \\
\hline
Default & Required / No default \\
\hline
Deprecated & No \\
\hline
\end{longtable}
 feff.symmetry{ UseSymmetry }
 If false, turn off symmetry considerations in path expansion.

\paragraph{Schema}
\begin{Verbatim}
|X| L
\end{Verbatim}

\subsection{feff.target}
\label{kw:feff-target}
\begin{longtable}{|p{1.5in}|p{4.5in}|}
\hline
Keyword & feff.target \\
\hline
Kind & scalar \\
\hline
Type & Integer \\
\hline
Default & Required / No default \\
\hline
Deprecated & No \\
\hline
\end{longtable}
 feff.target{ ic }
 Specifies the location of the absorber atom for reciprocal space calculations. It is entry ic of the
 feff.ATOMS card if an feff.ATOMS card and feff.LATTICE card are used. In conjunction with the cif\_input
 card it is entry ic the list of atoms as given in the ‘.cif’ file (i.e., a list of the crystallographically
 inequivalent atom positions in the unit cell). The target needs to be specified also for NOHOLE
 calculations. Note that this cannot be specified in the feff.POTENTIALS list because periodic
 boundary conditions would then produce an infinite number of core holes.
 * calculate a spectrum for the second atom in the feff.ATOMS list or CIF file.
 feff.TARGET{ 2 }

\paragraph{Schema}
\begin{Verbatim}
|X| I
\end{Verbatim}

\subsection{feff.tdlda}
\label{kw:feff-tdlda}
\begin{longtable}{|p{1.5in}|p{4.5in}|}
\hline
Keyword & feff.tdlda \\
\hline
Kind & scalar \\
\hline
Type & Integer \\
\hline
Default & Required / No default \\
\hline
Deprecated & No \\
\hline
\end{longtable}
 feff.tdlda{ ifxc }
 Use TDLDA to calculate x-ray absorption spectrum.
 ifxc - set algorithm for tdlda approximation
        0: use static approximation for screening of the x-ray field
        1: use approximate dynamic screening of x-ray field and core-hole.

\paragraph{Schema}
\begin{Verbatim}
|X| I
\end{Verbatim}

\subsection{feff.temperature}
\label{kw:feff-temperature}
\begin{longtable}{|p{1.5in}|p{4.5in}|}
\hline
Keyword & feff.temperature \\
\hline
Kind & vector \\
\hline
Type & Vector: Float, Integer \\
\hline
Default & Required / No default \\
\hline
Deprecated & No \\
\hline
\end{longtable}
 feff.temperature{ etemp iscfxc }
 The TEMP card sets the electronic temperature and the exchange-correlation potential. etemp
 = 0 (default). etemp is in eV. There are 4 different options for the exchange-correlation
 potential.
 iscfxc
    No temperature dependent exchange correlations
       11 von-Barth Hedin 1971 (default)
       12 Perdew-Zunger
    Explicitly temperature dependent exchange correlations
       21 Perrot Dharma-Wardana 1984
       22 KSDT (recommended)

\paragraph{Schema}
\begin{Verbatim}
|X| F I
\end{Verbatim}

\subsection{feff.title}
\label{kw:feff-title}
\begin{longtable}{|p{1.5in}|p{4.5in}|}
\hline
Keyword & feff.title \\
\hline
Kind & repeating-record \\
\hline
Type & Repeating record with columns: String \\
\hline
Default & Required / No default \\
\hline
Deprecated & No \\
\hline
\end{longtable}
 Set title for this feff calculation.

\paragraph{Schema}
\begin{Verbatim}
|X| S ...
|X| ...
\end{Verbatim}

\subsection{feff.tolscf}
\label{kw:feff-tolscf}
\begin{longtable}{|p{1.5in}|p{4.5in}|}
\hline
Keyword & feff.tolscf \\
\hline
Kind & vector \\
\hline
Type & Vector: Float, Float, Float \\
\hline
Default & Required / No default \\
\hline
Deprecated & No \\
\hline
\end{longtable}
 set tolerances for scf convergence
 feff.tolscf{ fermi\_tol charge\_tol partial\_charge\_tol }

\paragraph{Schema}
\begin{Verbatim}
|X| F F F
\end{Verbatim}

\subsection{feff.unfreezef}
\label{kw:feff-unfreezef}
\begin{longtable}{|p{1.5in}|p{4.5in}|}
\hline
Keyword & feff.unfreezef \\
\hline
Kind & scalar \\
\hline
Type & Logical \\
\hline
Default & Required / No default \\
\hline
Deprecated & No \\
\hline
\end{longtable}
 feff.unreezef{ UnfreezeFOrbitals }
 If true, allow f-orbitals to relax during SCF calculation.

\paragraph{Schema}
\begin{Verbatim}
|X| L
\end{Verbatim}

\subsection{feff.xanes}
\label{kw:feff-xanes}
\begin{longtable}{|p{1.5in}|p{4.5in}|}
\hline
Keyword & feff.xanes \\
\hline
Kind & scalar \\
\hline
Type & Float \\
\hline
Default & Required / No default \\
\hline
Deprecated & No \\
\hline
\end{longtable}
 feff.xanes{ xkmax }
 Calculate XANES spectrum.
   xkmax - calculate up to k = xkmax

\paragraph{Schema}
\begin{Verbatim}
|X| F
\end{Verbatim}

\subsection{feff.xes}
\label{kw:feff-xes}
\begin{longtable}{|p{1.5in}|p{4.5in}|}
\hline
Keyword & feff.xes \\
\hline
Kind & vector \\
\hline
Type & Vector: Float, Float, Float \\
\hline
Default & Required / No default \\
\hline
Deprecated & No \\
\hline
\end{longtable}
 feff.xes{ emin emax estep }
 Calculate x-ray emission spectrum.
   emin  - minimum energy of calculation.
   emax  - maximum energy of calculation
   estep - energy step of calculation

\paragraph{Schema}
\begin{Verbatim}
|X| F F F
\end{Verbatim}

\subsection{feff.xmcd}
\label{kw:feff-xmcd}
\begin{longtable}{|p{1.5in}|p{4.5in}|}
\hline
Keyword & feff.xmcd \\
\hline
Kind & vector \\
\hline
Type & Vector: Float, Float, Float \\
\hline
Default & Required / No default \\
\hline
Deprecated & No \\
\hline
\end{longtable}
 feff.xmcd{ xkmax xkstep estep }
 Caluclate x-ray magnetic (or natural) circular dichroism
   xkmax  - calculate up to k = xkmax
   xkstep - use steps of size xkstep
   estep  - near the edge, use steps calculated from estep

\paragraph{Schema}
\begin{Verbatim}
|X| F F F
\end{Verbatim}

\subsection{nuctemp}
\label{kw:nuctemp}
\begin{longtable}{|p{1.5in}|p{4.5in}|}
\hline
Keyword & nuctemp \\
\hline
Kind & scalar \\
\hline
Type & Float \\
\hline
Default & 300.0 \\
\hline
Deprecated & No \\
\hline
\end{longtable}

 nuctemp
 {
 temp

 Specify the temperature (in K) for the nuclear motion (DW factors).

\paragraph{Schema}
\begin{Verbatim}
300.0 |X| F
\end{Verbatim}

\subsection{opcons.drude}
\label{kw:opcons-drude}
\begin{longtable}{|p{1.5in}|p{4.5in}|}
\hline
Keyword & opcons.drude \\
\hline
Kind & vector \\
\hline
Type & Vector: Float, Float \\
\hline
Default & Required / No default \\
\hline
Deprecated & No \\
\hline
\end{longtable}
 Calculate a drude approximation for the low energy part of the spectrum in opcons.

\paragraph{Schema}
\begin{Verbatim}
|X| F F
\end{Verbatim}

\subsection{opcons.usesaved}
\label{kw:opcons-usesaved}
\begin{longtable}{|p{1.5in}|p{4.5in}|}
\hline
Keyword & opcons.usesaved \\
\hline
Kind & scalar \\
\hline
Type & Logical \\
\hline
Default & Required / No default \\
\hline
Deprecated & No \\
\hline
\end{longtable}
 feff.opcons{ }
 Calculate loss function using an atomic approximation.

\paragraph{Schema}
\begin{Verbatim}
|X| L
\end{Verbatim}

\subsection{run\_dmdw}
\label{kw:run-dmdw}
\begin{longtable}{|p{1.5in}|p{4.5in}|}
\hline
Keyword & run\_dmdw \\
\hline
Kind & scalar \\
\hline
Type & Logical \\
\hline
Default & False \\
\hline
Deprecated & No \\
\hline
\end{longtable}
 run standalone dmdw along with feff runs.

\paragraph{Schema}
\begin{Verbatim}
False |X| L
\end{Verbatim}

\section{ORCA}

\subsection{orca.basis.basis}
\label{kw:orca-basis-basis}
\begin{longtable}{|p{1.5in}|p{4.5in}|}
\hline
Keyword & orca.basis.basis \\
\hline
Kind & scalar \\
\hline
Type & String \\
\hline
Default & def2-TZVP \\
\hline
Deprecated & No \\
\hline
\end{longtable}
\paragraph{Schema}
\begin{Verbatim}
def2-TZVP |X| S
\end{Verbatim}

\subsection{orca.basis.decontract}
\label{kw:orca-basis-decontract}
\begin{longtable}{|p{1.5in}|p{4.5in}|}
\hline
Keyword & orca.basis.decontract \\
\hline
Kind & scalar \\
\hline
Type & Logical \\
\hline
Default & Required / No default \\
\hline
Deprecated & No \\
\hline
\end{longtable}
\paragraph{Schema}
\begin{Verbatim}
|X| L
\end{Verbatim}

\subsection{orca.ci}
\label{kw:orca-ci}
\begin{longtable}{|p{1.5in}|p{4.5in}|}
\hline
Keyword & orca.ci \\
\hline
Kind & scalar \\
\hline
Type & Logical \\
\hline
Default & F \\
\hline
Deprecated & No \\
\hline
\end{longtable}
\paragraph{Schema}
\begin{Verbatim}
F |X| L
\end{Verbatim}

\subsection{orca.citype}
\label{kw:orca-citype}
\begin{longtable}{|p{1.5in}|p{4.5in}|}
\hline
Keyword & orca.citype \\
\hline
Kind & scalar \\
\hline
Type & String \\
\hline
Default & CCSD \\
\hline
Deprecated & No \\
\hline
\end{longtable}
\paragraph{Schema}
\begin{Verbatim}
CCSD |X| S
\end{Verbatim}

\subsection{orca.coords.units}
\label{kw:orca-coords-units}
\begin{longtable}{|p{1.5in}|p{4.5in}|}
\hline
Keyword & orca.coords.units \\
\hline
Kind & scalar \\
\hline
Type & String \\
\hline
Default & Required / No default \\
\hline
Deprecated & No \\
\hline
\end{longtable}
\paragraph{Schema}
\begin{Verbatim}
|X| S
\end{Verbatim}

\subsection{orca.cpcm}
\label{kw:orca-cpcm}
\begin{longtable}{|p{1.5in}|p{4.5in}|}
\hline
Keyword & orca.cpcm \\
\hline
Kind & scalar \\
\hline
Type & String \\
\hline
Default & Required / No default \\
\hline
Deprecated & No \\
\hline
\end{longtable}
\paragraph{Schema}
\begin{Verbatim}
|X| S
\end{Verbatim}

\subsection{orca.expert}
\label{kw:orca-expert}
\begin{longtable}{|p{1.5in}|p{4.5in}|}
\hline
Keyword & orca.expert \\
\hline
Kind & repeating-record \\
\hline
Type & Repeating record with columns: String \\
\hline
Default & Required / No default \\
\hline
Deprecated & No \\
\hline
\end{longtable}
 Use verbatim expert setting for orca.

\paragraph{Schema}
\begin{Verbatim}
|X| S ...
|X| ...
\end{Verbatim}

\subsection{orca.geom.optimizationquality}
\label{kw:orca-geom-optimizationquality}
\begin{longtable}{|p{1.5in}|p{4.5in}|}
\hline
Keyword & orca.geom.optimizationquality \\
\hline
Kind & scalar \\
\hline
Type & String \\
\hline
Default & TIGHTOPT \\
\hline
Deprecated & No \\
\hline
\end{longtable}
\paragraph{Schema}
\begin{Verbatim}
TIGHTOPT |X| S
\end{Verbatim}

\subsection{orca.method.allowrhf}
\label{kw:orca-method-allowrhf}
\begin{longtable}{|p{1.5in}|p{4.5in}|}
\hline
Keyword & orca.method.allowrhf \\
\hline
Kind & scalar \\
\hline
Type & Logical \\
\hline
Default & F \\
\hline
Deprecated & No \\
\hline
\end{longtable}
\paragraph{Schema}
\begin{Verbatim}
F |X| L
\end{Verbatim}

\subsection{orca.method.amass}
\label{kw:orca-method-amass}
\begin{longtable}{|p{1.5in}|p{4.5in}|}
\hline
Keyword & orca.method.amass \\
\hline
Kind & scalar \\
\hline
Type & String \\
\hline
Default & Mass2016 \\
\hline
Deprecated & No \\
\hline
\end{longtable}
\paragraph{Schema}
\begin{Verbatim}
Mass2016 |X| S
\end{Verbatim}

\subsection{orca.method.frozencore}
\label{kw:orca-method-frozencore}
\begin{longtable}{|p{1.5in}|p{4.5in}|}
\hline
Keyword & orca.method.frozencore \\
\hline
Kind & scalar \\
\hline
Type & Logical \\
\hline
Default & F \\
\hline
Deprecated & No \\
\hline
\end{longtable}
\paragraph{Schema}
\begin{Verbatim}
F |X| L
\end{Verbatim}

\subsection{orca.method.grid}
\label{kw:orca-method-grid}
\begin{longtable}{|p{1.5in}|p{4.5in}|}
\hline
Keyword & orca.method.grid \\
\hline
Kind & scalar \\
\hline
Type & String \\
\hline
Default & GRID4 \\
\hline
Deprecated & No \\
\hline
\end{longtable}
\paragraph{Schema}
\begin{Verbatim}
GRID4 |X| S
\end{Verbatim}

\subsection{orca.method.gridx}
\label{kw:orca-method-gridx}
\begin{longtable}{|p{1.5in}|p{4.5in}|}
\hline
Keyword & orca.method.gridx \\
\hline
Kind & scalar \\
\hline
Type & String \\
\hline
Default & Required / No default \\
\hline
Deprecated & No \\
\hline
\end{longtable}
\paragraph{Schema}
\begin{Verbatim}
|X| S
\end{Verbatim}

\subsection{orca.method.method}
\label{kw:orca-method-method}
\begin{longtable}{|p{1.5in}|p{4.5in}|}
\hline
Keyword & orca.method.method \\
\hline
Kind & scalar \\
\hline
Type & String \\
\hline
Default & B3LYP \\
\hline
Deprecated & No \\
\hline
\end{longtable}
\paragraph{Schema}
\begin{Verbatim}
B3LYP |X| S
\end{Verbatim}

\subsection{orca.method.ri}
\label{kw:orca-method-ri}
\begin{longtable}{|p{1.5in}|p{4.5in}|}
\hline
Keyword & orca.method.ri \\
\hline
Kind & scalar \\
\hline
Type & Logical \\
\hline
Default & F \\
\hline
Deprecated & No \\
\hline
\end{longtable}
\paragraph{Schema}
\begin{Verbatim}
F |X| L
\end{Verbatim}

\subsection{orca.method.runtype}
\label{kw:orca-method-runtype}
\begin{longtable}{|p{1.5in}|p{4.5in}|}
\hline
Keyword & orca.method.runtype \\
\hline
Kind & scalar \\
\hline
Type & String \\
\hline
Default & OPT \\
\hline
Deprecated & No \\
\hline
\end{longtable}
\paragraph{Schema}
\begin{Verbatim}
OPT |X| S
\end{Verbatim}

\subsection{orca.method.usesymm}
\label{kw:orca-method-usesymm}
\begin{longtable}{|p{1.5in}|p{4.5in}|}
\hline
Keyword & orca.method.usesymm \\
\hline
Kind & scalar \\
\hline
Type & Logical \\
\hline
Default & F \\
\hline
Deprecated & No \\
\hline
\end{longtable}
\paragraph{Schema}
\begin{Verbatim}
F |X| L
\end{Verbatim}

\subsection{orca.mp2}
\label{kw:orca-mp2}
\begin{longtable}{|p{1.5in}|p{4.5in}|}
\hline
Keyword & orca.mp2 \\
\hline
Kind & scalar \\
\hline
Type & Logical \\
\hline
Default & F \\
\hline
Deprecated & No \\
\hline
\end{longtable}
\paragraph{Schema}
\begin{Verbatim}
F |X| L
\end{Verbatim}

\subsection{orca.mp2type}
\label{kw:orca-mp2type}
\begin{longtable}{|p{1.5in}|p{4.5in}|}
\hline
Keyword & orca.mp2type \\
\hline
Kind & scalar \\
\hline
Type & String \\
\hline
Default & MP2 \\
\hline
Deprecated & No \\
\hline
\end{longtable}
\paragraph{Schema}
\begin{Verbatim}
MP2 |X| S
\end{Verbatim}

\subsection{orca.output.pdbfile}
\label{kw:orca-output-pdbfile}
\begin{longtable}{|p{1.5in}|p{4.5in}|}
\hline
Keyword & orca.output.pdbfile \\
\hline
Kind & scalar \\
\hline
Type & Logical \\
\hline
Default & Required / No default \\
\hline
Deprecated & No \\
\hline
\end{longtable}
\paragraph{Schema}
\begin{Verbatim}
|X| L
\end{Verbatim}

\subsection{orca.output.print}
\label{kw:orca-output-print}
\begin{longtable}{|p{1.5in}|p{4.5in}|}
\hline
Keyword & orca.output.print \\
\hline
Kind & scalar \\
\hline
Type & String \\
\hline
Default & Required / No default \\
\hline
Deprecated & No \\
\hline
\end{longtable}
\paragraph{Schema}
\begin{Verbatim}
|X| S
\end{Verbatim}

\subsection{orca.output.printlevel}
\label{kw:orca-output-printlevel}
\begin{longtable}{|p{1.5in}|p{4.5in}|}
\hline
Keyword & orca.output.printlevel \\
\hline
Kind & scalar \\
\hline
Type & String \\
\hline
Default & Required / No default \\
\hline
Deprecated & No \\
\hline
\end{longtable}
\paragraph{Schema}
\begin{Verbatim}
|X| S
\end{Verbatim}

\subsection{orca.output.xyzfile}
\label{kw:orca-output-xyzfile}
\begin{longtable}{|p{1.5in}|p{4.5in}|}
\hline
Keyword & orca.output.xyzfile \\
\hline
Kind & scalar \\
\hline
Type & Logical \\
\hline
Default & Required / No default \\
\hline
Deprecated & No \\
\hline
\end{longtable}
\paragraph{Schema}
\begin{Verbatim}
|X| L
\end{Verbatim}

\subsection{orca.rel.method}
\label{kw:orca-rel-method}
\begin{longtable}{|p{1.5in}|p{4.5in}|}
\hline
Keyword & orca.rel.method \\
\hline
Kind & scalar \\
\hline
Type & String \\
\hline
Default & Required / No default \\
\hline
Deprecated & No \\
\hline
\end{longtable}
\paragraph{Schema}
\begin{Verbatim}
|X| S
\end{Verbatim}

\subsection{orca.rel.soctype}
\label{kw:orca-rel-soctype}
\begin{longtable}{|p{1.5in}|p{4.5in}|}
\hline
Keyword & orca.rel.soctype \\
\hline
Kind & scalar \\
\hline
Type & String \\
\hline
Default & Required / No default \\
\hline
Deprecated & No \\
\hline
\end{longtable}
\paragraph{Schema}
\begin{Verbatim}
|X| S
\end{Verbatim}

\subsection{orca.scf.autostart}
\label{kw:orca-scf-autostart}
\begin{longtable}{|p{1.5in}|p{4.5in}|}
\hline
Keyword & orca.scf.autostart \\
\hline
Kind & scalar \\
\hline
Type & Logical \\
\hline
Default & T \\
\hline
Deprecated & No \\
\hline
\end{longtable}
\paragraph{Schema}
\begin{Verbatim}
T |X| L
\end{Verbatim}

\subsection{orca.scf.cnvdamp}
\label{kw:orca-scf-cnvdamp}
\begin{longtable}{|p{1.5in}|p{4.5in}|}
\hline
Keyword & orca.scf.cnvdamp \\
\hline
Kind & scalar \\
\hline
Type & Logical \\
\hline
Default & Required / No default \\
\hline
Deprecated & No \\
\hline
\end{longtable}
\paragraph{Schema}
\begin{Verbatim}
|X| L
\end{Verbatim}

\subsection{orca.scf.cnvshift}
\label{kw:orca-scf-cnvshift}
\begin{longtable}{|p{1.5in}|p{4.5in}|}
\hline
Keyword & orca.scf.cnvshift \\
\hline
Kind & scalar \\
\hline
Type & Logical \\
\hline
Default & Required / No default \\
\hline
Deprecated & No \\
\hline
\end{longtable}
\paragraph{Schema}
\begin{Verbatim}
|X| L
\end{Verbatim}

\subsection{orca.scf.convergence}
\label{kw:orca-scf-convergence}
\begin{longtable}{|p{1.5in}|p{4.5in}|}
\hline
Keyword & orca.scf.convergence \\
\hline
Kind & scalar \\
\hline
Type & String \\
\hline
Default & NORMALSCF \\
\hline
Deprecated & No \\
\hline
\end{longtable}
\paragraph{Schema}
\begin{Verbatim}
NORMALSCF |X| S
\end{Verbatim}

\subsection{orca.scf.convergencestrategy}
\label{kw:orca-scf-convergencestrategy}
\begin{longtable}{|p{1.5in}|p{4.5in}|}
\hline
Keyword & orca.scf.convergencestrategy \\
\hline
Kind & scalar \\
\hline
Type & String \\
\hline
Default & NormalConv \\
\hline
Deprecated & No \\
\hline
\end{longtable}
\paragraph{Schema}
\begin{Verbatim}
NormalConv |X| S
\end{Verbatim}

\subsection{orca.scf.diis}
\label{kw:orca-scf-diis}
\begin{longtable}{|p{1.5in}|p{4.5in}|}
\hline
Keyword & orca.scf.diis \\
\hline
Kind & scalar \\
\hline
Type & Logical \\
\hline
Default & Required / No default \\
\hline
Deprecated & No \\
\hline
\end{longtable}
\paragraph{Schema}
\begin{Verbatim}
|X| L
\end{Verbatim}

\subsection{orca.scf.fracocc}
\label{kw:orca-scf-fracocc}
\begin{longtable}{|p{1.5in}|p{4.5in}|}
\hline
Keyword & orca.scf.fracocc \\
\hline
Kind & scalar \\
\hline
Type & String \\
\hline
Default & Required / No default \\
\hline
Deprecated & No \\
\hline
\end{longtable}
\paragraph{Schema}
\begin{Verbatim}
|X| S
\end{Verbatim}

\subsection{orca.scf.guess}
\label{kw:orca-scf-guess}
\begin{longtable}{|p{1.5in}|p{4.5in}|}
\hline
Keyword & orca.scf.guess \\
\hline
Kind & scalar \\
\hline
Type & String \\
\hline
Default & Required / No default \\
\hline
Deprecated & No \\
\hline
\end{longtable}
\paragraph{Schema}
\begin{Verbatim}
|X| S
\end{Verbatim}

\subsection{orca.scf.hftyp}
\label{kw:orca-scf-hftyp}
\begin{longtable}{|p{1.5in}|p{4.5in}|}
\hline
Keyword & orca.scf.hftyp \\
\hline
Kind & scalar \\
\hline
Type & String \\
\hline
Default & RHF \\
\hline
Deprecated & No \\
\hline
\end{longtable}
\paragraph{Schema}
\begin{Verbatim}
RHF |X| S
\end{Verbatim}

\subsection{orca.scf.jmatrix}
\label{kw:orca-scf-jmatrix}
\begin{longtable}{|p{1.5in}|p{4.5in}|}
\hline
Keyword & orca.scf.jmatrix \\
\hline
Kind & scalar \\
\hline
Type & String \\
\hline
Default & Required / No default \\
\hline
Deprecated & No \\
\hline
\end{longtable}
\paragraph{Schema}
\begin{Verbatim}
|X| S
\end{Verbatim}

\subsection{orca.scf.kdiis}
\label{kw:orca-scf-kdiis}
\begin{longtable}{|p{1.5in}|p{4.5in}|}
\hline
Keyword & orca.scf.kdiis \\
\hline
Kind & scalar \\
\hline
Type & Logical \\
\hline
Default & Required / No default \\
\hline
Deprecated & No \\
\hline
\end{longtable}
\paragraph{Schema}
\begin{Verbatim}
|X| L
\end{Verbatim}

\subsection{orca.scf.keepdens}
\label{kw:orca-scf-keepdens}
\begin{longtable}{|p{1.5in}|p{4.5in}|}
\hline
Keyword & orca.scf.keepdens \\
\hline
Kind & scalar \\
\hline
Type & Logical \\
\hline
Default & Required / No default \\
\hline
Deprecated & No \\
\hline
\end{longtable}
\paragraph{Schema}
\begin{Verbatim}
|X| L
\end{Verbatim}

\subsection{orca.scf.keepints}
\label{kw:orca-scf-keepints}
\begin{longtable}{|p{1.5in}|p{4.5in}|}
\hline
Keyword & orca.scf.keepints \\
\hline
Kind & scalar \\
\hline
Type & Logical \\
\hline
Default & Required / No default \\
\hline
Deprecated & No \\
\hline
\end{longtable}
\paragraph{Schema}
\begin{Verbatim}
|X| L
\end{Verbatim}

\subsection{orca.scf.kmatrix}
\label{kw:orca-scf-kmatrix}
\begin{longtable}{|p{1.5in}|p{4.5in}|}
\hline
Keyword & orca.scf.kmatrix \\
\hline
Kind & scalar \\
\hline
Type & String \\
\hline
Default & Required / No default \\
\hline
Deprecated & No \\
\hline
\end{longtable}
\paragraph{Schema}
\begin{Verbatim}
|X| S
\end{Verbatim}

\subsection{orca.scf.maxiter}
\label{kw:orca-scf-maxiter}
\begin{longtable}{|p{1.5in}|p{4.5in}|}
\hline
Keyword & orca.scf.maxiter \\
\hline
Kind & scalar \\
\hline
Type & String \\
\hline
Default & Required / No default \\
\hline
Deprecated & No \\
\hline
\end{longtable}
\paragraph{Schema}
\begin{Verbatim}
|X| S
\end{Verbatim}

\subsection{orca.scf.nr}
\label{kw:orca-scf-nr}
\begin{longtable}{|p{1.5in}|p{4.5in}|}
\hline
Keyword & orca.scf.nr \\
\hline
Kind & scalar \\
\hline
Type & Logical \\
\hline
Default & Required / No default \\
\hline
Deprecated & No \\
\hline
\end{longtable}
\paragraph{Schema}
\begin{Verbatim}
|X| L
\end{Verbatim}

\subsection{orca.scf.readints}
\label{kw:orca-scf-readints}
\begin{longtable}{|p{1.5in}|p{4.5in}|}
\hline
Keyword & orca.scf.readints \\
\hline
Kind & scalar \\
\hline
Type & Logical \\
\hline
Default & Required / No default \\
\hline
Deprecated & No \\
\hline
\end{longtable}
\paragraph{Schema}
\begin{Verbatim}
|X| L
\end{Verbatim}

\subsection{orca.scf.scfmode}
\label{kw:orca-scf-scfmode}
\begin{longtable}{|p{1.5in}|p{4.5in}|}
\hline
Keyword & orca.scf.scfmode \\
\hline
Kind & scalar \\
\hline
Type & String \\
\hline
Default & Required / No default \\
\hline
Deprecated & No \\
\hline
\end{longtable}
\paragraph{Schema}
\begin{Verbatim}
|X| S
\end{Verbatim}

\subsection{orca.scf.smeartemp}
\label{kw:orca-scf-smeartemp}
\begin{longtable}{|p{1.5in}|p{4.5in}|}
\hline
Keyword & orca.scf.smeartemp \\
\hline
Kind & scalar \\
\hline
Type & Logical \\
\hline
Default & Required / No default \\
\hline
Deprecated & No \\
\hline
\end{longtable}
\paragraph{Schema}
\begin{Verbatim}
|X| L
\end{Verbatim}

\subsection{orca.scf.soscf}
\label{kw:orca-scf-soscf}
\begin{longtable}{|p{1.5in}|p{4.5in}|}
\hline
Keyword & orca.scf.soscf \\
\hline
Kind & scalar \\
\hline
Type & Logical \\
\hline
Default & Required / No default \\
\hline
Deprecated & No \\
\hline
\end{longtable}
\paragraph{Schema}
\begin{Verbatim}
|X| L
\end{Verbatim}

\subsection{orca.scf.uno}
\label{kw:orca-scf-uno}
\begin{longtable}{|p{1.5in}|p{4.5in}|}
\hline
Keyword & orca.scf.uno \\
\hline
Kind & scalar \\
\hline
Type & Logical \\
\hline
Default & Required / No default \\
\hline
Deprecated & No \\
\hline
\end{longtable}
\paragraph{Schema}
\begin{Verbatim}
|X| L
\end{Verbatim}

\subsection{orca.scf.usecheapints}
\label{kw:orca-scf-usecheapints}
\begin{longtable}{|p{1.5in}|p{4.5in}|}
\hline
Keyword & orca.scf.usecheapints \\
\hline
Kind & scalar \\
\hline
Type & Logical \\
\hline
Default & Required / No default \\
\hline
Deprecated & No \\
\hline
\end{longtable}
\paragraph{Schema}
\begin{Verbatim}
|X| L
\end{Verbatim}

\subsection{orca.scf.valformat}
\label{kw:orca-scf-valformat}
\begin{longtable}{|p{1.5in}|p{4.5in}|}
\hline
Keyword & orca.scf.valformat \\
\hline
Kind & scalar \\
\hline
Type & String \\
\hline
Default & Required / No default \\
\hline
Deprecated & No \\
\hline
\end{longtable}
\paragraph{Schema}
\begin{Verbatim}
|X| S
\end{Verbatim}

\subsection{orca.tole}
\label{kw:orca-tole}
\begin{longtable}{|p{1.5in}|p{4.5in}|}
\hline
Keyword & orca.tole \\
\hline
Kind & scalar \\
\hline
Type & Float \\
\hline
Default & Required / No default \\
\hline
Deprecated & No \\
\hline
\end{longtable}
\paragraph{Schema}
\begin{Verbatim}
|X| F
\end{Verbatim}

\subsection{orca.tolmaxg}
\label{kw:orca-tolmaxg}
\begin{longtable}{|p{1.5in}|p{4.5in}|}
\hline
Keyword & orca.tolmaxg \\
\hline
Kind & scalar \\
\hline
Type & Float \\
\hline
Default & Required / No default \\
\hline
Deprecated & No \\
\hline
\end{longtable}
\paragraph{Schema}
\begin{Verbatim}
|X| F
\end{Verbatim}

\subsection{orca.tolrmaxd}
\label{kw:orca-tolrmaxd}
\begin{longtable}{|p{1.5in}|p{4.5in}|}
\hline
Keyword & orca.tolrmaxd \\
\hline
Kind & scalar \\
\hline
Type & Float \\
\hline
Default & Required / No default \\
\hline
Deprecated & No \\
\hline
\end{longtable}
\paragraph{Schema}
\begin{Verbatim}
|X| F
\end{Verbatim}

\subsection{orca.tolrmsd}
\label{kw:orca-tolrmsd}
\begin{longtable}{|p{1.5in}|p{4.5in}|}
\hline
Keyword & orca.tolrmsd \\
\hline
Kind & scalar \\
\hline
Type & Float \\
\hline
Default & Required / No default \\
\hline
Deprecated & No \\
\hline
\end{longtable}
\paragraph{Schema}
\begin{Verbatim}
|X| F
\end{Verbatim}

\subsection{orca.tolrmsg}
\label{kw:orca-tolrmsg}
\begin{longtable}{|p{1.5in}|p{4.5in}|}
\hline
Keyword & orca.tolrmsg \\
\hline
Kind & scalar \\
\hline
Type & Float \\
\hline
Default & Required / No default \\
\hline
Deprecated & No \\
\hline
\end{longtable}
\paragraph{Schema}
\begin{Verbatim}
|X| F
\end{Verbatim}

\section{SIESTA}

\subsection{phsf.broadening}
\label{kw:phsf-broadening}
\begin{longtable}{|p{1.5in}|p{4.5in}|}
\hline
Keyword & phsf.broadening \\
\hline
Kind & scalar \\
\hline
Type & Float \\
\hline
Default & 0.5 \\
\hline
Deprecated & No \\
\hline
\end{longtable}
\paragraph{Schema}
\begin{Verbatim}
0.5 |X| F
\end{Verbatim}

\subsection{phsf.ekeqp}
\label{kw:phsf-ekeqp}
\begin{longtable}{|p{1.5in}|p{4.5in}|}
\hline
Keyword & phsf.ekeqp \\
\hline
Kind & vector \\
\hline
Type & Vector: Float, Float \\
\hline
Default & 0.0 0.0 \\
\hline
Deprecated & No \\
\hline
\end{longtable}
\paragraph{Schema}
\begin{Verbatim}
0.0 0.0 |X| F F
\end{Verbatim}

\subsection{phsf.numtimepoints}
\label{kw:phsf-numtimepoints}
\begin{longtable}{|p{1.5in}|p{4.5in}|}
\hline
Keyword & phsf.numtimepoints \\
\hline
Kind & scalar \\
\hline
Type & Integer \\
\hline
Default & 40000 \\
\hline
Deprecated & No \\
\hline
\end{longtable}
\paragraph{Schema}
\begin{Verbatim}
40000 |X| I
\end{Verbatim}

\subsection{siesta.AtomicCoordinatesFormat}
\label{kw:siesta-AtomicCoordinatesFormat}
\begin{longtable}{|p{1.5in}|p{4.5in}|}
\hline
Keyword & siesta.AtomicCoordinatesFormat \\
\hline
Kind & scalar \\
\hline
Type & String \\
\hline
Default & Fractional \\
\hline
Deprecated & No \\
\hline
\end{longtable}
 Specify units and format of atomic coordinates.
 Borh            - cartesian in Bohr.
 Ang             - cartesian in Angstrom.
 ScaledCartesian - cartesian scaled by lattice constant.
 Fractional      - coordinates are given in units of the lattice vectors.

\paragraph{Schema}
\begin{Verbatim}
Fractional |X| S
\end{Verbatim}

\subsection{siesta.Beta.Broadening}
\label{kw:siesta-Beta-Broadening}
\begin{longtable}{|p{1.5in}|p{4.5in}|}
\hline
Keyword & siesta.Beta.Broadening \\
\hline
Kind & scalar \\
\hline
Type & Float \\
\hline
Default & 0.1 \\
\hline
Deprecated & No \\
\hline
\end{longtable}
\paragraph{Schema}
\begin{Verbatim}
0.1 |X| F
\end{Verbatim}

\subsection{siesta.Block.AtomicCoordinatesAndAtomicSpecies}
\label{kw:siesta-Block-AtomicCoordinatesAndAtomicSpecies}
\begin{longtable}{|p{1.5in}|p{4.5in}|}
\hline
Keyword & siesta.Block.AtomicCoordinatesAndAtomicSpecies \\
\hline
Kind & repeating-record \\
\hline
Type & Repeating record with columns: Float, Float, Float, Integer \\
\hline
Default & Required / No default \\
\hline
Deprecated & No \\
\hline
\end{longtable}
 Specify coordinates of each atom as well as the species index.
 x1 y1 z1 SpeciesIndex1
 x2 y2 z2 SpeciesIndex2
 ...

\paragraph{Schema}
\begin{Verbatim}
|X| F F F I
|X| ...
\end{Verbatim}

\subsection{siesta.Block.ChemicalSpeciesLabel}
\label{kw:siesta-Block-ChemicalSpeciesLabel}
\begin{longtable}{|p{1.5in}|p{4.5in}|}
\hline
Keyword & siesta.Block.ChemicalSpeciesLabel \\
\hline
Kind & vector \\
\hline
Type & Vector: Integer, Integer, String \\
\hline
Default & Required / No default \\
\hline
Deprecated & No \\
\hline
\end{longtable}
 Indicate the different chemical species in the unit cell.
 Format:
 SpeciesIndex1 AtomicNumber1 AtomicSymbol1
 SpeciesIndex2 AtomicNumber2 AtomicSymbol2
 ...

\paragraph{Schema}
\begin{Verbatim}
|X| I I S
\end{Verbatim}

\subsection{siesta.Block.LatticeParameters}
\label{kw:siesta-Block-LatticeParameters}
\begin{longtable}{|p{1.5in}|p{4.5in}|}
\hline
Keyword & siesta.Block.LatticeParameters \\
\hline
Kind & vector \\
\hline
Type & Vector: Float, Float, Float, Float, Float, Float \\
\hline
Default & Required / No default \\
\hline
Deprecated & No \\
\hline
\end{longtable}
 Set lattice parameters, i.e.,
 a b c alpha beta gamma

\paragraph{Schema}
\begin{Verbatim}
|X| F F F F F F
\end{Verbatim}

\subsection{siesta.Block.PAO.Basis}
\label{kw:siesta-Block-PAO-Basis}
\begin{longtable}{|p{1.5in}|p{4.5in}|}
\hline
Keyword & siesta.Block.PAO.Basis \\
\hline
Kind & repeating-record \\
\hline
Type & Repeating record with columns: String \\
\hline
Default & Required / No default \\
\hline
Deprecated & No \\
\hline
\end{longtable}
 Set siesta basis manually

\paragraph{Schema}
\begin{Verbatim}
|X| S ...
|X| ...
\end{Verbatim}

\subsection{siesta.Coreresponse.Broadening}
\label{kw:siesta-Coreresponse-Broadening}
\begin{longtable}{|p{1.5in}|p{4.5in}|}
\hline
Keyword & siesta.Coreresponse.Broadening \\
\hline
Kind & scalar \\
\hline
Type & Float \\
\hline
Default & 0.5 \\
\hline
Deprecated & No \\
\hline
\end{longtable}
\paragraph{Schema}
\begin{Verbatim}
0.5 |X| F
\end{Verbatim}

\subsection{siesta.Diag.DivideAndConquer}
\label{kw:siesta-Diag-DivideAndConquer}
\begin{longtable}{|p{1.5in}|p{4.5in}|}
\hline
Keyword & siesta.Diag.DivideAndConquer \\
\hline
Kind & scalar \\
\hline
Type & String \\
\hline
Default & Required / No default \\
\hline
Deprecated & No \\
\hline
\end{longtable}
\paragraph{Schema}
\begin{Verbatim}
|X| S
\end{Verbatim}

\subsection{siesta.DM.InitSpin.AF}
\label{kw:siesta-DM-InitSpin-AF}
\begin{longtable}{|p{1.5in}|p{4.5in}|}
\hline
Keyword & siesta.DM.InitSpin.AF \\
\hline
Kind & scalar \\
\hline
Type & Logical \\
\hline
Default & Required / No default \\
\hline
Deprecated & No \\
\hline
\end{longtable}
 Set magnetic behavior of system to antiferromagnetic if true.
 Default is false and will be set up as ferromagnetic.

\paragraph{Schema}
\begin{Verbatim}
|X| L
\end{Verbatim}

\subsection{siesta.DM.MixingWeight}
\label{kw:siesta-DM-MixingWeight}
\begin{longtable}{|p{1.5in}|p{4.5in}|}
\hline
Keyword & siesta.DM.MixingWeight \\
\hline
Kind & scalar \\
\hline
Type & Float \\
\hline
Default & 0.10 \\
\hline
Deprecated & No \\
\hline
\end{longtable}
 Mixing weight for the density matrix.

\paragraph{Schema}
\begin{Verbatim}
0.10 |X| F
\end{Verbatim}

\subsection{siesta.DM.NumberBroyden}
\label{kw:siesta-DM-NumberBroyden}
\begin{longtable}{|p{1.5in}|p{4.5in}|}
\hline
Keyword & siesta.DM.NumberBroyden \\
\hline
Kind & scalar \\
\hline
Type & Integer \\
\hline
Default & Required / No default \\
\hline
Deprecated & No \\
\hline
\end{longtable}
 Set number of broyden steps

\paragraph{Schema}
\begin{Verbatim}
|X| I
\end{Verbatim}

\subsection{siesta.DM.NumberBroyden}
\label{kw:siesta-DM-NumberBroyden}
\begin{longtable}{|p{1.5in}|p{4.5in}|}
\hline
Keyword & siesta.DM.NumberBroyden \\
\hline
Kind & scalar \\
\hline
Type & Integer \\
\hline
Default & Required / No default \\
\hline
Deprecated & No \\
\hline
\end{longtable}
\paragraph{Schema}
\begin{Verbatim}
|X| I
\end{Verbatim}

\subsection{siesta.DM.NumberPulay}
\label{kw:siesta-DM-NumberPulay}
\begin{longtable}{|p{1.5in}|p{4.5in}|}
\hline
Keyword & siesta.DM.NumberPulay \\
\hline
Kind & scalar \\
\hline
Type & Integer \\
\hline
Default & Required / No default \\
\hline
Deprecated & No \\
\hline
\end{longtable}
 Set number of Pulay steps

\paragraph{Schema}
\begin{Verbatim}
|X| I
\end{Verbatim}

\subsection{siesta.DM.Tolerance}
\label{kw:siesta-DM-Tolerance}
\begin{longtable}{|p{1.5in}|p{4.5in}|}
\hline
Keyword & siesta.DM.Tolerance \\
\hline
Kind & scalar \\
\hline
Type & Float \\
\hline
Default & 1.0e-4 \\
\hline
Deprecated & No \\
\hline
\end{longtable}
 Convergence tolerance for elements of the density matrix

\paragraph{Schema}
\begin{Verbatim}
1.0e-4 |X| F
\end{Verbatim}

\subsection{siesta.ElectronicTemperature}
\label{kw:siesta-ElectronicTemperature}
\begin{longtable}{|p{1.5in}|p{4.5in}|}
\hline
Keyword & siesta.ElectronicTemperature \\
\hline
Kind & vector \\
\hline
Type & Vector: Float, String \\
\hline
Default & 1.0 K \\
\hline
Deprecated & No \\
\hline
\end{longtable}
 Set electronic teperature used to define occupation numbers during SCF.

\paragraph{Schema}
\begin{Verbatim}
1.0 K |X| F S
\end{Verbatim}

\subsection{siesta.expert}
\label{kw:siesta-expert}
\begin{longtable}{|p{1.5in}|p{4.5in}|}
\hline
Keyword & siesta.expert \\
\hline
Kind & repeating-record \\
\hline
Type & Repeating record with columns: String \\
\hline
Default & Required / No default \\
\hline
Deprecated & No \\
\hline
\end{longtable}
\paragraph{Schema}
\begin{Verbatim}
|X| S ...
|X| ...
\end{Verbatim}

\subsection{siesta.LatticeConstant}
\label{kw:siesta-LatticeConstant}
\begin{longtable}{|p{1.5in}|p{4.5in}|}
\hline
Keyword & siesta.LatticeConstant \\
\hline
Kind & vector \\
\hline
Type & Vector: Float, String \\
\hline
Default & Required / No default \\
\hline
Deprecated & No \\
\hline
\end{longtable}
 Set an overall multiplicative constant for the lattice parameters.

\paragraph{Schema}
\begin{Verbatim}
|X| F S
\end{Verbatim}

\subsection{siesta.LongOutput}
\label{kw:siesta-LongOutput}
\begin{longtable}{|p{1.5in}|p{4.5in}|}
\hline
Keyword & siesta.LongOutput \\
\hline
Kind & scalar \\
\hline
Type & Logical \\
\hline
Default & T \\
\hline
Deprecated & No \\
\hline
\end{longtable}
 Output lots of info.

\paragraph{Schema}
\begin{Verbatim}
T |X| L
\end{Verbatim}

\subsection{siesta.MaxSCFIterations}
\label{kw:siesta-MaxSCFIterations}
\begin{longtable}{|p{1.5in}|p{4.5in}|}
\hline
Keyword & siesta.MaxSCFIterations \\
\hline
Kind & scalar \\
\hline
Type & Integer \\
\hline
Default & 100 \\
\hline
Deprecated & No \\
\hline
\end{longtable}
 Max number of SCF iterations allowed.

\paragraph{Schema}
\begin{Verbatim}
100 |X| I
\end{Verbatim}

\subsection{siesta.MeshCutoff}
\label{kw:siesta-MeshCutoff}
\begin{longtable}{|p{1.5in}|p{4.5in}|}
\hline
Keyword & siesta.MeshCutoff \\
\hline
Kind & vector \\
\hline
Type & Vector: Float, String \\
\hline
Default & 120.0 Ry \\
\hline
Deprecated & No \\
\hline
\end{longtable}
 Energy that defines the real-space grid. Point spacing scales as $1/\sqrt(E_c)$.

\paragraph{Schema}
\begin{Verbatim}
120.0 Ry |X| F S
\end{Verbatim}

\subsection{siesta.NumberOfAtoms}
\label{kw:siesta-NumberOfAtoms}
\begin{longtable}{|p{1.5in}|p{4.5in}|}
\hline
Keyword & siesta.NumberOfAtoms \\
\hline
Kind & scalar \\
\hline
Type & Integer \\
\hline
Default & Required / No default \\
\hline
Deprecated & No \\
\hline
\end{longtable}
 Indicate number of atoms in unit cell.

\paragraph{Schema}
\begin{Verbatim}
|X| I
\end{Verbatim}

\subsection{siesta.NumberOfSpecies}
\label{kw:siesta-NumberOfSpecies}
\begin{longtable}{|p{1.5in}|p{4.5in}|}
\hline
Keyword & siesta.NumberOfSpecies \\
\hline
Kind & scalar \\
\hline
Type & Integer \\
\hline
Default & Required / No default \\
\hline
Deprecated & No \\
\hline
\end{longtable}
 Indicate number of species in unit cell.

\paragraph{Schema}
\begin{Verbatim}
|X| I
\end{Verbatim}

\subsection{siesta.PAO.BasisSize}
\label{kw:siesta-PAO-BasisSize}
\begin{longtable}{|p{1.5in}|p{4.5in}|}
\hline
Keyword & siesta.PAO.BasisSize \\
\hline
Kind & scalar \\
\hline
Type & String \\
\hline
Default & DZP \\
\hline
Deprecated & No \\
\hline
\end{longtable}
 Define basis size. Can be SZ, DZ, SZP, DZP.

\paragraph{Schema}
\begin{Verbatim}
DZP |X| S
\end{Verbatim}

\subsection{siesta.PAO.BasisType}
\label{kw:siesta-PAO-BasisType}
\begin{longtable}{|p{1.5in}|p{4.5in}|}
\hline
Keyword & siesta.PAO.BasisType \\
\hline
Kind & scalar \\
\hline
Type & String \\
\hline
Default & split \\
\hline
Deprecated & No \\
\hline
\end{longtable}
 Define basis type
 split: Split-valence scheme for multiple-zeta. The split is based on different radii.
 splitgauss: Same as split but using Gaussian functions.
 nodes: Generalized PAO’s.
 nonodes: The original PAO’s are used, multiple-zeta is generated by changing the
 scale-factors, instead of using the excited orbitals.
 filteret: Use the filterets as a systematic basis set. The size of the basis set is
 controlled by the filter cut-off for the orbitals.

\paragraph{Schema}
\begin{Verbatim}
split |X| S
\end{Verbatim}

\subsection{siesta.SolutionMethod}
\label{kw:siesta-SolutionMethod}
\begin{longtable}{|p{1.5in}|p{4.5in}|}
\hline
Keyword & siesta.SolutionMethod \\
\hline
Kind & scalar \\
\hline
Type & String \\
\hline
Default & diagon \\
\hline
Deprecated & No \\
\hline
\end{longtable}
 Method for solving the eigenvalue problem.

\paragraph{Schema}
\begin{Verbatim}
diagon |X| S
\end{Verbatim}

\subsection{siesta.SpinPolarized}
\label{kw:siesta-SpinPolarized}
\begin{longtable}{|p{1.5in}|p{4.5in}|}
\hline
Keyword & siesta.SpinPolarized \\
\hline
Kind & scalar \\
\hline
Type & String \\
\hline
Default & Required / No default \\
\hline
Deprecated & No \\
\hline
\end{longtable}
 Include spin in calculation.
 Value should be yes or no.

\paragraph{Schema}
\begin{Verbatim}
|X| S
\end{Verbatim}

\subsection{siesta.TD.CoreExcitedAtom}
\label{kw:siesta-TD-CoreExcitedAtom}
\begin{longtable}{|p{1.5in}|p{4.5in}|}
\hline
Keyword & siesta.TD.CoreExcitedAtom \\
\hline
Kind & scalar \\
\hline
Type & Integer \\
\hline
Default & Required / No default \\
\hline
Deprecated & No \\
\hline
\end{longtable}
 Set which atom will be the absorbing atom, i.e., which atoms the
 core-hole will be centered on

\paragraph{Schema}
\begin{Verbatim}
|X| I
\end{Verbatim}

\subsection{siesta.TD.CorePerturbationCharge}
\label{kw:siesta-TD-CorePerturbationCharge}
\begin{longtable}{|p{1.5in}|p{4.5in}|}
\hline
Keyword & siesta.TD.CorePerturbationCharge \\
\hline
Kind & scalar \\
\hline
Type & Float \\
\hline
Default & Required / No default \\
\hline
Deprecated & No \\
\hline
\end{longtable}
 Set the core-hole charge.

\paragraph{Schema}
\begin{Verbatim}
|X| F
\end{Verbatim}

\subsection{siesta.TD.mxPC}
\label{kw:siesta-TD-mxPC}
\begin{longtable}{|p{1.5in}|p{4.5in}|}
\hline
Keyword & siesta.TD.mxPC \\
\hline
Kind & scalar \\
\hline
Type & Integer \\
\hline
Default & 2 \\
\hline
Deprecated & No \\
\hline
\end{longtable}
 Maximum number of points to use for predictor-corrector algorith.

\paragraph{Schema}
\begin{Verbatim}
2 |X| I
\end{Verbatim}

\subsection{siesta.TD.NumberOfTimeSteps}
\label{kw:siesta-TD-NumberOfTimeSteps}
\begin{longtable}{|p{1.5in}|p{4.5in}|}
\hline
Keyword & siesta.TD.NumberOfTimeSteps \\
\hline
Kind & scalar \\
\hline
Type & Integer \\
\hline
Default & Required / No default \\
\hline
Deprecated & No \\
\hline
\end{longtable}
 Specify number of time steps for RT-TDDFT run.

\paragraph{Schema}
\begin{Verbatim}
|X| I
\end{Verbatim}

\subsection{siesta.TD.ShapeOfEfield}
\label{kw:siesta-TD-ShapeOfEfield}
\begin{longtable}{|p{1.5in}|p{4.5in}|}
\hline
Keyword & siesta.TD.ShapeOfEfield \\
\hline
Kind & scalar \\
\hline
Type & String \\
\hline
Default & Required / No default \\
\hline
Deprecated & No \\
\hline
\end{longtable}
 Specify the shape of the field for the RT-TDDFT run.
 Options are:
 step     - Turn off dipole field at time 0.
 halfstep - Dipole field changes from full strength to 1/2
            strength at time 0.
 delta    - Dipole field with time dependence delta(t).
 const    - Turn on dipole field at time 0.
 sine     - Dipole field with damped sine wave time dependence,
            with optional Gaussian envelope to create quasi-harmonic
            wave packet.
 core     - turn on core-hole potential at time 0.
 core2    - GS with full core-hole. Excited states from linear response after.

\paragraph{Schema}
\begin{Verbatim}
|X| S
\end{Verbatim}

\subsection{siesta.TD.TimeStep}
\label{kw:siesta-TD-TimeStep}
\begin{longtable}{|p{1.5in}|p{4.5in}|}
\hline
Keyword & siesta.TD.TimeStep \\
\hline
Kind & scalar \\
\hline
Type & Float \\
\hline
Default & Required / No default \\
\hline
Deprecated & No \\
\hline
\end{longtable}
 Specify the time step for RT-TDDFT run.

\paragraph{Schema}
\begin{Verbatim}
|X| F
\end{Verbatim}

\subsection{siesta.UseSaveData}
\label{kw:siesta-UseSaveData}
\begin{longtable}{|p{1.5in}|p{4.5in}|}
\hline
Keyword & siesta.UseSaveData \\
\hline
Kind & scalar \\
\hline
Type & Logical \\
\hline
Default & F \\
\hline
Deprecated & No \\
\hline
\end{longtable}
 Use saved data to start siesta calculation

\paragraph{Schema}
\begin{Verbatim}
F |X| L
\end{Verbatim}

\subsection{siesta.XC.authors}
\label{kw:siesta-XC-authors}
\begin{longtable}{|p{1.5in}|p{4.5in}|}
\hline
Keyword & siesta.XC.authors \\
\hline
Kind & scalar \\
\hline
Type & String \\
\hline
Default & PBE \\
\hline
Deprecated & No \\
\hline
\end{longtable}
 Set authors of exchange correlation functional.

\paragraph{Schema}
\begin{Verbatim}
PBE |X| S
\end{Verbatim}

\subsection{siesta.XC.functional}
\label{kw:siesta-XC-functional}
\begin{longtable}{|p{1.5in}|p{4.5in}|}
\hline
Keyword & siesta.XC.functional \\
\hline
Kind & scalar \\
\hline
Type & String \\
\hline
Default & GGA \\
\hline
Deprecated & No \\
\hline
\end{longtable}
 Set flavor of exchange-correlation functional, i.e., LDA or GGA.

\paragraph{Schema}
\begin{Verbatim}
GGA |X| S
\end{Verbatim}

\section{cif2cell}

\subsection{cif2cell.atomicunits}
\label{kw:cif2cell-atomicunits}
\begin{longtable}{|p{1.5in}|p{4.5in}|}
\hline
Keyword & cif2cell.atomicunits \\
\hline
Kind & scalar \\
\hline
Type & Logical \\
\hline
Default & Required / No default \\
\hline
Deprecated & No \\
\hline
\end{longtable}
 print coordinates in bohr rather than angstroms

\paragraph{Schema}
\begin{Verbatim}
|X| L
\end{Verbatim}

\subsection{cif2cell.cartesian}
\label{kw:cif2cell-cartesian}
\begin{longtable}{|p{1.5in}|p{4.5in}|}
\hline
Keyword & cif2cell.cartesian \\
\hline
Kind & scalar \\
\hline
Type & Logical \\
\hline
Default & Required / No default \\
\hline
Deprecated & No \\
\hline
\end{longtable}
 make output in cartesian form

\paragraph{Schema}
\begin{Verbatim}
|X| L
\end{Verbatim}

\subsection{cif2cell.cif\_input}
\label{kw:cif2cell-cif-input}
\begin{longtable}{|p{1.5in}|p{4.5in}|}
\hline
Keyword & cif2cell.cif\_input \\
\hline
Kind & scalar \\
\hline
Type & String \\
\hline
Default & Required / No default \\
\hline
Deprecated & No \\
\hline
\end{longtable}
 Crystallographic information file to use.

\paragraph{Schema}
\begin{Verbatim}
|X| S
\end{Verbatim}

\subsection{cif2cell.outfile}
\label{kw:cif2cell-outfile}
\begin{longtable}{|p{1.5in}|p{4.5in}|}
\hline
Keyword & cif2cell.outfile \\
\hline
Kind & scalar \\
\hline
Type & String \\
\hline
Default & Required / No default \\
\hline
Deprecated & No \\
\hline
\end{longtable}
 Write output to outfile

\paragraph{Schema}
\begin{Verbatim}
|X| S
\end{Verbatim}

\subsection{cif2cell.program}
\label{kw:cif2cell-program}
\begin{longtable}{|p{1.5in}|p{4.5in}|}
\hline
Keyword & cif2cell.program \\
\hline
Kind & scalar \\
\hline
Type & String \\
\hline
Default & Required / No default \\
\hline
Deprecated & No \\
\hline
\end{longtable}
\paragraph{Schema}
\begin{Verbatim}
% Program for cif2cell to write input for.
|X| S
\end{Verbatim}

\subsection{cif2cell.supercell}
\label{kw:cif2cell-supercell}
\begin{longtable}{|p{1.5in}|p{4.5in}|}
\hline
Keyword & cif2cell.supercell \\
\hline
Kind & vector \\
\hline
Type & Vector: Integer, Integer, Integer \\
\hline
Default & Required / No default \\
\hline
Deprecated & No \\
\hline
\end{longtable}
\paragraph{Schema}
\begin{Verbatim}
|X| I I I
\end{Verbatim}

\section{Loop Workflows}

\subsection{loop\_labels}
\label{kw:loop-labels}
\begin{longtable}{|p{1.5in}|p{4.5in}|}
\hline
Keyword & loop\_labels \\
\hline
Kind & repeating-record \\
\hline
Type & Repeating record with columns: String \\
\hline
Default & Required / No default \\
\hline
Deprecated & No \\
\hline
\end{longtable}
 define the labels that will be used for the
 loop directories. There should be one label
 for each value of the loop parameter.

\paragraph{Schema}
\begin{Verbatim}
|X| S
|X| ...
\end{Verbatim}

\subsection{loop\_parameter}
\label{kw:loop-parameter}
\begin{longtable}{|p{1.5in}|p{4.5in}|}
\hline
Keyword & loop\_parameter \\
\hline
Kind & repeating-record \\
\hline
Type & Repeating record with columns: String \\
\hline
Default & Required / No default \\
\hline
Deprecated & No \\
\hline
\end{longtable}
 Define loop over a set of parameter values
 for a certain calculation.
 First line is the keyword to loop over.
 The rest of the lines are the values
 associated with each iteration of the loop,
 with one line per set of values.
 Note: This only works with single line
 keywords.

\paragraph{Schema}
\begin{Verbatim}
|X| S
|X| S ...
|X| ...
\end{Verbatim}

\subsection{loop\_target}
\label{kw:loop-target}
\begin{longtable}{|p{1.5in}|p{4.5in}|}
\hline
Keyword & loop\_target \\
\hline
Kind & scalar \\
\hline
Type & String \\
\hline
Default & Required / No default \\
\hline
Deprecated & No \\
\hline
\end{longtable}
 Set the target property to calculate during
 the loop.

\paragraph{Schema}
\begin{Verbatim}
|X| S
\end{Verbatim}

\section{Configurational Averaging}

\subsection{cfavg\_choose\_random\_absorbers}
\label{kw:cfavg-choose-random-absorbers}
\begin{longtable}{|p{1.5in}|p{4.5in}|}
\hline
Keyword & cfavg\_choose\_random\_absorbers \\
\hline
Kind & scalar \\
\hline
Type & Logical \\
\hline
Default & False \\
\hline
Deprecated & No \\
\hline
\end{longtable}
\paragraph{Schema}
\begin{Verbatim}
False |X| L
\end{Verbatim}

\subsection{cfavg\_egrid}
\label{kw:cfavg-egrid}
\begin{longtable}{|p{1.5in}|p{4.5in}|}
\hline
Keyword & cfavg\_egrid \\
\hline
Kind & scalar \\
\hline
Type & String \\
\hline
Default & "first" \\
\hline
Deprecated & No \\
\hline
\end{longtable}
 Choose energy grid type for configurational
 averaging of spectra. "regular" would be
 regular in energy, while "first" uses
 the grid of the first calculation.

\paragraph{Schema}
\begin{Verbatim}
"first" |X| S
\end{Verbatim}

\subsection{cfavg\_max\_configurations}
\label{kw:cfavg-max-configurations}
\begin{longtable}{|p{1.5in}|p{4.5in}|}
\hline
Keyword & cfavg\_max\_configurations \\
\hline
Kind & scalar \\
\hline
Type & Integer \\
\hline
Default & 10 \\
\hline
Deprecated & No \\
\hline
\end{longtable}
 Set the number of configurations used in disordered systems.

\paragraph{Schema}
\begin{Verbatim}
10 |X| I
\end{Verbatim}

\subsection{cfavg\_max\_configurations}
\label{kw:cfavg-max-configurations}
\begin{longtable}{|p{1.5in}|p{4.5in}|}
\hline
Keyword & cfavg\_max\_configurations \\
\hline
Kind & scalar \\
\hline
Type & Integer \\
\hline
Default & Required / No default \\
\hline
Deprecated & No \\
\hline
\end{longtable}
\paragraph{Schema}
\begin{Verbatim}
|X| I
\end{Verbatim}

\subsection{cfavg\_target}
\label{kw:cfavg-target}
\begin{longtable}{|p{1.5in}|p{4.5in}|}
\hline
Keyword & cfavg\_target \\
\hline
Kind & scalar \\
\hline
Type & String \\
\hline
Default & Required / No default \\
\hline
Deprecated & No \\
\hline
\end{longtable}
\paragraph{Schema}
\begin{Verbatim}
|X| S ...
\end{Verbatim}

\section{Fitting}

\subsection{fit.bond}
\label{kw:fit-bond}
\begin{longtable}{|p{1.5in}|p{4.5in}|}
\hline
Keyword & fit.bond \\
\hline
Kind & repeating-record \\
\hline
Type & Repeating record with columns: Integer \\
\hline
Default & Required / No default \\
\hline
Deprecated & No \\
\hline
\end{longtable}
\paragraph{Schema}
\begin{Verbatim}
|X| I ...
|X| ...
\end{Verbatim}

\subsection{fit.datafile}
\label{kw:fit-datafile}
\begin{longtable}{|p{1.5in}|p{4.5in}|}
\hline
Keyword & fit.datafile \\
\hline
Kind & scalar \\
\hline
Type & String \\
\hline
Default & Required / No default \\
\hline
Deprecated & No \\
\hline
\end{longtable}
\paragraph{Schema}
\begin{Verbatim}
|X| S
\end{Verbatim}

\subsection{fit.parameters}
\label{kw:fit-parameters}
\begin{longtable}{|p{1.5in}|p{4.5in}|}
\hline
Keyword & fit.parameters \\
\hline
Kind & repeating-record \\
\hline
Type & Repeating record with columns: String, Float \\
\hline
Default & Required / No default \\
\hline
Deprecated & No \\
\hline
\end{longtable}
\paragraph{Schema}
\begin{Verbatim}
|X| S F
|X| ...
\end{Verbatim}

\subsection{fit.target}
\label{kw:fit-target}
\begin{longtable}{|p{1.5in}|p{4.5in}|}
\hline
Keyword & fit.target \\
\hline
Kind & scalar \\
\hline
Type & String \\
\hline
Default & Required / No default \\
\hline
Deprecated & No \\
\hline
\end{longtable}
\paragraph{Schema}
\begin{Verbatim}
|X| S
\end{Verbatim}

\section{OCEAN}

\subsection{ocean.abpad}
\label{kw:ocean-abpad}
\begin{longtable}{|p{1.5in}|p{4.5in}|}
\hline
Keyword & ocean.abpad \\
\hline
Kind & scalar \\
\hline
Type & Integer \\
\hline
Default & Required / No default \\
\hline
Deprecated & No \\
\hline
\end{longtable}
 For the calculation of the final state wavefunctions, the convergence of the
 highest bands can be very slow. The calculation can be sped up by adding some
 throwaway bands at teh top. These bands are not considered when checking for convergence. ocean.abpad adds bands to the calculation. Only used for abinit.

\paragraph{Schema}
\begin{Verbatim}
|X| I
\end{Verbatim}

\subsection{ocean.acell}
\label{kw:ocean-acell}
\begin{longtable}{|p{1.5in}|p{4.5in}|}
\hline
Keyword & ocean.acell \\
\hline
Kind & vector \\
\hline
Type & Vector: Float, Float, Float \\
\hline
Default & Required / No default \\
\hline
Deprecated & No \\
\hline
\end{longtable}
 Scaling in Bohr for the primitive vectors of the unit cell

\paragraph{Schema}
\begin{Verbatim}
|X| F F F
\end{Verbatim}

\subsection{ocean.cnbse.broaden}
\label{kw:ocean-cnbse-broaden}
\begin{longtable}{|p{1.5in}|p{4.5in}|}
\hline
Keyword & ocean.cnbse.broaden \\
\hline
Kind & scalar \\
\hline
Type & Float \\
\hline
Default & Required / No default \\
\hline
Deprecated & No \\
\hline
\end{longtable}
 Half? width of broadening in eV for the spectrum. Must be larger than zero.
 Suggested value is the core-hole broadening.

\paragraph{Schema}
\begin{Verbatim}
|X| F
\end{Verbatim}

\subsection{ocean.cnbse.gmres.elist}
\label{kw:ocean-cnbse-gmres-elist}
\begin{longtable}{|p{1.5in}|p{4.5in}|}
\hline
Keyword & ocean.cnbse.gmres.elist \\
\hline
Kind & scalar \\
\hline
Type & Float \\
\hline
Default & Required / No default \\
\hline
Deprecated & No \\
\hline
\end{longtable}
 List of energies at which to run the GMRES algorithm in eV. Lowest
 unoccupied state is set at zero without the core-hole.

\paragraph{Schema}
\begin{Verbatim}
|X| F ...
\end{Verbatim}

\subsection{ocean.cnbse.gmres.erange}
\label{kw:ocean-cnbse-gmres-erange}
\begin{longtable}{|p{1.5in}|p{4.5in}|}
\hline
Keyword & ocean.cnbse.gmres.erange \\
\hline
Kind & scalar \\
\hline
Type & Float \\
\hline
Default & Required / No default \\
\hline
Deprecated & No \\
\hline
\end{longtable}
 Use energies at with even spacing specified by
 emin emax estep

\paragraph{Schema}
\begin{Verbatim}
|X| F
\end{Verbatim}

\subsection{ocean.cnbse.gmres.ffff}
\label{kw:ocean-cnbse-gmres-ffff}
\begin{longtable}{|p{1.5in}|p{4.5in}|}
\hline
Keyword & ocean.cnbse.gmres.ffff \\
\hline
Kind & scalar \\
\hline
Type & Float \\
\hline
Default & Required / No default \\
\hline
Deprecated & No \\
\hline
\end{longtable}
 Sets convergence criterion for GMRES.

\paragraph{Schema}
\begin{Verbatim}
|X| F
\end{Verbatim}

\subsection{ocean.cnbse.gmres.gprc}
\label{kw:ocean-cnbse-gmres-gprc}
\begin{longtable}{|p{1.5in}|p{4.5in}|}
\hline
Keyword & ocean.cnbse.gmres.gprc \\
\hline
Kind & scalar \\
\hline
Type & Float \\
\hline
Default & Required / No default \\
\hline
Deprecated & No \\
\hline
\end{longtable}
 Set Lorenzian broadening in Hartree in the GMRES algorithm when
 pre-conditioning. Default is 0.5, which should be reasonable for K edges.

\paragraph{Schema}
\begin{Verbatim}
|X| F
\end{Verbatim}

\subsection{ocean.cnbse.gmres.nloop}
\label{kw:ocean-cnbse-gmres-nloop}
\begin{longtable}{|p{1.5in}|p{4.5in}|}
\hline
Keyword & ocean.cnbse.gmres.nloop \\
\hline
Kind & scalar \\
\hline
Type & Integer \\
\hline
Default & Required / No default \\
\hline
Deprecated & No \\
\hline
\end{longtable}
 Restart GMRES algorithm will restart if the subspace grows to size nloop.

\paragraph{Schema}
\begin{Verbatim}
|X| I
\end{Verbatim}

\subsection{ocean.cnbse.mode}
\label{kw:ocean-cnbse-mode}
\begin{longtable}{|p{1.5in}|p{4.5in}|}
\hline
Keyword & ocean.cnbse.mode \\
\hline
Kind & scalar \\
\hline
Type & String \\
\hline
Default & Required / No default \\
\hline
Deprecated & No \\
\hline
\end{longtable}
 Which spectroscopy to calculate: "XAS", "XES"
 NRIXS and XRS are treated the same as XAS.

\paragraph{Schema}
\begin{Verbatim}
|X| S
\end{Verbatim}

\subsection{ocean.cnbse.niter}
\label{kw:ocean-cnbse-niter}
\begin{longtable}{|p{1.5in}|p{4.5in}|}
\hline
Keyword & ocean.cnbse.niter \\
\hline
Kind & scalar \\
\hline
Type & Integer \\
\hline
Default & Required / No default \\
\hline
Deprecated & No \\
\hline
\end{longtable}
 Number of Haydock iterations. Default is 100.

\paragraph{Schema}
\begin{Verbatim}
|X| I
\end{Verbatim}

\subsection{ocean.cnbse.rad}
\label{kw:ocean-cnbse-rad}
\begin{longtable}{|p{1.5in}|p{4.5in}|}
\hline
Keyword & ocean.cnbse.rad \\
\hline
Kind & scalar \\
\hline
Type & Float \\
\hline
Default & Required / No default \\
\hline
Deprecated & No \\
\hline
\end{longtable}
 One of the screening radius as defined by ocean.screen.shells that will be
 used in the BSE calculation.

\paragraph{Schema}
\begin{Verbatim}
|X| F
\end{Verbatim}

\subsection{ocean.cnbse.solver}
\label{kw:ocean-cnbse-solver}
\begin{longtable}{|p{1.5in}|p{4.5in}|}
\hline
Keyword & ocean.cnbse.solver \\
\hline
Kind & scalar \\
\hline
Type & String \\
\hline
Default & Required / No default \\
\hline
Deprecated & No \\
\hline
\end{longtable}
 Choose between Haydock algorithm (full spectrum, no excitons) or GMRES
 (single energy, exciton density).

\paragraph{Schema}
\begin{Verbatim}
|X| S
\end{Verbatim}

\subsection{ocean.cnbse.spect\_range}
\label{kw:ocean-cnbse-spect-range}
\begin{longtable}{|p{1.5in}|p{4.5in}|}
\hline
Keyword & ocean.cnbse.spect\_range \\
\hline
Kind & vector \\
\hline
Type & Vector: Integer, Float, Float \\
\hline
Default & Required / No default \\
\hline
Deprecated & No \\
\hline
\end{longtable}
 Sets energies over which the spectrum will be calculated.
 nsteps emin emax

\paragraph{Schema}
\begin{Verbatim}
|X| I F F
\end{Verbatim}

\subsection{ocean.cnbse.strength}
\label{kw:ocean-cnbse-strength}
\begin{longtable}{|p{1.5in}|p{4.5in}|}
\hline
Keyword & ocean.cnbse.strength \\
\hline
Kind & scalar \\
\hline
Type & Float \\
\hline
Default & Required / No default \\
\hline
Deprecated & No \\
\hline
\end{longtable}
 Sets the strength for the two interaction terms in the BSE Hamiltonian.
 Default is 1, and for emission 0. This can be changed for analysis.

\paragraph{Schema}
\begin{Verbatim}
|X| F
\end{Verbatim}

\subsection{ocean.cnbse.xmesh}
\label{kw:ocean-cnbse-xmesh}
\begin{longtable}{|p{1.5in}|p{4.5in}|}
\hline
Keyword & ocean.cnbse.xmesh \\
\hline
Kind & vector \\
\hline
Type & Vector: Integer, Integer, Integer \\
\hline
Default & Required / No default \\
\hline
Deprecated & No \\
\hline
\end{longtable}
 Wave-functions are converted into the NIST BSE format and condensed onto a
 grid of size ocean.cnbse.xmesh. Default is guessed in the code.

\paragraph{Schema}
\begin{Verbatim}
|X| I I I
\end{Verbatim}

\subsection{ocean.core\_offset}
\label{kw:ocean-core-offset}
\begin{longtable}{|p{1.5in}|p{4.5in}|}
\hline
Keyword & ocean.core\_offset \\
\hline
Kind & scalar \\
\hline
Type & String \\
\hline
Default & Required / No default \\
\hline
Deprecated & No \\
\hline
\end{longtable}
 If true, core-level shift will be calculated, false, no shift, or a number
 to specify the shift by hand.

\paragraph{Schema}
\begin{Verbatim}
|X| S
\end{Verbatim}

\subsection{ocean.degauss}
\label{kw:ocean-degauss}
\begin{longtable}{|p{1.5in}|p{4.5in}|}
\hline
Keyword & ocean.degauss \\
\hline
Kind & scalar \\
\hline
Type & Float \\
\hline
Default & Required / No default \\
\hline
Deprecated & No \\
\hline
\end{longtable}
 Broadening for the smearing function in Ryd. for metallic occupations in the
 DFT calculations.

\paragraph{Schema}
\begin{Verbatim}
|X| F
\end{Verbatim}

\subsection{ocean.dft}
\label{kw:ocean-dft}
\begin{longtable}{|p{1.5in}|p{4.5in}|}
\hline
Keyword & ocean.dft \\
\hline
Kind & scalar \\
\hline
Type & String \\
\hline
Default & abi \\
\hline
Deprecated & No \\
\hline
\end{longtable}
 Which DFT code to use: qe - QuantumEspresso, abi - ABINIT

\paragraph{Schema}
\begin{Verbatim}
abi |X| S
\end{Verbatim}

\subsection{ocean.dft\_energy\_range}
\label{kw:ocean-dft-energy-range}
\begin{longtable}{|p{1.5in}|p{4.5in}|}
\hline
Keyword & ocean.dft\_energy\_range \\
\hline
Kind & scalar \\
\hline
Type & Float \\
\hline
Default & Required / No default \\
\hline
Deprecated & No \\
\hline
\end{longtable}
 Instead of setting the number of bands (ocean.nbands) the user may request
 an energy range in eV for the final state wave-functions. This is an estimate
 using the volume of the unit cell and may be unreliable. Default is 25 eV.

\paragraph{Schema}
\begin{Verbatim}
|X| F
\end{Verbatim}

\subsection{ocean.diemac}
\label{kw:ocean-diemac}
\begin{longtable}{|p{1.5in}|p{4.5in}|}
\hline
Keyword & ocean.diemac \\
\hline
Kind & scalar \\
\hline
Type & Float \\
\hline
Default & Required / No default \\
\hline
Deprecated & No \\
\hline
\end{longtable}
 Macroscopic dielectric matrix.

\paragraph{Schema}
\begin{Verbatim}
|X| F
\end{Verbatim}

\subsection{ocean.ecut}
\label{kw:ocean-ecut}
\begin{longtable}{|p{1.5in}|p{4.5in}|}
\hline
Keyword & ocean.ecut \\
\hline
Kind & scalar \\
\hline
Type & Float \\
\hline
Default & Required / No default \\
\hline
Deprecated & No \\
\hline
\end{longtable}
 Plane wave basis truncation energy in Rydberg

\paragraph{Schema}
\begin{Verbatim}
|X| F
\end{Verbatim}

\subsection{ocean.edges}
\label{kw:ocean-edges}
\begin{longtable}{|p{1.5in}|p{4.5in}|}
\hline
Keyword & ocean.edges \\
\hline
Kind & repeating-record \\
\hline
Type & Repeating record with columns: Integer, Integer, Integer \\
\hline
Default & Required / No default \\
\hline
Deprecated & No \\
\hline
\end{longtable}
 Specify the atoms and edges to calculate. Each edge entry consists of 3
 integers. The first, if greater than zero, is the index of the atom for
 which the edge should be calculated. The second and third numbers are the
 principle quantum number and angular momentum.
 '1 2 1' denotes the L edges of the first atom
 If the first number is negative, then it is interpreted as -Z, where
 Z is the atomic number, and in that case edges of all atoms with that
 atomic number will be calculated.
 '-22 2 1' will run the L edges of every titanium atom in the cell.

\paragraph{Schema}
\begin{Verbatim}
|X| I I I
|X| ...
\end{Verbatim}

\subsection{ocean.fband}
\label{kw:ocean-fband}
\begin{longtable}{|p{1.5in}|p{4.5in}|}
\hline
Keyword & ocean.fband \\
\hline
Kind & scalar \\
\hline
Type & Float \\
\hline
Default & Required / No default \\
\hline
Deprecated & No \\
\hline
\end{longtable}
 Determines the number of occupied bands to be included in the SCF calculation
 of the density. For insulators the default (fband = 0.125) is fine, but
 for metals the highest band the in the density calculation should have no
 occupation weight. Using fband, the number of bands is determined by
 n = natom*fband.

\paragraph{Schema}
\begin{Verbatim}
|X| F
\end{Verbatim}

\subsection{ocean.k0}
\label{kw:ocean-k0}
\begin{longtable}{|p{1.5in}|p{4.5in}|}
\hline
Keyword & ocean.k0 \\
\hline
Kind & vector \\
\hline
Type & Vector: Float, Float, Float \\
\hline
Default & Required / No default \\
\hline
Deprecated & No \\
\hline
\end{longtable}
 DFT states are calculated using a shifted k-point grid. The first k-point is
 given by:
   $1/(N_{k}(1)*k0(1))$, $1/(N_{k}(2)*k0(2)$, $1/(N_{k}(3)*k0(3)$
 where $N_k$ is given by ocean.nkpt. The default shift is $k0 = 1/8$, $2/8$, $3/8$.

\paragraph{Schema}
\begin{Verbatim}
|X| F F F
\end{Verbatim}

\subsection{ocean.ldau}
\label{kw:ocean-ldau}
\begin{longtable}{|p{1.5in}|p{4.5in}|}
\hline
Keyword & ocean.ldau \\
\hline
Kind & repeating-record \\
\hline
Type & Repeating record with columns: Logical \\
\hline
Default & Required / No default \\
\hline
Deprecated & No \\
\hline
\end{longtable}
 Controls the U parameters for the DFT run. Only works with QuantumEspresso
 currently. Not sure what the parameters are, but I believe they are
 logical flag to run/not run LDA+U, U for each atom that has spin
 associated with smag.

\paragraph{Schema}
\begin{Verbatim}
|X| L
|X| F
|X| ...
\end{Verbatim}

\subsection{ocean.metal}
\label{kw:ocean-metal}
\begin{longtable}{|p{1.5in}|p{4.5in}|}
\hline
Keyword & ocean.metal \\
\hline
Kind & scalar \\
\hline
Type & Logical \\
\hline
Default & Required / No default \\
\hline
Deprecated & No \\
\hline
\end{longtable}
 If True, the code will complain unless occopt is 3. If False, the codes
 expects occopt = 1.

\paragraph{Schema}
\begin{Verbatim}
|X| L
\end{Verbatim}

\subsection{ocean.natom}
\label{kw:ocean-natom}
\begin{longtable}{|p{1.5in}|p{4.5in}|}
\hline
Keyword & ocean.natom \\
\hline
Kind & scalar \\
\hline
Type & Integer \\
\hline
Default & Required / No default \\
\hline
Deprecated & No \\
\hline
\end{longtable}
 Number of atoms in the unit cell

\paragraph{Schema}
\begin{Verbatim}
|X| I
\end{Verbatim}

\subsection{ocean.nbands}
\label{kw:ocean-nbands}
\begin{longtable}{|p{1.5in}|p{4.5in}|}
\hline
Keyword & ocean.nbands \\
\hline
Kind & scalar \\
\hline
Type & Integer \\
\hline
Default & Required / No default \\
\hline
Deprecated & No \\
\hline
\end{longtable}
 Number of bands for the final state calculations. This includes valence and conduction bands.

\paragraph{Schema}
\begin{Verbatim}
|X| I
\end{Verbatim}

\subsection{ocean.ngkpt}
\label{kw:ocean-ngkpt}
\begin{longtable}{|p{1.5in}|p{4.5in}|}
\hline
Keyword & ocean.ngkpt \\
\hline
Kind & vector \\
\hline
Type & Vector: Integer, Integer, Integer \\
\hline
Default & Required / No default \\
\hline
Deprecated & No \\
\hline
\end{longtable}
 Number of k-points in each direction used to sample the cell for
 the ground-state calculation of the density.

\paragraph{Schema}
\begin{Verbatim}
|X| I I I
\end{Verbatim}

\subsection{ocean.nkpt}
\label{kw:ocean-nkpt}
\begin{longtable}{|p{1.5in}|p{4.5in}|}
\hline
Keyword & ocean.nkpt \\
\hline
Kind & vector \\
\hline
Type & Vector: Integer, Integer, Integer \\
\hline
Default & Required / No default \\
\hline
Deprecated & No \\
\hline
\end{longtable}
 Number of k-points in each direction used to sample the
 cell for the calculation of the
 final states.

\paragraph{Schema}
\begin{Verbatim}
|X| I I I
\end{Verbatim}

\subsection{ocean.nspin}
\label{kw:ocean-nspin}
\begin{longtable}{|p{1.5in}|p{4.5in}|}
\hline
Keyword & ocean.nspin \\
\hline
Kind & scalar \\
\hline
Type & Integer \\
\hline
Default & Required / No default \\
\hline
Deprecated & No \\
\hline
\end{longtable}
 Choose paramagnetic (nspin = 1), or spin dependent (nspin = 2).

\paragraph{Schema}
\begin{Verbatim}
|X| I
\end{Verbatim}

\subsection{ocean.nstep}
\label{kw:ocean-nstep}
\begin{longtable}{|p{1.5in}|p{4.5in}|}
\hline
Keyword & ocean.nstep \\
\hline
Kind & scalar \\
\hline
Type & Integer \\
\hline
Default & Required / No default \\
\hline
Deprecated & No \\
\hline
\end{longtable}
 Maximum number of iterations for the DFT SCF calculations. Default is 50

\paragraph{Schema}
\begin{Verbatim}
|X| I
\end{Verbatim}

\subsection{ocean.ntypat}
\label{kw:ocean-ntypat}
\begin{longtable}{|p{1.5in}|p{4.5in}|}
\hline
Keyword & ocean.ntypat \\
\hline
Kind & scalar \\
\hline
Type & Integer \\
\hline
Default & Required / No default \\
\hline
Deprecated & No \\
\hline
\end{longtable}
 Number of different types (usually species) of atom in the unit cell.

\paragraph{Schema}
\begin{Verbatim}
|X| I
\end{Verbatim}

\subsection{ocean.occopt}
\label{kw:ocean-occopt}
\begin{longtable}{|p{1.5in}|p{4.5in}|}
\hline
Keyword & ocean.occopt \\
\hline
Kind & scalar \\
\hline
Type & Integer \\
\hline
Default & Required / No default \\
\hline
Deprecated & No \\
\hline
\end{longtable}
 Controls determination of occupation. Most important
 values are:
 1 - States are all doubly degenerate and either occupied or empty depending on band.
 3 - States are all doubly degenerate but can have fractional occupations
     depending on Fermi function. Suitable for metals.

\paragraph{Schema}
\begin{Verbatim}
|X| I
\end{Verbatim}

\subsection{ocean.opf.fill}
\label{kw:ocean-opf-fill}
\begin{longtable}{|p{1.5in}|p{4.5in}|}
\hline
Keyword & ocean.opf.fill \\
\hline
Kind & vector \\
\hline
Type & Vector: Integer, String \\
\hline
Default & Required / No default \\
\hline
Deprecated & No \\
\hline
\end{longtable}
 atomic number followed by the name of the .fill file to use for this calculation.

\paragraph{Schema}
\begin{Verbatim}
|X| I S
\end{Verbatim}

\subsection{ocean.opf.hfkgrid}
\label{kw:ocean-opf-hfkgrid}
\begin{longtable}{|p{1.5in}|p{4.5in}|}
\hline
Keyword & ocean.opf.hfkgrid \\
\hline
Kind & vector \\
\hline
Type & Vector: Integer, Integer \\
\hline
Default & Required / No default \\
\hline
Deprecated & No \\
\hline
\end{longtable}
 The first parameter sets the grid for hfk.x, and should be left at 2000.
 The second determines the maximum number of proejectors per angular momentum
 channel. The default of 20 will work for many calculations.

\paragraph{Schema}
\begin{Verbatim}
|X| I I
\end{Verbatim}

\subsection{ocean.opf.opts}
\label{kw:ocean-opf-opts}
\begin{longtable}{|p{1.5in}|p{4.5in}|}
\hline
Keyword & ocean.opf.opts \\
\hline
Kind & vector \\
\hline
Type & Vector: Integer, String \\
\hline
Default & Required / No default \\
\hline
Deprecated & No \\
\hline
\end{longtable}
 atomic number followed by the name of the .opts file to use for this calculation.

\paragraph{Schema}
\begin{Verbatim}
|X| I S
\end{Verbatim}

\subsection{ocean.opf.program}
\label{kw:ocean-opf-program}
\begin{longtable}{|p{1.5in}|p{4.5in}|}
\hline
Keyword & ocean.opf.program \\
\hline
Kind & scalar \\
\hline
Type & String \\
\hline
Default & Required / No default \\
\hline
Deprecated & No \\
\hline
\end{longtable}
 Use oncvpsp UPS pseudopotentials with qe

\paragraph{Schema}
\begin{Verbatim}
|X| S
\end{Verbatim}

\subsection{ocean.para\_prefix}
\label{kw:ocean-para-prefix}
\begin{longtable}{|p{1.5in}|p{4.5in}|}
\hline
Keyword & ocean.para\_prefix \\
\hline
Kind & scalar \\
\hline
Type & String \\
\hline
Default & Required / No default \\
\hline
Deprecated & No \\
\hline
\end{longtable}
 Set the command to use for parallel runs, i.e., mpirun, mpiexec etc,
 along with any flags to use.

\paragraph{Schema}
\begin{Verbatim}
|X| S ...
\end{Verbatim}

\subsection{ocean.photon.energy}
\label{kw:ocean-photon-energy}
\begin{longtable}{|p{1.5in}|p{4.5in}|}
\hline
Keyword & ocean.photon.energy \\
\hline
Kind & scalar \\
\hline
Type & Float \\
\hline
Default & Required / No default \\
\hline
Deprecated & No \\
\hline
\end{longtable}
 Set energy of x-rays at edge

\paragraph{Schema}
\begin{Verbatim}
|X| F
\end{Verbatim}

\subsection{ocean.photon.operator}
\label{kw:ocean-photon-operator}
\begin{longtable}{|p{1.5in}|p{4.5in}|}
\hline
Keyword & ocean.photon.operator \\
\hline
Kind & scalar \\
\hline
Type & String \\
\hline
Default & dipole \\
\hline
Deprecated & No \\
\hline
\end{longtable}
 Select the operator to be used to calculate the (diple or quad)

\paragraph{Schema}
\begin{Verbatim}
dipole |X| S
\end{Verbatim}

\subsection{ocean.photon.polarization}
\label{kw:ocean-photon-polarization}
\begin{longtable}{|p{1.5in}|p{4.5in}|}
\hline
Keyword & ocean.photon.polarization \\
\hline
Kind & vector \\
\hline
Type & Vector: Float, Float, Float \\
\hline
Default & Required / No default \\
\hline
Deprecated & No \\
\hline
\end{longtable}
 Direction of polarization for XANES or XES calculations

\paragraph{Schema}
\begin{Verbatim}
|X| F F F
\end{Verbatim}

\subsection{ocean.photon.qhat}
\label{kw:ocean-photon-qhat}
\begin{longtable}{|p{1.5in}|p{4.5in}|}
\hline
Keyword & ocean.photon.qhat \\
\hline
Kind & vector \\
\hline
Type & Vector: Float, Float, Float \\
\hline
Default & Required / No default \\
\hline
Deprecated & No \\
\hline
\end{longtable}
 Direction of q-vector for quadrupole or NRIXS calculations

\paragraph{Schema}
\begin{Verbatim}
|X| F F F
\end{Verbatim}

\subsection{ocean.photon\_q}
\label{kw:ocean-photon-q}
\begin{longtable}{|p{1.5in}|p{4.5in}|}
\hline
Keyword & ocean.photon\_q \\
\hline
Kind & vector \\
\hline
Type & Vector: Float, Float, Float \\
\hline
Default & Required / No default \\
\hline
Deprecated & No \\
\hline
\end{longtable}
 For UV/VIS and RIXS calculations, the valence and conduction band states
 must be offset from each other (by the momentum that is absorbed or
 transferred). This sets that momentum offset in units of the reciprocal
 lattice vectors of the system. Default is 0 0 0.

\paragraph{Schema}
\begin{Verbatim}
|X| F F F
\end{Verbatim}

\subsection{ocean.pp\_list}
\label{kw:ocean-pp-list}
\begin{longtable}{|p{1.5in}|p{4.5in}|}
\hline
Keyword & ocean.pp\_list \\
\hline
Kind & repeating-record \\
\hline
Type & Repeating record with columns: String \\
\hline
Default & Required / No default \\
\hline
Deprecated & No \\
\hline
\end{longtable}
 Names of the pseudopotential files to be used listed in the same order as
 ocean.znucl

\paragraph{Schema}
\begin{Verbatim}
|X| S
|X| ...
\end{Verbatim}

\subsection{ocean.rprim}
\label{kw:ocean-rprim}
\begin{longtable}{|p{1.5in}|p{4.5in}|}
\hline
Keyword & ocean.rprim \\
\hline
Kind & matrix \\
\hline
Type & 3 rows of Float, Float, Float \\
\hline
Default & Required / No default \\
\hline
Deprecated & No \\
\hline
\end{longtable}
 Primitive vectors of unit cell

\paragraph{Schema}
\begin{Verbatim}
|X| F F F
|X| F F F
|X| F F F
\end{Verbatim}

\subsection{ocean.scfac}
\label{kw:ocean-scfac}
\begin{longtable}{|p{1.5in}|p{4.5in}|}
\hline
Keyword & ocean.scfac \\
\hline
Kind & scalar \\
\hline
Type & Float \\
\hline
Default & Required / No default \\
\hline
Deprecated & No \\
\hline
\end{longtable}
 The is the Slater integral factor to scale the slater integrals. For 3d
 transition metals 0.8 is a reasonable value, while 0.6 is better for
 f-electron atoms.

\paragraph{Schema}
\begin{Verbatim}
|X| F
\end{Verbatim}

\subsection{ocean.scratch}
\label{kw:ocean-scratch}
\begin{longtable}{|p{1.5in}|p{4.5in}|}
\hline
Keyword & ocean.scratch \\
\hline
Kind & scalar \\
\hline
Type & String \\
\hline
Default & Required / No default \\
\hline
Deprecated & No \\
\hline
\end{longtable}
 Where to write scratch files

\paragraph{Schema}
\begin{Verbatim}
|X| S
\end{Verbatim}

\subsection{ocean.screen.nbands}
\label{kw:ocean-screen-nbands}
\begin{longtable}{|p{1.5in}|p{4.5in}|}
\hline
Keyword & ocean.screen.nbands \\
\hline
Kind & scalar \\
\hline
Type & Integer \\
\hline
Default & Required / No default \\
\hline
Deprecated & No \\
\hline
\end{longtable}
 Number of bands to use for the screening calculations. This requires a large
 number of bands. Should be about 100eV above the Fermi level, but convergence
 should be checked.

\paragraph{Schema}
\begin{Verbatim}
|X| I
\end{Verbatim}

\subsection{ocean.screen.nkpt}
\label{kw:ocean-screen-nkpt}
\begin{longtable}{|p{1.5in}|p{4.5in}|}
\hline
Keyword & ocean.screen.nkpt \\
\hline
Kind & vector \\
\hline
Type & Vector: Integer, Integer, Integer \\
\hline
Default & Required / No default \\
\hline
Deprecated & No \\
\hline
\end{longtable}
 Number of k-point to be used for the screening calculations. Default is 2 2 2,
 and is sufficient for a wide variety of systems, although very small unit
 cells might require more, and large unit cells may only require gamma point.

\paragraph{Schema}
\begin{Verbatim}
|X| I I I
\end{Verbatim}

\subsection{ocean.screen.shells}
\label{kw:ocean-screen-shells}
\begin{longtable}{|p{1.5in}|p{4.5in}|}
\hline
Keyword & ocean.screen.shells \\
\hline
Kind & scalar \\
\hline
Type & Float \\
\hline
Default & Required / No default \\
\hline
Deprecated & No \\
\hline
\end{longtable}
 The screening calculation is RPA at small radius and a model at large radius.
 Teh cross-over between the RPA and model is set by shells in Bohr.
 Several different radii can be chosen to look at the convergence.
 Convergence is usually reached at arount 3-4 Bohr. The supercell defined by
 ocean.paw.nkpt must have dimensions larger than the radius specified by
 ocean.screen.shells or results may be unreliable. The default of 3.5 Bohr is
 usually OK. Values > 6 Bohr should not be used.

\paragraph{Schema}
\begin{Verbatim}
|X| F ...
\end{Verbatim}

\subsection{ocean.screen\_energy\_range}
\label{kw:ocean-screen-energy-range}
\begin{longtable}{|p{1.5in}|p{4.5in}|}
\hline
Keyword & ocean.screen\_energy\_range \\
\hline
Kind & scalar \\
\hline
Type & Float \\
\hline
Default & Required / No default \\
\hline
Deprecated & No \\
\hline
\end{longtable}
 Instead of setting number of bands for the screening, the user may request
 an energy range in eV. This is only an estimate, and may need to be converged.

\paragraph{Schema}
\begin{Verbatim}
|X| F
\end{Verbatim}

\subsection{ocean.ser\_prefix}
\label{kw:ocean-ser-prefix}
\begin{longtable}{|p{1.5in}|p{4.5in}|}
\hline
Keyword & ocean.ser\_prefix \\
\hline
Kind & scalar \\
\hline
Type & String \\
\hline
Default & Required / No default \\
\hline
Deprecated & No \\
\hline
\end{longtable}
 prefix for serial parts of the code, for example, cut3d should be run
 with only 1 processor, i.e., mpirun -n 1

\paragraph{Schema}
\begin{Verbatim}
|X| S ...
\end{Verbatim}

\subsection{ocean.smag}
\label{kw:ocean-smag}
\begin{longtable}{|p{1.5in}|p{4.5in}|}
\hline
Keyword & ocean.smag \\
\hline
Kind & repeating-record \\
\hline
Type & Repeating record with columns: Integer, Integer, Integer \\
\hline
Default & Required / No default \\
\hline
Deprecated & No \\
\hline
\end{longtable}
 Controls initial magnetism for the DFT calculations.

\paragraph{Schema}
\begin{Verbatim}
|X| I I I
|X| ...
\end{Verbatim}

\subsection{ocean.smag.ixc}
\label{kw:ocean-smag-ixc}
\begin{longtable}{|p{1.5in}|p{4.5in}|}
\hline
Keyword & ocean.smag.ixc \\
\hline
Kind & scalar \\
\hline
Type & Integer \\
\hline
Default & Required / No default \\
\hline
Deprecated & No \\
\hline
\end{longtable}
 Used for abinit runs in smag card. Don't know the meaning at the moment.
 This should be combiened with ocean.smag to write the card to the ocean
 input.

\paragraph{Schema}
\begin{Verbatim}
|X| I
\end{Verbatim}

\subsection{ocean.spin-orbit}
\label{kw:ocean-spin-orbit}
\begin{longtable}{|p{1.5in}|p{4.5in}|}
\hline
Keyword & ocean.spin-orbit \\
\hline
Kind & scalar \\
\hline
Type & Float \\
\hline
Default & Required / No default \\
\hline
Deprecated & No \\
\hline
\end{longtable}
 If ocean.spin-orbit >= 0, the code will not automatically calculate the
 spin-orbit splitting.

\paragraph{Schema}
\begin{Verbatim}
|X| F
\end{Verbatim}

\subsection{ocean.toldfe}
\label{kw:ocean-toldfe}
\begin{longtable}{|p{1.5in}|p{4.5in}|}
\hline
Keyword & ocean.toldfe \\
\hline
Kind & scalar \\
\hline
Type & Float \\
\hline
Default & Required / No default \\
\hline
Deprecated & No \\
\hline
\end{longtable}
 Total energy convergence parameter for the density run. Default is $10^{-6}$

\paragraph{Schema}
\begin{Verbatim}
|X| F
\end{Verbatim}

\subsection{ocean.tolwfr}
\label{kw:ocean-tolwfr}
\begin{longtable}{|p{1.5in}|p{4.5in}|}
\hline
Keyword & ocean.tolwfr \\
\hline
Kind & scalar \\
\hline
Type & Float \\
\hline
Default & Required / No default \\
\hline
Deprecated & No \\
\hline
\end{longtable}
 Convergence criterium for the non-scf wave-function calculations.
 Default is $10^{-16}$

\paragraph{Schema}
\begin{Verbatim}
|X| F
\end{Verbatim}

\subsection{ocean.typat}
\label{kw:ocean-typat}
\begin{longtable}{|p{1.5in}|p{4.5in}|}
\hline
Keyword & ocean.typat \\
\hline
Kind & scalar \\
\hline
Type & Integer \\
\hline
Default & Required / No default \\
\hline
Deprecated & No \\
\hline
\end{longtable}
 Type of each atom in the unit cell denoted by the integer index corresponding
 to the ocean.znucl list of atomic numbers. This list must be in the same
 as the coordinates listed in ocean.xred

\paragraph{Schema}
\begin{Verbatim}
|X| I ...
\end{Verbatim}

\subsection{ocean.xred}
\label{kw:ocean-xred}
\begin{longtable}{|p{1.5in}|p{4.5in}|}
\hline
Keyword & ocean.xred \\
\hline
Kind & repeating-record \\
\hline
Type & Repeating record with columns: Float, Float, Float \\
\hline
Default & Required / No default \\
\hline
Deprecated & No \\
\hline
\end{longtable}
 Reduced coordinats of atoms in the unit cell.

\paragraph{Schema}
\begin{Verbatim}
|X| F F F
|X| ...
\end{Verbatim}

\subsection{ocean.znucl}
\label{kw:ocean-znucl}
\begin{longtable}{|p{1.5in}|p{4.5in}|}
\hline
Keyword & ocean.znucl \\
\hline
Kind & scalar \\
\hline
Type & Integer \\
\hline
Default & Required / No default \\
\hline
Deprecated & No \\
\hline
\end{longtable}
 List of length ntypat with atomic numbers of each type of atom in the unit cell

\paragraph{Schema}
\begin{Verbatim}
|X| I ...
\end{Verbatim}

\subsection{ocean.zymb}
\label{kw:ocean-zymb}
\begin{longtable}{|p{1.5in}|p{4.5in}|}
\hline
Keyword & ocean.zymb \\
\hline
Kind & scalar \\
\hline
Type & String \\
\hline
Default & Required / No default \\
\hline
Deprecated & No \\
\hline
\end{longtable}
 For use with QE. Each element listed in ocean.znucle can have a symbol
 associated with it.

\paragraph{Schema}
\begin{Verbatim}
|X| S ...
\end{Verbatim}

\section{Miscellaneous}

\subsection{ene\_int}
\label{kw:ene-int}
\begin{longtable}{|p{1.5in}|p{4.5in}|}
\hline
Keyword & ene\_int \\
\hline
Kind & scalar \\
\hline
Type & Float \\
\hline
Default & Required / No default \\
\hline
Deprecated & No \\
\hline
\end{longtable}
 Internal (electronic) energy of the system, in au.

\paragraph{Schema}
\begin{Verbatim}
|X| F
\end{Verbatim}

\subsection{mac\_diel\_const}
\label{kw:mac-diel-const}
\begin{longtable}{|p{1.5in}|p{4.5in}|}
\hline
Keyword & mac\_diel\_const \\
\hline
Kind & scalar \\
\hline
Type & Float \\
\hline
Default & Required / No default \\
\hline
Deprecated & No \\
\hline
\end{longtable}
 Approximate value of the macroscopic dielectric constant of the system used in
 some methods to accelerate convergence of the SCF cycle.

\paragraph{Schema}
\begin{Verbatim}
|X| F
\end{Verbatim}

\subsection{magmom\_by\_label}
\label{kw:magmom-by-label}
\begin{longtable}{|p{1.5in}|p{4.5in}|}
\hline
Keyword & magmom\_by\_label \\
\hline
Kind & repeating-record \\
\hline
Type & Repeating record with columns: String, Float \\
\hline
Default & Required / No default \\
\hline
Deprecated & No \\
\hline
\end{longtable}
 set spin for elements of a cif file by the label

\paragraph{Schema}
\begin{Verbatim}
|X| S F
|X| ...
\end{Verbatim}

\subsection{mp\_apikey}
\label{kw:mp-apikey}
\begin{longtable}{|p{1.5in}|p{4.5in}|}
\hline
Keyword & mp\_apikey \\
\hline
Kind & scalar \\
\hline
Type & String \\
\hline
Default & Required / No default \\
\hline
Deprecated & No \\
\hline
\end{longtable}
 Specify your mp API key. This should probably be done with an environment variable instead for security.

\paragraph{Schema}
\begin{Verbatim}
|X| S
\end{Verbatim}

\subsection{mp\_id}
\label{kw:mp-id}
\begin{longtable}{|p{1.5in}|p{4.5in}|}
\hline
Keyword & mp\_id \\
\hline
Kind & scalar \\
\hline
Type & String \\
\hline
Default & Required / No default \\
\hline
Deprecated & No \\
\hline
\end{longtable}
 Single material project id. Can be mp-NNNN or just the number.

\paragraph{Schema}
\begin{Verbatim}
|X| S
\end{Verbatim}

\subsection{spectralFunction\_file}
\label{kw:spectralFunction-file}
\begin{longtable}{|p{1.5in}|p{4.5in}|}
\hline
Keyword & spectralFunction\_file \\
\hline
Kind & scalar \\
\hline
Type & String \\
\hline
Default & Required / No default \\
\hline
Deprecated & No \\
\hline
\end{longtable}
 file to read spectral function from - 2 column file

\paragraph{Schema}
\begin{Verbatim}
|X| S
\end{Verbatim}

\subsection{xanes\_file}
\label{kw:xanes-file}
\begin{longtable}{|p{1.5in}|p{4.5in}|}
\hline
Keyword & xanes\_file \\
\hline
Kind & scalar \\
\hline
Type & String \\
\hline
Default & Required / No default \\
\hline
Deprecated & No \\
\hline
\end{longtable}
 file to read xanes from - 2 column file

\paragraph{Schema}
\begin{Verbatim}
|X| S
\end{Verbatim}

\end{document}
